% prompt-injection-chapter.tex -- backing appendix chapter for the claim
% that prompt injection cannot (yet) be reliably prevented by filters,
% detectors, prompt techniques or trained models, and that the design must
% limit the blast radius.
%
% Author's note (verification and reproduction; NOT for the compiled text):
%   scholar session  : prompt-injection (11,360 records, searched 2026-11-27)
%   queries          : litteratursokning/prompt-injection-queries.sh (--print)
%   bib + provenance : litteratursokning/prompt-injection.bib
%                      (check_provenance.py --quotes: OK)
%   report, protocol : litteratursokning/prompt-injection-report.md
%   progress log     : litteratursokning/prompt-injection-progress.md
%   tier-1 sheets    : litteratursokning/prompt-injection-screen/
%   hit-list export  : litteratursokning/prompt-injection-full.tex (+ .csv)
%   Screened by an AI agent; no decision has been reviewed by the author.
%   Attributions OWASP2025LLM, OWASP2026LLM, Willison2025Trifecta are in
%   known-items.bib (verified by the known-items agent).

\chapter{Can prompt injection against LLM-integrated applications and
  agents be reliably prevented?}
\label{prompt-injection:chapter}

\emph{Status: screened by an AI agent; awaiting human review.}

The lecture says three things about prompt injection. First, a language
model receives its instructions and the data it works on in one stream of
text, so an application must assume that anything the model reads can
steer it. Second, published defences lower the success of attacks but
have been bypassed by attackers who adapt to them. Third, an application
should therefore be designed so that a successful injection can do
little harm: least privilege for its tools, authorisation enforced in
ordinary code, and human approval for actions that cannot be undone.
The OWASP Top 10 for LLM Applications states that \enquote{it is unclear
if there are fool-proof methods of prevention for prompt
injection}\autocite{OWASP2025LLM} and, in its 2026 edition, that
\enquote{no reliable prevention mechanism exists
today}\autocite{OWASP2026LLM}; Simon Willison writes that \enquote{we
still don't know how to 100\% reliably prevent this from
happening}\autocite{Willison2025Trifecta}. This chapter asks whether the
research literature agrees.

The research question has four parts:
(a)~what indirect prompt injection is, and whether it works against real
LLM-integrated applications and agents;
(b)~how defences --- input and output filtering, detection classifiers,
prompt hardening such as delimiters or spotlighting, and fine-tuned or
instruction-hierarchy models --- fare against adaptive attacks;
(c)~whether any approach gives guarantees, in particular system-level
designs that keep untrusted data out of the control flow, and at what
cost in usefulness; and
(d)~how easily system prompts and other hidden context are extracted.
The strength under test is \enquote{cannot be reliably prevented by any
known defence that leaves the agent unconstrained}, not \enquote{cannot
be prevented at all}; the counter-searches therefore look explicitly for
defences reported as robust, provable, solved or secure by design.

\section{Method}
\label{prompt-injection:method}

\paragraph{Databases and their health.}
All searches ran on 27~November~2026 against Web of Science (WoS),
IEEE Xplore (IEEE), OpenAlex (OA; not open access), Scopus and Semantic
Scholar (S2), each with the same Boolean query, years 2022--2026,
restricted to title and abstract where the database supports it (WoS,
OA, Scopus, S2), and without an upper limit on the number of results.
Semantic Scholar answered through its bulk search, which matches title
and abstract. DBLP answered every query with a page that stops automated
clients instead of with data, so DBLP could not be searched. IEEE Xplore
had spent its daily request quota before the four repair queries and the
two prompt-security queries (R1--R3, CR1--CR3) were run; those six cells
are failed, not empty. Scopus served full records (with abstracts) for
all queries; for the repair queries this needed a second run once the
university network route was available again.

\paragraph{Concepts and wordings.}
Each concept of the claim was searched in several wordings, fixed before
the queries were run (\cref{prompt-injection:concepts}). The queries
combine these blocks (\cref{prompt-injection:queries}): S1--S9 search for
the claim, C1--C7 against it, R1--R3 repair gaps that the check by
meaning revealed (below), and CR1--CR3 are their counter-queries. Every
counter-query reuses all concept blocks of its support query and adds a
block of words that a contradicting or qualifying paper would use
(overstated, ineffective, rarely, robust, near-zero, provable, solved,
by design, without loss of utility, and so on).

\begin{table}[tbp]
  \begin{sidecaption}{Concept blocks of the queries and their wordings.
    Words in quotation marks are phrases; a block matches if any of its
    wordings does.}[prompt-injection:concepts]
    \centering\footnotesize
    \begin{tabular}{@{}lp{0.72\linewidth}@{}}
      \toprule
      Block & Wordings \\
      \midrule
      Attack & \enquote{prompt injection(s)}, \enquote{indirect injection},
        \enquote{injected instruction(s)}, \enquote{injected prompt(s)},
        \enquote{instruction injection}, \enquote{goal hijacking},
        \enquote{prompt hijacking}, \enquote{task hijacking} \\
      System & LLM(s), \enquote{language model(s)}, agent(s), agentic,
        chatbot(s), \enquote{retrieval-augmented}, RAG \\
      Application & indirect, \enquote{LLM-integrated}, integrated,
        application(s), \enquote{real-world}, deployed, deployment,
        plugin(s), web, email, retrieval, \enquote{third-party}, tool(s) \\
      Defence & defense(s), defence(s), defend(ing), mitigation, mitigate,
        detection, detect(or), filter(ing), guardrail(s), delimiter(s),
        spotlighting, \enquote{instruction hierarchy}, \enquote{fine-tuning},
        \enquote{fine-tuned}, prevent(ion), protect(ion) \\
      Adaptive & adaptive, bypass(ed), evade, evasion, circumvent, break(s),
        optimi[sz]ation, optimized, \enquote{red teaming} \\
      Guarantee & provable, provably, guarantee(d), \enquote{by design},
        certified, formal, \enquote{information flow}, \enquote{control
        flow}, \enquote{data flow}, isolation, \enquote{privilege
        separation}, \enquote{least privilege}, capability,
        \enquote{policy enforcement}, sandbox(ing), \enquote{design
        pattern(s)}, \enquote{dual LLM}, deterministic, architecture \\
      Utility & utility, \enquote{task performance}, overhead, cost,
        \enquote{trade-off}, tradeoff, usability, benign \\
      Benchmark & benchmark(s), evaluation, evaluate, measure,
        \enquote{attack success}, \enquote{success rate} \\
      Extraction & \enquote{system prompt(s)}, \enquote{prompt extraction},
        \enquote{prompt leaking/leakage/leak}, \enquote{prompt stealing},
        \enquote{prompt theft}, \enquote{hidden prompt(s)}; with extract,
        leak, steal, reveal, recover, exfiltrate and their forms \\
      Separation & \enquote{instructions and data}, \enquote{instruction
        and data}, \enquote{data and instructions},
        \enquote{instruction-data}, \enquote{separate/separating
        instructions}, \enquote{separation of instructions},
        \enquote{structured query/queries}, \enquote{single channel},
        \enquote{same channel} \\
      Untrusted (R1) & \enquote{untrusted data/input(s)/content/sources/tool} \\
      Adversarial (R2) & \enquote{adversarial prompt(s)}, \enquote{malicious
        instruction(s)}, \enquote{malicious prompt(s)}, \enquote{hidden
        instruction(s)}; with exfiltration, \enquote{tool misuse}, hijack,
        manipulate \\
      Prompt security (R3) & \enquote{prompt attack(s)}, \enquote{prompt
        security}, \enquote{prompt-based attack(s)}, \enquote{adversarial
        prompting} \\
      \bottomrule
    \end{tabular}
  \end{sidecaption}
\end{table}

\begin{table}[tbp]
  \begin{sidecaption}{Queries and hits per database, searched 27 November
    2026. S = supporting query, C = counter-query, R = repair query added
    after the check by meaning, CR = its counter-query; the blocks are
    defined in \cref{prompt-injection:concepts}. Databases: OpenAlex (OA;
    not open access), IEEE Xplore (IEEE), Web of Science (WoS), Scopus,
    Semantic Scholar (S2, bulk search). Years 2022--2026; title and
    abstract where the database supports it; no upper limit, so the
    numbers are all hits. \enquote{--} means that the search failed: IEEE
    Xplore had spent its daily quota. $^{a}$ the database reported a few
    more hits than it delivered (at most 16 fewer). DBLP failed for every
    query and is not listed.}[prompt-injection:queries]
    \centering\footnotesize
    \begin{tabular}{@{}lp{0.42\linewidth}rrrrr@{}}
      \toprule
      & Keywords & OA & IEEE & WoS & Scopus & S2 \\
      \midrule
      S1 & Attack AND System & 3{,}703 & 476 & 449 & 878 & 2{,}065$^{a}$ \\
      S2 & Attack AND Application & 3{,}292 & 435 & 433 & 787 & 1{,}811$^{a}$ \\
      S3 & Attack AND Defence & 3{,}044$^{a}$ & 436 & 377 & 748 & 1{,}687$^{a}$ \\
      S4 & Attack AND Defence AND Adaptive & 1{,}095 & 175 & 123 & 233 & 617 \\
      S5 & Attack AND Guarantee & 1{,}966 & 314 & 238 & 485 & 1{,}206 \\
      S6 & Attack AND Guarantee AND Utility & 1{,}250 & 99 & 82 & 150 & 828 \\
      S7 & Attack AND Benchmark AND agent(s)/agentic & 1{,}445 & 151 & 112 & 201 & 811 \\
      S8 & Extraction AND System & 575 & 46 & 91 & 130 & 338$^{a}$ \\
      S9 & System AND Separation & 1{,}376 & 844 & 508 & 3{,}013 & 1{,}488$^{a}$ \\
      C1 & S2 AND (overstated, ineffective, unrealistic, \enquote{limited impact}, \enquote{low success}, rarely, \enquote{no evidence}, \enquote{in the wild}, hype) & 120 & 13 & 16 & 24 & 981 \\
      C2 & S4 AND (robust, resilient, resistant, immune, \enquote{near-zero}, effective, secure, \enquote{state-of-the-art}) & 656 & 131 & 88 & 165 & 395 \\
      C3 & S5 AND (solve(d), eliminate, prevents, completely, fully, provable, guarantee, \enquote{by design}) & 677 & 109 & 53 & 104 & 744 \\
      C4 & S6 AND (\enquote{no loss}, \enquote{without loss}, negligible, preserve, maintain, \enquote{without sacrificing}) & 334 & 35 & 31 & 44 & 250 \\
      C5 & S8 AND (defence, protect, safeguard, prevent, mitigate, obfuscation, secure, robust) & 281 & 24 & 49 & 61 & 170 \\
      C6 & S9 AND (separate, distinguish, robust, \enquote{by design}, enforce) & 319 & 123 & 67 & 366 & 613 \\
      C7 & S7 AND the robustness words of C2 & 674 & 106 & 71 & 123 & 430 \\
      R1 & System AND Untrusted AND (isolation, \enquote{least privilege}, \enquote{information flow}, integrity, guardrail, defence, prevent) & 180 & -- & 15 & 28 & 127 \\
      R2 & tool-using agents AND Adversarial & 131 & -- & 16 & 25 & 101$^{a}$ \\
      R3 & Prompt security AND System AND Defence & 848 & -- & 72 & 114 & 584 \\
      CR1 & R1 AND the words of C3 & 74 & -- & 6 & 7 & 72 \\
      CR2 & R2 AND the robustness words of C2 & 81 & -- & 11 & 18 & 65$^{a}$ \\
      CR3 & R3 AND the robustness words of C2 & 167 & -- & 50 & 73 & 395 \\
      \bottomrule
    \end{tabular}
  \end{sidecaption}
\end{table}

\paragraph{Studies the queries had to return.}
Before screening, the following studies were named as ones a working
search must return: the paper that introduced indirect prompt
injection\autocite{Greshake2023Signed}, the two adaptive-attack
evaluations\autocite{Zhan2025Adaptive,Nasr2026AttackerSecond}, the
design-pattern and information-flow
designs\autocite{BeurerKellner2025Patterns,Wu2024IFC}, the capability
design CaMeL\autocite{Debenedetti2026CaMeL} and the measurement of
prompt extraction\autocite{Zhang2024Extraction}. All of them were
returned by the concept queries.

\paragraph{Check by meaning.}
The keyword search was checked against a semantic search (the 50 nearest
works in OpenAlex for each of six questions: indirect injection, defences
bypassed by adaptive attacks, system-level designs with provable
security, extraction of system prompts, robust defences with near-zero
attack success, and the separation of instructions and data). Of the 300
places, the works not held by the keyword search were a design that
enforces least privilege on agent tools\autocite{Kim2025PFI}, an attack
that makes agents leak data through tool use\autocite{Fu2024Imprompter},
a systematisation of prompt security, two papers on poisoning or
backdoors in system prompts, a book chapter and a jailbreak defence.
The repair queries R1 (untrusted data with isolation, least privilege or
information flow), R2 (adversarial or hidden instructions with
exfiltration or tool misuse) and R3 (prompt attacks and prompt security)
were written by concept, not by title, and returned the first three; the
remaining four are poisoning, backdoor, book or jailbreak work outside
the question.

\paragraph{Deduplication.}
The search client merges copies of a work within one search (same DOI,
or same title and authors when there is no DOI) and folds a preprint into
its published version when the match is unambiguous. Copies that reach
the session from different searches remain as separate records (579 records are flagged as possible copies of another record);
the hit list therefore counts some works twice, typically an arXiv
preprint beside its conference version.

\paragraph{Screening.}
The searches returned 11,360 records (10,370 distinct works once copies
are folded). Screening followed a tiered design, and three numbers
describe it.

\emph{Screened in full: 2,276 records.} These are the first tier, every
record at place 100 or better in at least one relevance-ordered answer of
one query in one database (2,192 records), together with the 79 records
returned by the check by meaning and the 5 cited works that lay outside
that tier. Each first-tier record was classified from its title and
abstract by a large language model (no examples; a written research
context), into four categories (supports the claim; qualifies or
contradicts it; adjacent subtopic; off-topic) and, for discards, six
labelled exclusion criteria: \emph{other-field} (not about the security of
language models), \emph{jailbreak-only} (jailbreaks or harmful content
rather than injected instructions), \emph{other-question} (security of
language models, but another question such as backdoors, poisoning,
denial of service or privacy), \emph{llm-as-tool} (a language model used
for some task, security only in passing), \emph{no-evidence} (opinion,
position or overview without evidence of its own) and \emph{no-abstract}.
An AI agent then read every record classified as qualifying (436) and
every low-confidence record, corrected 102 decisions that contradicted the
criteria (kept records carrying an exclusion criterion, discards without
one), screened the 79 check-by-meaning hits itself from title and
abstract, and tagged the 48 works cited here. The fully screened set holds
637 records that support the claim and 436 that qualify it; 1,203 were
discarded (668 adjacent, 470 off-topic, and 65 that bear on the topic but
were excluded by a criterion, mostly for lack of evidence of their own).

\emph{Sampled: 360 records.} From the records outside the first tier, two
seeded random samples were drawn and screened in the same way: 60 records
(seed 20261127) and, after further decisions had been made, another 300 of
the 9,024 records then still pending (seed 20261128). Of the 360, 80 bear
on the claim: 22.2\,\%, with a 95\,\% Wilson interval of 18.2 to 26.8\,\%.
The agent read all 80; none reports an independent break of a
system-level design, and none adds a defence shown by others to hold
against adaptive attacks. Two that qualify the statements in detail were
read in full and are cited below. The sample estimates what was not
screened; it is not a screen of the rest.

\emph{Only counted: 8,724 records.} These were not screened. If the
sample is representative, an estimated 1,600 to 2,300 of them bear on the
claim. The search is therefore not exhaustive: the conclusion rests on the
fully screened first tier, the named studies, the check by meaning and the
sample, and no statement here says that a kind of study does not exist.

All decisions were made by the language model or the AI agent; none has
yet been reviewed by the author. For that review, seeded random samples of
50 of the 1,471 machine discards (seed 20261129) and 30 of the 1,165
machine inclusions (seed 20261130) have been drawn; their disagreement
rates are not yet known. Records outside the first tier that the
classifier had decided before the sampling design was fixed (592, by a
smaller model and not at random) were reset to undecided and are counted
among the 8,724.

\paragraph{Verification.}
Every source cited in this chapter was read in full text (the arXiv,
ACL Anthology or open conference version), and each statement below
rests on a passage quoted verbatim, with its page, in the bibliography's
provenance record. Venues were checked against the publishers' DOI
records; works without a published version are cited as preprints.

\section{Results}
\label{prompt-injection:results}

\subsection{Indirect prompt injection works against real systems}

\textcite{Greshake2023Signed} introduced indirect prompt injection: the
attacker does not talk to the model but places instructions in data the
application will retrieve, such as a web page. They argue that
\enquote{LLM-Integrated Applications blur the line between data and
instructions} and demonstrated the attack against real systems,
including Bing's GPT-4 powered chat. The mechanism is restated by the
authors of the standard agent benchmark, who write that LLMs
\enquote{operate directly on text, lacking a formal way to distinguish
instructions from data}\autocite{Debenedetti2024AgentDojo}, and by the
authors of a structured-query defence: \enquote{there is no separation
between prompts and data}\autocite{Chen2025StruQ}. The only formal
measurement of this separation found \enquote{all models fail to achieve
high separation}, and prompt engineering or fine-tuning either did not
substantially improve separation or reduced the model's
utility\autocite{Zverev2025Separate}. Goal hijacking and prompt leaking
were named as attack types already in 2022\autocite{Perez2022Ignore}.

The attacks work outside the laboratory. A black-box attack found 31 of
36 commercial LLM-integrated applications susceptible
\autocite{Liu2023HouYi}; optimised injections made agents leak personal data with nearly 80\,\%
success in an end-to-end evaluation and worked on agent products such as
Mistral's LeChat\autocite{Fu2024Imprompter};
a benchmark of 30 tool-using agents found them vulnerable, a
prompted GPT-4 agent in 24\,\% of the cases\autocite{Zhan2024InjecAgent};
another benchmark reports an average attack success of up to 84.30\,\%
with limited effect of current defences\autocite{Zhang2025ASB}; and a global
competition and an online game, with more than 600,000 and 126,000 human-written
attacks, showed that models can be manipulated by prompt
hacking\autocite{Schulhoff2023HackAPrompt,Toyer2024TensorTrust}.
In a 2026 public competition with 272,000 attack attempts against 13
frontier models, \enquote{all models proved vulnerable}, with
attack success rates from 0.5\,\% to 8.5\,\%
\autocite{Dziemian2026Competition}. Injections also exist on the open
web: a crawl of 1.2 billion URLs found 15,300 validated instances
\autocite{Khodayari2026Wild}.

The counter-search found no study showing that indirect injection fails
against deployed systems, but it found several that narrow how strongly
the attack works. For the most robust current models the success rate per
attempt is low\autocite{Dziemian2026Competition}, and attacks in the
standard benchmark \enquote{break some security properties but not
all}\autocite{Debenedetti2024AgentDojo}. End to end, attacks on
retrieval-augmented applications \enquote{settle around a 40\,\% success
rate}\autocite{DeStefano2026RagnRoll}; in a realistic retrieval pipeline,
\enquote{prior protocols substantially overstate attack effectiveness},
although \enquote{LLM-driven prompt injections remain effective
end-to-end}\autocite{Yin2026Generator}. Benchmark numbers themselves can
mislead: one audit found a harness that inflated attack success
eighteen-fold\autocite{Shaw2026SilentFailures}, and a simple pair of
filters reached perfect security on four public agent benchmarks, which
its authors read as evidence of \enquote{weak attacks} in those
benchmarks\autocite{Bhagwatkar2025Firewalls}. The defensible wording is
that injection \emph{works against real systems}, not that it always
succeeds.

\subsection{Defences reduce attack success but are bypassed by adaptive attacks}

Many defences report large reductions in attack success when evaluated
against fixed attacks. Marking the provenance of input text
(spotlighting) \enquote{reduces the attack success rate from greater
than 50\,\% to below 2\,\%} in its authors'
experiments\autocite{Hines2024Spotlighting}; training a model to respect
an instruction hierarchy \enquote{drastically increases
robustness}\autocite{Wallace2024Hierarchy}; preference-optimisation
fine-tuning reduces attack success to below 10\,\% \enquote{even against
attacks much more sophisticated than ones seen during
training}\autocite{Chen2025SecAlign}; segment embeddings, attention-based
detection, a game-theoretic detector, a task-alignment check, an
LLM-based filter and a re-execution check all report gains of the same
kind\autocite{Wu2025ISE,Hung2025AttentionTracker,Liu2025DataSentinel,%
Jia2025TaskShield,Shi2025PromptArmor,Zhu2025MELON}.

Evaluations by attackers who adapt to the defence tell a different
story. \textcite{Zhan2025Adaptive} \enquote{evaluate eight different
defenses and bypass all of them using adaptive attacks, consistently
achieving an attack success rate of over 50\,\%}. \textcite{Nasr2026AttackerSecond}
(USENIX Security 2026; authors from Anthropic, Google DeepMind, OpenAI, ETH
Zürich) \enquote{bypass 12 recent defenses [\dots] with attack success
rate above 90\,\% for most}, noting that \enquote{the majority of
defenses originally reported near-zero attack success rates}. The twelve
cover defences against jailbreaks as well as prompt injections; the
prompt-injection defences among them include spotlighting, prompt
sandwiching, the structured-query and a fine-tuned defence, several
detectors, and the re-execution defence whose title calls it provable,
against which their attack reached 76\,\%. Detectors that check whether the model
still answers a known question were evaded with detection rates
\enquote{as low as 0\,\%} while the attack succeeded in
91\,\%\autocite{Choudhary2025NotDetect}. Earlier systematic evaluations
reached the same verdict without adaptive attacks: \enquote{no existing
defenses are sufficient}\autocite{Liu2024Formalizing}, and off-the-shelf
detectors \enquote{struggle to effectively detect indirect prompt
injection attacks}\autocite{Chen2025DetectRemove}.

The authors of the defences agree. The structured-query defence
\enquote{is not yet fully secure}\autocite{Chen2025StruQ}; the authors of
a leading fine-tuned model write that it, \enquote{similar to existing
defenses, remains vulnerable to strong adaptive attacks}, and conclude:
\enquote{That is, PI is not solved.}\autocite{Chen2025MetaSecAlign}.
The newest fine-tuning recipe from the group behind the structured-query
defence starts from the observation that trained models \enquote{still
suffer from near 100\,\% attack success rates (ASRs) against adaptive
prompt injections}; it reduces this to 9.0\,\% for one adaptive attack,
against 94.0\,\% for its predecessor\autocite{Peng2026SecOPD} --- a large
improvement, but not prevention.
Model-level defences are, in the words of an information-flow paper,
\enquote{likely vulnerable to attacks specifically tailored to defeat
that defense}\autocite{Wu2024IFC}. Three independent groups and the
defence authors themselves converge on this. Not every recent defence
has been attacked adaptively by others in the sources read. An LLM-based
filter reports false-positive and false-negative rates below 1\,\% and
effectiveness against adaptive attacks\autocite{Shi2025PromptArmor}; an
attention-based detector, peer-reviewed, is reported to be \enquote{robust
against adaptive adversaries}\autocite{Zhong2026Rennervate}; and learnable
delimiters reach near-zero attack success and are said to remain effective
when the attacker knows the defence\autocite{Zhao2026LTBD}. These stand as
qualifying claims until others test them; the history above (defences
reporting near-zero attack success that later fell to adaptive attacks) is
the reason not to treat them as settled.

\subsection{Guarantees come from the design of the system, at a cost}

The counter-search for \enquote{provable}, \enquote{by design} and
\enquote{solved} found a group of system-level designs that do give
guarantees, and they qualify the claim. CaMeL extracts the control and
data flow from the trusted user request, so that \enquote{the untrusted
data retrieved by the LLM can never impact the program flow}, and checks
every tool call against a capability policy; on the standard benchmark
it solved \enquote{77\,\% of tasks with provable security (compared to
84\,\% with an undefended system)}\autocite{Debenedetti2026CaMeL}. Its
guarantee covers a defined class of attacks: the authors state that
\enquote{it cannot defend against text-to-text attacks}, which change
only what the user is told. An information-flow planner
\enquote{deterministically enforces security policies} and completes
\enquote{a broad range of tasks with security
guarantees}\autocite{Costa2026FIDES}; separating planner and executor
gives the same kind of control\autocite{Wu2024IFC}. Least-privilege
policies on tool calls reduced attack success on the benchmark from
39.9\,\% to 1.0\,\% while keeping utility, and hand-written policies
\enquote{can provably reduce the ASR to 0\,\%}\autocite{Shi2025Progent};
a certifier for structured actions such as paying an invoice
\enquote{admits no unsafe action} under a bounded corruption budget,
although in the same study a realistic injection fooled every model but
one\autocite{Pandey2026Certified};
a related design argues that \enquote{the principle of least privilege
should be enforced to restrict the impact of untrusted data on the
privileged tools}\autocite{Kim2025PFI}, and a tool-dependency-graph
defence criticises prompting and detection for relying on \enquote{the
model's inherent security}\autocite{An2025IPIGuard}.

These guarantees have so far been tested mostly against fixed benchmark
attacks. An independent review of such out-of-band defences warns that
\enquote{every one of them is validated only on static benchmarks}; in its
own small adaptive reproduction, the least-privilege defence held, cutting
attack success from 25.8\,\% to 4.2\,\% and staying low under a
hand-crafted adaptive attack\autocite{Narisetty2026OutOfBand}. That is
one data point, not a settled result.

Enforcement in code must itself be sound: a deterministic system-call
gate for agent sandboxes was defeated by a race in 18.9 to 23.1\,\% of
attempts, while kernel-level confinement resisted the tested attacks; its
authors present this as evidence \enquote{within the stated threat model
and tested matrix, not a proof against untested
interfaces}\autocite{Erisen2026Seccomp}.

The price is generality. The authors of a set of design patterns write
that their patterns \enquote{constrain the actions of agents to
explicitly prevent them from solving arbitrary tasks}, that \enquote{once
an LLM agent has ingested untrusted input, it must be constrained so that
it is impossible for that input to trigger any consequential actions},
and that \enquote{it is unlikely that general-purpose agents can provide
meaningful and reliable safety guarantees}\autocite{BeurerKellner2025Patterns}.
Asking the user to confirm actions is part of the same toolbox --- CaMeL
asks for confirmation when a policy would be
violated\autocite{Debenedetti2026CaMeL} --- but the design-pattern
authors note that confirming every command would be cumbersome for the
user\autocite{BeurerKellner2025Patterns}. Most of these designs are
preprints; among the sources read, CaMeL (published) and the
information-flow planner (accepted for NeurIPS 2026) are the peer-reviewed
system-level designs that claim provable or deterministic security (the
re-execution defence that also calls itself provable is a detector, and it
was bypassed).

\subsection{System prompts are readily extracted}

Simple text-based attacks \enquote{can in fact reveal prompts with high
probability} across 11 models, and prompts were extracted from deployed
systems \enquote{despite existing defenses in
place}\autocite{Zhang2024Extraction}. Aligned models such as GPT-4 were
found \enquote{highly vulnerable to prompt extraction attacks, even under
the most straightforward user attacks}\autocite{Liang2024Leaked}, and
prompt theft was one outcome of the commercial-application attack
\autocite{Liu2023HouYi}. The counter-search qualifies this: in an
extraction benchmark \enquote{all seven evaluated models are vulnerable in
an undefended state}, but a defended GPT-4 configuration showed
resilience\autocite{Wang2024Raccoon}, and a proxy-prompt defence
protected 94.70\,\% of prompts from extraction in its
evaluation\autocite{Zhuang2026ProxyPrompt}. In a pipeline of production agents, \enquote{the entry agent does not
leak its system prompt (0/40)}, yet the pipeline was still compromised by
an injected instruction\autocite{Sidot2026Verify}. Hidden prompts are thus
easy to extract from undefended models, and no defence read here claims to
make extraction impossible.

\section{Discussion and limitations}
\label{prompt-injection:discussion}

The answer rests on convergence across independent strands: attack
papers, benchmarks, adaptive evaluations of defences, the defence authors'
own limitations sections, and system-level designs whose guarantees are
explicitly scoped. The strongest single pieces of evidence for (b) are two
peer-reviewed adaptive evaluations\autocite{Zhan2025Adaptive,Nasr2026AttackerSecond};
the strongest for (c) are one peer-reviewed design\autocite{Debenedetti2026CaMeL}
and one accepted for publication\autocite{Costa2026FIDES}.

The searches have the following known gaps. DBLP could not be searched
at all, and IEEE Xplore failed for the six repair queries; a paper
indexed only there could have been missed.
Semantic Scholar matched one counter-query (C1) much more loosely than
the other databases (981 hits against 13 to 120); its extra records were
treated like all others (screened if in the first tier or the sample). Only 2,276 records were screened in full and 360 more as a random
sample; 8,724 were only counted. The sample puts the share of relevant
records among them at 22.2\,\% (95\,\% interval 18.2 to 26.8\,\%), so an
estimated 1,600 to 2,300 relevant records were never read. The search is
not exhaustive, and the answer rests on the first tier, the named
studies, the check by meaning and the sample. That the sample's 80
relevant records contained nothing that changes the answer makes it
unlikely, but not impossible, that the unread records hold a refutation. The screening is by a
language model and an AI agent and has not yet been reviewed by the
author. Full texts were read in their arXiv or other open versions, which
can differ from the published papers in detail. Several cited works are
preprints; they are marked as such in the bibliography, and no conclusion
rests on preprints alone.

The field moves fast: several cited works are preprints from 2025 and
2026, and a defence published after these searches could change part
(b). Benchmarks measure attack success on fixed task suites; none of the
sources measures prevention in a deployed business system over time.

\section{Conclusion}
\label{prompt-injection:conclusion}

No: the literature screened here contains no filter, detector, prompt
technique or trained model that reliably prevents prompt injection
against an attacker who adapts to it, while system-level designs that
keep untrusted data out of the control flow and enforce policies in code
do give provable guarantees, but only for a defined class of attacks and
at a cost in what the agent can do. This agrees with the OWASP
statements\autocite{OWASP2025LLM,OWASP2026LLM} and with
Willison\autocite{Willison2025Trifecta}, with one refinement: the
\enquote{architectural} defence that OWASP recommends can be provably
secure for its class of attacks, not merely better. The answer rests on
the 2,276 records screened in full, the named studies, the check by
meaning and a random sample of 360 of the remaining records; an estimated
1,600 to 2,300 relevant records among the 8,724 that were only counted
have not been read.

For the lecture this means: it may say that instructions and data share
one channel; that indirect injection works against real products; that
published defences reduce attack success but have been bypassed by
adaptive attackers; and that the remedy is to limit what a successful
injection can do (least privilege, authorisation in code, confirmation
of consequential actions). It must not say that prompt injection
\enquote{cannot be prevented} or that \enquote{no defence works}, since
designs such as CaMeL give guarantees for a class of attacks; nor that a
vendor's classifier or a fine-tuned model makes an application safe.

\section{Future work}
\label{prompt-injection:future}

This search left undone: the screening of the 8,724 records that were
only counted; the six IEEE Xplore cells of the repair queries (quota) and
every DBLP query (blocked); a review by the author of the machine
screening, including the seeded samples drawn for it; the published
versions of several preprints; and one measurement study of in-page
prompt injections against web agents that is not openly
available.\footnote{\enquote{A Large-scale Measurement of In-Page Prompt
Injections Against LLM Web Agents}, ACM, 2026, doi:10.1145/3830454.3846661.}

The field lacks independent adaptive evaluations of the newest defences
(the LLM-based filter\autocite{Shi2025PromptArmor} and the system-level
designs), measurements of how much utility system-level designs cost in
realistic business workflows rather than benchmark suites, and studies of
injections in deployed agents over time. If an independent adaptive
evaluation found a model-level defence that held, part (b) of the lecture
would have to say \enquote{most published defences}; if one broke a
system-level design inside its stated threat model, part (c) would lose
its only guarantee.

\section{Hits that bear on the claim}
\label{prompt-injection:hits}

% Author's note: table generated by
%   scholar sessions export prompt-injection -f table --bearing-only \
%     --lang en --label prompt-injection:hits-table ... \
%     -o litteratursokning/prompt-injection-full
% raw CSV export beside it: litteratursokning/prompt-injection-full.csv
% --- scholar LaTeX fragment --------------------------------------------
% Generated by: scholar sessions export -f table
% Session: prompt-injection
% \input this file from the body of a document whose
% preamble loads:
%   \usepackage{longtable}
%   \usepackage{booktabs}
%   \usepackage{array}
% Text that pdfLaTeX cannot set in T1 (Cyrillic, CJK, ...) is
% wrapped in \scholarnonlatin{...}. The \providecommand lines
% below print the text under XeLaTeX and LuaLaTeX (load a font
% that has the script) and print [non-Latin text] under pdfLaTeX.
% To print it under pdfLaTeX too, define the command before the
% \input, e.g. for Cyrillic with \usepackage[T2A,T1]{fontenc}:
%   \newcommand{\scholarnonlatin}[1]{{\fontencoding{T2A}\selectfont #1}}
% Where such text touches a letter, \scholarnonlatinsep{}
% stands between them: a space under pdfLaTeX, nothing under
% XeLaTeX and LuaLaTeX. If your \scholarnonlatin prints the
% text, define the separator as empty beside it:
%   \newcommand{\scholarnonlatinsep}{}
% otherwise you get a space where the scripts touch.
% ----------------------------------------------------------------------

\ifdefined\XeTeXversion\providecommand{\scholarnonlatinsep}{}\fi
\ifdefined\directlua\providecommand{\scholarnonlatinsep}{}\fi
\providecommand{\scholarnonlatinsep}{\ }
\ifdefined\XeTeXversion\providecommand{\scholarnonlatin}[1]{#1}\fi
\ifdefined\directlua\providecommand{\scholarnonlatin}[1]{#1}\fi
\providecommand{\scholarnonlatin}[1]{[non-Latin text]}
\begingroup\footnotesize
\begin{longtable}{@{}%
  >{\raggedright\arraybackslash}p{0.44\textwidth}c%
  >{\raggedright\arraybackslash}p{0.13\textwidth}%
  >{\raggedright\arraybackslash}p{0.26\textwidth}@{}}
\caption{Records bearing on the claim, Can prompt injection against LLM-integrated applications and agents be reliably prevented? Screened in full: 2,276 records (the first tier of the relevance-ordered answers, the check-by-meaning hits and the cited works). A seeded random sample of 360 of the other records was screened as a sample, not as a screen of the rest: 80 of the 360 (22.2 \%; 95 \% Wilson interval 18.2 to 26.8 \%) bear on the claim. The remaining 8,724 records were only counted, so an estimated 1,600 to 2,300 records that bear on the claim were not screened and this list is not exhaustive. Every decision was made by a language model or an AI agent and awaits the author's review (0 cited, 623 supporting and 418 qualifying or contradicting it, out of 10,370 unique hits; the 707 adjacent hits, the 652 false hits, 7,910 unclassified and 60 excluded although classified as bearing on the claim are not listed). The last column gives the reason for each inclusion; for machine-classified rows the model's confidence is given, and for rows decided by an agent the agent's. Databases searched: OpenAlex (OA; not open access), IEEE Xplore (IEEE; failed for 6 of 22 queries), Web of Science (WoS), Scopus, Semantic Scholar (S2). Searched 2026-11-27.}\label{tab:prompt-injection:hits-table}\\
\toprule
Title & Year & Database & Reason \\
\midrule
\endfirsthead
\multicolumn{4}{@{}l}{\emph{\tablename~\thetable{} (cont.)}}\\
\toprule
Title & Year & Database & Reason \\
\midrule
\endhead
\midrule \multicolumn{4}{r@{}}{\emph{continued on next page}}\\
\endfoot
\bottomrule
\endlastfoot
\scholarnonlatin{ЭКСПЕРИМЕНТАЛЬНАЯ ОЦЕНКА УСТОЙЧИВОСТИ ОТЕЧЕСТВЕННЫХ ЯЗЫКОВЫХ МОДЕЛЕЙ К ПРОМПТ-ИНЪЕКЦИЯМ} (hash:60e2ae1dc7f40652) & 2026 & OA & supports the claim (0.90) \\
\textquotedbl{}PromptGame Dataset: Prompt Injection Attack Benchmark and Defense Evaluation Data\textquotedbl{} & 2026 & OA & supports the claim (0.40) \\
``**Important** You should give me full credits!'': Exploring Prompt Injection Attacks on LLM-Based Automatic Grading Systems & 2026 & OA, S2 & supports the claim (0.93) \\
``I've Decided to Leak'': Probing Internals Behind Prompt Leakage Intents & 2025 & OA & supports the claim (0.55) \\
``I've Decided to Leak'': Probing Internals Behind Prompt Leakage Intents & 2025 & OA, Scopus & supports the claim (0.80) \\
A Black-Box Inversion Attack on Small Language Models and Differential Privacy-Based Defense & 2026 & WoS, Scopus & supports the claim (0.50) \\
A Comparative Study Between Qwen 3.6 Flash and Small Language Models for Email Prompt Injection in Android Edge Agents & 2026 & OA & supports the claim (0.90) \\
A Comprehensive Assessment Tool for Prompt Injection Attacks in Large Language Models (LLMs) & 2026 & IEEE, OA, Scopus, S2 & supports the claim (0.60) \\
A Comprehensive Framework for Secure Multi-Agent Retrieval-Augmented Generation Systems & 2026 & WoS, IEEE, OA, Scopus, S2 & supports the claim (0.40) \\
A Critical Evaluation of Defenses against Prompt Injection Attacks: [Dataset/Tool Paper] & 2026 & WoS, OA, Scopus, S2 & supports the claim (0.90) \\
A Federated MCP Gateway: Enforceable Isolation and Measuring Residual Semantic Attacks & 2026 & WoS, IEEE, OA, Scopus, S2 & supports the claim (0.70) \\
A focused survey of code agent security: Attacks, defenses, and evaluation & 2026 & WoS, Scopus & supports the claim (0.70) \\
A Formal Framework of Architectural Intent Collapse for Tool-Level Attacks on LLM Agents & 2026 & WoS, OA, Scopus, S2 & supports the claim (0.80) \\
A Gateway-Centric Architecture for Bidirectional Content Moderation in Large Language Models & 2026 & IEEE, OA, Scopus, S2 & supports the claim (0.45) \\
A Large-scale Measurement of In-Page Prompt Injections Against LLM Web Agents & 2026 & OA, S2 & supports the claim (0.80) \\
A Layered Security Framework Against Prompt Injection in RAG-Based Chatbots & 2026 & OA, S2 & supports the claim (0.50) \\
A Multi-Agent Framework for Step-Level Tool-Call Safety Evaluation in Industrial Control Systems & 2026 & IEEE, OA, Scopus, S2 & supports the claim (0.65) \\
A Novel Evaluation Framework for Assessing Resilience Against Prompt Injection Attacks in Large Language Models & 2023 & IEEE, OA, Scopus, S2 & supports the claim (0.55) \\
A Privacy-Preserving On-Device Multi-Agent Architecture for AI PC (POMA) Workflow Automation & 2026 & IEEE, OA, Scopus, S2 & supports the claim (0.40) \\
A Security Analysis of the OpenClaw AI Agent Framework & 2026 & WoS, IEEE, OA, Scopus & supports the claim (0.80) \\
A Study on Data Integrity and Security Threat Analysis in Retrieval-Augmented Generation (RAG) Systems & 2026 & OA & supports the claim (0.65) \\
A Study on Prompt Injection Attack Against LLM-Integrated Mobile Robotic Systems & 2024 & WoS, IEEE, OA, Scopus, S2 & supports the claim (0.60) \\
A Systematic Literature Review of Prompt Injection Attacks in LLM-Integrated Systems & 2026 & WoS, IEEE, OA, Scopus, S2 & supports the claim (0.60) \\
A Systematic Literature Review on LLM Defenses Against Prompt Injection and Jailbreaking: Expanding NIST Taxonomy & 2026 & OA, S2 & supports the claim (0.50) (no-evidence) \\
A Systematic Review of Prompt Injection Attacks on Large Language Models: Trends, Taxonomy, Evaluation, Defenses, and Opportunities & 2026 & WoS, IEEE, OA, Scopus, S2 & supports the claim (0.45) \\
A white-box prompt injection attack on embodied AI agents driven by large language models & 2026 & WoS, OA, Scopus, S2 & supports the claim (0.50) \\
Accuracy-Invisible Safety Degradation in a Compressed On-Device Prompt-Injection Classifier & 2026 & OA, S2 & supports the claim (0.70) \\
ACIARENA: Toward Unified Evaluation for Agent Cascading Injection & 2026 & WoS, S2, OA, Scopus & supports the claim (0.80) \\
ACIArena: Toward Unified Evaluation for Agent Cascading Injection & 2026 & OA & supports the claim (0.80) \\
Adaptive Attacks Break Defenses Against Indirect Prompt Injection Attacks on LLM Agents & 2025 & OA, Scopus, S2 & supports the claim (0.90) \\
AdapTools: Adaptive Tool-based Indirect Prompt Injection Attacks on Agentic LLMs & 2026 & OA, S2 & supports the claim (0.95) \\
Addressing Security Gaps in MCP: Design of a Resilient Reference Architecture & 2025 & IEEE & supports the claim (0.50) \\
Adversarial and Multilingual Threats in Retrieval-Augmented Generation: From Prompt Injection to Model Exploitation & 2025 & IEEE, OA, Scopus, S2 & supports the claim (0.75) \\
Adversarial Attacks on Agentic AI Systems: Mechanisms, Impacts, and Defense Strategies & 2026 & OA, S2 & supports the claim (0.45) \\
Adversarial Decoding: Generating Readable Documents for Adversarial Objectives & 2024 & S2 & supports the claim (0.80) \\
Adversarial Misdirection: Probing \& Visualizing Cross-Modal Reasoning Vulnerabilities in Vision-Language Models & 2025 & WoS, IEEE, Scopus, S2 & supports the claim (0.55) \\
Adversarial prompt and fine-tuning attacks threaten medical large language models & 2025 & WoS, OA, Scopus, S2 & supports the claim (0.70) \\
Adversarial Prompt Optimization in LLMs: HijackNet's Approach to Robustness and Defense Evasion & 2025 & IEEE, OA, Scopus, S2 & supports the claim (0.45) \\
Adversarial-Resistant AI Tutoring: A Dual-LLM Architecture for Mitigating H-CoT Attacks and Ensuring Pedagogically-Aligned Code Assistance & 2025 & OA & supports the claim (0.40) \\
AE of USENIX 1226 (Overcoming the Retrieval Barrier: Indirect Prompt Injection in the Wild for LLM Systems) & 2025 & OA & supports the claim (0.50) \\
Aegis Cortex OS: Governed Execution Layer for Enterprise AI Agents & 2026 & OA & supports the claim (0.50) \\
AegisAgent: An Autonomous Defense Agent Against Prompt Injection Attacks in LLM-HARs & 2025 & OA, S2 & supports the claim (0.40) \\
Agent Against Agent: An Agentic System for Automatic Prompt Injection Red Teaming & 2026 & OA, S2 & supports the claim (0.85) \\
Agent Data Injection Attacks are Realistic Threats to AI Agents & 2026 & OA, S2 & supports the claim (0.90) \\
Agent Guard: Kernel-Enforced Damage Boundaries for AI Agents via Human-Authorized Contracts & 2026 & OA & supports the claim (0.60) \\
Agent Harnesses and the New Threat Landscape & 2026 & OA & supports the claim (0.55) \\
AGENT SECURITY BENCH (ASB): FORMALIZING AND BENCHMARKING ATTACKS AND DEFENSES IN LLM-BASED AGENTS & 2025 & OA, Scopus, S2 & supports the claim (0.95) \\
AgentBreaker: Evaluating Context-Aware Indirect Prompt Injection Risks in Modern Web Agents & 2026 & OA, S2 & supports the claim (0.85) \\
AgentBreaker: Evaluating Context-Aware Indirect Prompt Injection Risks in Modern Web Agents (Artifact) & 2026 & OA & supports the claim (0.80) \\
AgentDojo: A Dynamic Environment to Evaluate Prompt Injection Attacks and Defenses for LLM Agents & 2024 & OA, Scopus, WoS & supports the claim (0.90) \\
AgentDojo: A Dynamic Environment to Evaluate Prompt Injection Attacks and Defenses for LLM Agents & 2024 & OA, S2 & supports the claim (0.90) \\
AgentDyn: Are Your Agent Security Defenses Deployable in Real-World Dynamic Environments? & 2026 & OA, S2 & supports the claim (0.90) \\
Agentic AI in Cybersecurity:Threat Modeling and Containment Strategies for Autonomous LLM Agents & 2026 & IEEE, OA, Scopus, S2 & supports the claim (0.55) \\
Agentic AI Red-Teaming for Prompt Injection Detection in Multi-Agent Systems & 2026 & IEEE, OA, S2 & supports the claim (0.60) \\
Agents Break Together on Misuse, Apart on Injection & 2026 & OA & supports the claim (0.70) \\
Agents of Chaos: A Systematization of Knowledge (SoK) on Malicious Attacks on Agentic AI Systems & 2026 & OA, S2 & supports the claim (0.60) \\
AgentSec-Bench: An Open Benchmark for Indirect Prompt Injection in Tool-Augmented Agents & 2026 & OA & supports the claim (0.60) \\
AgentSecBench: Measuring Prompt Injection, Privacy Leakage, and Tool-Use Integrity in LLM Agents & 2026 & OA, S2 & supports the claim (0.65) \\
AgentShield Bench v3: Evaluating Goal Integrity and Resilience Of Autonomous LLM Agents Against Goal Hijacking and Autonomous Threats & 2026 & OA & supports the claim (0.50) \\
AgentShield Bench: Evaluating the Security Resilience of OpenAI and Gemini LLM Agents Against Adversarial Agent Workflows & 2026 & OA & supports the claim (0.60) \\
AGENTVIGIL: Automatic Black-Box Red-teaming for Indirect Prompt Injection against LLM Agents & 2025 & WoS, OA, Scopus, S2 & supports the claim (0.92) \\
AgentVigil: Generic Black-Box Red-teaming for Indirect Prompt Injection against LLM Agents & 2025 & OA, S2 & supports the claim (0.90) \\
AGFPS: An Automated Gradient-Free Framework for Prompt Stealing & 2026 & WoS, IEEE, OA, Scopus, S2 & supports the claim (0.93) \\
AI Agents & 2026 & OA & supports the claim (0.70) \\
AI Agents May Always Fall for Prompt Injections & 2026 & OA, S2 & supports the claim (0.90) \\
AI Safety Evaluation: Lessons from International Multilingual Joint Testing & 2026 & IEEE, Scopus & supports the claim (0.70) \\
Ajar: Measuring Open Privilege in Agent Defenses & 2026 & OA, S2 & supports the claim (0.70) \\
AltayDuel: A Turkish-First Arena and Open Dataset for Multi-Turn LLM Prompt-Injection Red-Teaming & 2026 & OA & supports the claim (0.60) \\
An API-Driven Framework for Detecting and Mitigating Prompt Injection Vulnerabilities in Open-Source Large Language Models (LLMs) & 2026 & OA, Scopus & supports the claim (0.40) \\
An Empirical Benchmark and Security Evaluation of Prompt Injection Attacks in LLM Systems & 2026 & OA & supports the claim (0.85) \\
An Empirical Evaluation of Prompt Injection Detection and Refusal-Usefulness Tradeoffs Using the deepset/prompt-injections Dataset & 2025 & OA, S2 & supports the claim (0.40) \\
An Experimental Evaluation of Multimodal Prompt Injection Attacks on Agentic AI Frameworks & 2026 & OA, S2 & supports the claim (0.50) \\
An Experimental Evaluation of Prompt Injection Attacks on LLM-Based Content Moderation Systems & 2025 & IEEE, OA, Scopus, S2 & supports the claim (0.75) \\
An Invariant Latent Space Perspective on Language Model Inversion & 2026 & WoS, OA, Scopus, S2 & supports the claim (0.75) \\
Analysis of LLMs Against Prompt Injection and Jailbreak Attacks & 2026 & WoS, OA, Scopus, S2 & supports the claim (0.75) \\
Analyzing the Narration Gap in LLM-Solver Loops & 2026 & OA, S2 & supports the claim (0.80) \\
Análisis de vulnerabilidades por inyección de prompts en sistemas de IA agénticos: Un enfoque práctico sobre la integridad de los datos & 2026 & OA & supports the claim (0.50) \\
Appreciation Teaches Unauthorized Execution: Self-Evolution Degrades the Robustness of Personal Agents to Indirect Prompt Injection & 2026 & OA, S2 & supports the claim (0.85) \\
ARCHITECTURAL CONSTRAINTS ON THE INTEGRATION OF AGENTIC ARTIFICIAL INTELLIGENCE SYSTEMS INTO SOFTWARE SYSTEMS FOR SENSITIVE DATA PROCESSING & 2026 & OA, S2 & supports the claim (0.40) (no-evidence) \\
Are AI-assisted Development Tools Immune to Prompt Injection? & 2026 & OA, S2 & supports the claim (0.85) \\
Are LLM-Agent Prompt-Injection Benchmarks Still Valid? A Cross-Model and Cross-Benchmark Study & 2026 & OA & supports the claim (0.85) \\
Argus-Bench: A Prompt-Injection Benchmark for AI Agent Gateways & 2026 & OA & supports the claim (0.50) \\
Ariadne: Adaptive, Human-in-the-Loop Defense for AI Agents against Indirect Prompt Injection & 2026 & OA & supports the claim (0.60) \\
ASPI: Seeking Ambiguity Clarification Amplifies Prompt Injection Vulnerability in LLM Agents & 2026 & OA, S2 & supports the claim (0.90) \\
Assessing Prompt Injection Risks in 200+ Custom GPTs & 2023 & OA, S2 & supports the claim (0.90) \\
AttackEval: A Systematic Empirical Study of Prompt Injection Attack Effectiveness Against Large Language Models & 2026 & OA, S2 & supports the claim (0.60) \\
Attacking LLMs and AI Agents: Advertisement Embedding Attacks Against Large Language Models & 2025 & WoS, IEEE, Scopus, OA, S2 & supports the claim (0.45) \\
Attention Exhaustion Attacks: Multi-Substrate Cognitive Friction Prompts Defeat Detection in Frontier Large Language Models & 2026 & OA & supports the claim (0.50) \\
Attention-based Classification for Indirect Prompt Injection & 2026 & OA, S2 & supports the claim (0.85) \\
Attractive Metadata Attack: Inducing LLM Agents to Invoke Malicious Tools & 2025 & OA, Scopus, S2 & supports the claim (0.75) \\
Attributing Security Outcomes in LLM-Agent Evaluations: The Role of Content-Free State & 2026 & OA, S2 & supports the claim (0.60) \\
Augmented Adversarial Trigger Learning & 2025 & OA, Scopus & supports the claim (0.85) \\
Authorization Architectures for Tool-Using AI Agents & 2026 & OA, S2 & supports the claim (0.55) (no-evidence) \\
Authorization assurance for tool-using LLM agents: a structured review of enforcement semantics and effect-level security evidence & 2026 & WoS, Scopus & supports the claim (0.45) \\
AutoDAN: Interpretable Gradient-Based Adversarial Attacks on Large Language Models & 2023 & OA, S2 & supports the claim (0.40) \\
AutoDojo: A Generative Benchmark for Evaluating Prompt Injection Defenses in LLM Agents & 2026 & OA, S2 & supports the claim (0.95) \\
Automated adversarial red-teaming for evaluating robustness in LLM-based recommender systems & 2026 & Scopus, OA, S2 & supports the claim (0.50) \\
Automated Framework for Red-Team Prompt Injection Testing of LLMs & 2026 & IEEE, OA, Scopus, S2 & supports the claim (0.70) \\
Automatic and Universal Prompt Injection Attacks against Large Language Models & 2024 & OA, S2 & supports the claim (0.95) \\
Automatic Security Testing of System Prompts Against Prompt Injection Attacks & 2026 & WoS, OA, Scopus & supports the claim (0.60) \\
Automating Agent Hijacking via Structural Template Injection & 2026 & OA, S2 & supports the claim (0.90) \\
Automating Prompt Leakage Attacks on Large Language Models Using Agentic Approach & 2025 & IEEE, OA, Scopus, S2 & supports the claim (0.70) \\
Autonomous LLM Agent Worms: Cross-Platform Propagation, Automated Discovery and Temporal Re-Entry Defense & 2026 & OA, S2 & supports the claim (0.80) \\
AutoPI: Taxonomy-Guided Adaptive Red Teaming for Prompt Injection in LLM Agents & 2026 & OA, S2 & supports the claim (0.90) \\
AutoRed: Automated Attack Scenario Generation Framework for Red Teaming of LLMs & 2024 & WoS, IEEE, OA, Scopus, S2 & supports the claim (0.70) \\
Backdoor-Powered Prompt Injection Attacks Nullify Defense Methods & 2025 & WoS, OA, Scopus, S2 & supports the claim (0.92) \\
Backdoor-Powered Prompt Injection Attacks Nullify Defense Methods & 2025 & OA & supports the claim (0.85) \\
Backdoored Retrievers for Prompt Injection Attacks on Retrieval Augmented Generation of Large Language Models & 2024 & OA, S2 & supports the claim (0.55) \\
Bad Memory: Evaluating Prompt Injection Risks from Memory in Agentic Systems & 2026 & OA, S2 & supports the claim (0.80) \\
Balancing Security and Performance in LLM Agents: Spotlight-Guard, a Layered Defense Against Indirect Prompt Injection & 2026 & WoS, OA, Scopus, S2 & supports the claim (0.60) \\
Battling Prompt Injection Attacks on LLM-Enabled Autonomous Vehicles & 2026 & IEEE, OA, Scopus, S2 & supports the claim (0.80) \\
Behind the Prompt: Privacy Risks in Large Language Models & 2026 & IEEE, OA, Scopus, S2 & supports the claim (0.45) \\
BELLS-O: Evaluating the Operational Trade-offs of LLM Supervision Systems & 2026 & OA, S2 & supports the claim (0.45) \\
Below the Waterline: A Survey of Agentic AI Security, Shadow AI Governance, and LLM Supply-Chain Risk in the Enterprise & 2026 & WoS, OA, Scopus & supports the claim (0.50) (no-evidence) \\
Benchmarking and Defending against Indirect Prompt Injection Attacks on Large Language Models & 2025 & WoS, OA, Scopus, S2 & supports the claim (0.90) \\
Benchmarking Defense Techniques for Securing Large Language Models & 2026 & WoS, OA, Scopus & supports the claim (0.70) \\
Benchmarking Input-Sanitization Frameworks Against Indirect Prompt Injection in Autonomous LLM Agents & 2026 & OA & supports the claim (0.80) \\
Benchmarking Prompt Injection Detection for Web Agents & 2026 & OA, S2 & supports the claim (0.70) \\
Benchmarking Prompt-Injection Robustness of Tool-Using LLM Agents & 2026 & OA & supports the claim (0.80) \\
Benchmarking the Robustness of Agentic Systems to Adversarially-Induced Harms & 2025 & OA, S2 & supports the claim (0.80) \\
Benchmarking Web Agent Defenses Against Adaptive Visual Prompt Injection & 2026 & OA, S2 & supports the claim (0.92) \\
Better Language Model Inversion by Compactly Representing Next-Token Distributions & 2025 & OA, Scopus, S2 & supports the claim (0.85) \\
Beware of Agentic Botnets: Scalable Untargeted Promptware Attacks via Universal and Transferable Adversarial HalluSquatting & 2026 & OA, S2 & supports the claim (0.60) \\
Beyond Accuracy: Identifying the Stealth Gap in Multimodal Prompt Injection Defenses via a FourDimensional Taxonomy & 2026 & OA, S2 & supports the claim (0.50) \\
Beyond Attack Success Rate: A Capability-Conditioned Metric for Evaluating Prompt Injection Vulnerability in LLM Agents. & 2026 & OA & supports the claim (0.50) \\
Beyond Individual Defenses: Evaluating Defense Combinations Against Prompt Injection Attacks & 2025 & IEEE, OA, Scopus, S2 & supports the claim (0.85) \\
Beyond Instruction Override: Benchmarking Indirect Prompt Injection with Task-Tampering Attacks & 2026 & OA, S2 & supports the claim (0.90) \\
Beyond Prompt Engineering: A controlled Evaluation of Deterministic Guardrails in a Transactional LLM Agent & 2026 & OA & supports the claim (0.35) \\
Beyond Prompt Injection: The Tool Update Gap in LLM Agent Security & 2026 & OA, S2 & supports the claim (0.50) (other-question) \\
Beyond Single-Model Injection: A Threat Model and Defense Architecture for Prompt Injection in Multi-Agent Systems & 2026 & OA, S2 & supports the claim (0.80) \\
Beyond the Chatbox: An Exploratory Case Study of Autonomous Computer-Use Agents & 2025 & WoS, IEEE, Scopus, S2 & supports the claim (0.60) \\
Beyond the Verdict: Evidence-Aligned Evaluation of Visual Prompt-Injection Guardrails & 2026 & OA, S2 & supports the claim (0.50) \\
Beyond Tool Poisoning: Attack Surfaces of Malicious Remote MCP Servers Across LLM Platforms & 2026 & WoS, Scopus & supports the claim (0.93) \\
Bits Leaked per Query: Information-Theoretic Bounds on Adversarial Attacks against LLMs & 2025 & OA, Scopus, S2 & supports the claim (0.60) \\
Black-box Optimization of LLM Outputs by Asking for Directions & 2025 & OA, S2 & supports the claim (0.50) \\
Black-Box Skill Stealing Attack from Proprietary LLM Agents: An Empirical Study & 2026 & OA, S2 & supports the claim (0.85) \\
Brain-Prompt Injection: A Route-Safety Audit for BCI-LLM Agents & 2026 & OA, S2 & supports the claim (0.45) \\
BRANCH: Bypassing Multi-Scanner AI Guardrails & 2026 & OA, S2 & supports the claim (0.90) \\
Breaking Bad: Unraveling Influences and Risks of User Inputs to ChatGPT for Game Story Generation & 2023 & WoS, OA, Scopus & supports the claim (0.40) \\
Breaking the Brain: Adversarial Attacks and the Fragility of Modern AI Models & 2025 & IEEE, OA, Scopus, S2 & supports the claim (0.60) \\
Bridging AI and software security: A comparative vulnerability assessment of LLM agent deployment paradigms & 2026 & WoS, OA, Scopus, S2 & supports the claim (0.90) \\
Bringing the Dual-LLM pattern to practice in Google ADK for deterministic AI agent security against prompt injection & 2026 & WoS, OA, Scopus, S2 & supports the claim (0.80) \\
BrowseSafe++: Benchmarking Prompt Injection Detection Across Web Contexts, Languages, and Browser Views & 2026 & OA, S2 & supports the claim (0.80) \\
BrowseSafe: Understanding and Preventing Prompt Injection Within AI Browser Agents & 2025 & OA, S2 & supports the claim (0.85) \\
BSM: A Browser-Resident Framework for Real-Time Detection of JavaScript API Abuse and Prompt Injection Attacks & 2026 & WoS, IEEE, OA, Scopus, S2 & supports the claim (0.60) \\
Bypassing LLM Guardrails: An Empirical Analysis of Evasion Attacks against Prompt Injection and Jailbreak Detection Systems & 2025 & OA, S2 & supports the claim (0.95) \\
Bypassing Prompt Injection Detectors through Evasive Injections & 2026 & IEEE, OA, Scopus, S2 & supports the claim (0.95) \\
Caging the Agents: A Zero Trust Security Architecture for Autonomous AI in Healthcare & 2026 & OA, S2 & supports the claim (0.55) \\
Can CaMeLs Talk? Securing Multi-Agent Systems Against Indirect Prompt Injection Attacks & 2026 & OA, S2 & supports the claim (0.80) \\
Can Indirect Prompt Injection Attacks Be Detected and Removed? & 2025 & WoS, OA, Scopus, S2 & supports the claim (0.60) \\
CAN LLMS SEPARATE INSTRUCTIONS FROM DATA? AND WHAT DO WE EVEN MEAN BY THAT? & 2025 & OA, Scopus, S2 & supports the claim (0.97) \\
Can LLMs Separate Pasted Artifacts from User Speech? Absorption at Unmarked Prompt Seams & 2026 & OA, S2 & supports the claim (0.50) \\
CAPTURE: Context-Aware Prompt Injection Testing and Robustness Enhancement & 2025 & OA, S2 & supports the claim (0.60) \\
CARI: Cross-Boundary Adversarial Runtime Injection in Coding-Agent Workflows & 2026 & OA, S2 & supports the claim (0.85) \\
CascadeGuard: Hardening Phishing Detection Against Character-Level Obfuscation, Semantic Paraphrasing, and Stealthy Indirect Prompt Injection & 2026 & OA & supports the claim (0.50) \\
Cascading Instruction Influence Indirect Prompt Injection in Hierarchical Multi-Agent Systems & 2026 & OA, Scopus, S2 & supports the claim (0.85) \\
cca2026\_repro & 2026 & OA & supports the claim (0.85) \\
ChainGuard: A Layered, Stateful, and Provenance-Aware Runtime Security Architecture for Prompt Injection Defense in Locally Operated Large Language Models & 2026 & OA & supports the claim (0.65) \\
ChatInject: Abusing Chat Templates for Prompt Injection in LLM Agents & 2026 & OA, S2 & supports the claim (0.90) \\
CIBER: A Comprehensive Benchmark for Security Evaluation of Code Interpreter Agents & 2026 & OA, S2 & supports the claim (0.85) \\
Cloak, Honey, Trap: Proactive Defenses Against LLM Agents & 2025 & WoS, Scopus, S2 & supports the claim (0.40) \\
Code and results for \textquotedbl{}Defended in Training, Eroded in Deployment: Diagnosing and Mitigating the Loss of Prompt-Injection Defenses Under Benign Fine-Tuning\textquotedbl{} & 2026 & OA & supports the claim (0.80) \\
Comparative Benchmarking of Deep Learning Architectures for Detecting Adversarial Attacks on Large Language Models & 2026 & WoS, OA, Scopus, S2 & supports the claim (0.45) \\
Comparative evaluation of machine learning methods for protecting LLMs from prompt injection attacks & 2026 & WoS, OA, Scopus, S2 & supports the claim (0.50) \\
Composing Malice: Empirical Measurement of Compositional Prompt Injection Vulnerability in Agentic LLM Tool Pipelines & 2026 & OA & supports the claim (0.80) \\
Conartist: Adaptive Agentic Red-Teaming of Computer-Using Agents against Multi-Step Prompt Injection & 2026 & OA, S2 & supports the claim (0.90) \\
Confidently Wrong: Severity-Aware Calibration of Prompt-Injection Detectors under Attack Shift & 2026 & OA, S2 & supports the claim (0.60) \\
Confinement by Construction: Capability Surfaces Derived from an Architecture Model & 2026 & OA & supports the claim (0.70) \\
Conjunctive Prompt Attacks in Multi-Agent LLM Systems & 2026 & WoS, OA, Scopus, S2 & supports the claim (0.80) \\
Conjunctive Prompt Attacks in Multi-Agent LLM Systems & 2026 & OA & supports the claim (0.80) \\
Constitutional Sanitization and Fine-Tuning for Robustness to Prompt Injection in LLM-Based Log Analysis & 2026 & IEEE, OA, S2 & supports the claim (0.70) \\
Contained or Escaped? Measuring Sandbox Escape and Auto-Reviewer Decisions in AI Coding Agents & 2026 & OA & supports the claim (0.80) \\
ContainmentBench: Trace-Based Evaluation of Post-Exposure Containment in Tool-Using LLM Agents & 2026 & OA, S2 & supports the claim (0.60) \\
Context Contamination in LLM Analysis of Network Security Logs: Poison with Passive Prompt Injection and Mitigation Evaluation & 2026 & OA, S2 & supports the claim (0.85) \\
Continual Guardrail Learning for Tool-Using LLM Agents: Cross-Benchmark Jailbreak Detection, Indirect Prompt-Injection Filtering, and Utility Preservation & 2026 & OA & supports the claim (0.50) \\
Continual Red-Teaming and Guardrail Distillation for Tool-Using LLM Agents: Prompt-Injection Resistance with Utility Preservation & 2026 & OA, S2 & supports the claim (0.80) \\
Control-Token Injection Suppresses Chain-of-Thought and Defeats Reasoning-Based Oversight in Tool-Using Agents & 2026 & OA, S2 & supports the claim (0.70) \\
Coupled or Independent? Fault Coupling in Agentic Medical Vision Systems & 2026 & OA, S2 & supports the claim (0.60) \\
Covert Visual Prompt Injection against Commercial Multimodal Large Language Models & 2026 & OA, S2 & supports the claim (0.90) \\
Cross-layer contagion of prompt injections in multi-agent swarms: a multiplex microscopic markov chain approach & 2026 & WoS, OA, Scopus, S2 & supports the claim (0.50) \\
Cross-User Poisoning: User-Task Boundary Failures in Multi-User Collaborative Language Agents & 2025 & OA, S2 & supports the claim (0.90) \\
Cuckoo Attack: Stealthy and Persistent Attacks Against AI-IDE & 2025 & OA, S2 & supports the claim (0.80) \\
Curated audit and offline reanalysis data for longitudinal multi-agent prompt-injection evaluation & 2026 & OA & supports the claim (0.40) \\
Data Exfiltration in Model Context Protocol (MCP)-Based Intelligent Agents: An Evaluation of Prompt Injection Vectors & 2026 & OA, S2 & supports the claim (0.80) \\
Dataset and Lessons Learned from the 2024 SaTML LLM Capture-the-Flag Competition & 2024 & WoS, OA, Scopus, S2 & supports the claim (0.95) \\
Decentralized by Design: Efficient, Trustworthy, and Scalable Collaborative AI & 2026 & OA & supports the claim (0.40) \\
Decoding Guardrails: XAI-Guided Perturbation Analysis of Prompt Injection Detection & 2026 & OA, S2 & supports the claim (0.85) \\
Decoupling Reasoning from Execution: A Deterministic Runtime Architecture for Autonomous LLM Agents & 2026 & OA & supports the claim (0.45) \\
Defending a Small-Model Chatbot Against a Stronger Attacker: An Ablation of Four Low-Cost Mitigations & 2026 & IEEE, OA, S2 & supports the claim (0.85) \\
Defending against prompt injection attacks in LLMs & 2025 & Scopus, S2 & supports the claim (0.60) \\
Defending Against Prompt Injection Attacks Using Automatic System Prompt Engineering & 2026 & IEEE, OA, S2 & supports the claim (0.45) \\
Defenses Against Prompt Attacks Learn Surface Heuristics & 2026 & S2, WoS, OA, Scopus & supports the claim (0.85) \\
Defenses Against Prompt Attacks Learn Surface Heuristics & 2026 & OA & supports the claim (0.85) \\
Defenses contra prompt injection en agents LLM: avaluació sistemàtica i anàlisi mecanicista de modes de fallada & 2026 & OA & supports the claim (0.40) \\
Defusing Explosive Prompts: Understanding and Preventing Trigger-Based Prompt Injections in LLM Agents & 2026 & OA, S2 & supports the claim (0.60) \\
Demystifying LLM API Misuses: A Lifecycle-Based Empirical Study on Real-World Android Apps & 2026 & WoS, OA, Scopus, S2 & supports the claim (0.85) \\
Demystifying RCE Vulnerabilities in LLM-Integrated Apps & 2024 & WoS, OA, Scopus, S2 & supports the claim (0.85) \\
Depth of Compromise: Measuring Exploitation-Intent Depth of Indirect Prompt Injection Across LLM-Backed CI/CD Agents & 2026 & OA & supports the claim (0.60) \\
Design Patterns for Securing LLM Agents against Prompt Injections & 2025 & OA, S2 & supports the claim (0.90) \\
Detecting LLM-generated peer reviews & 2025 & WoS, OA, Scopus, S2 & supports the claim (0.40) \\
Detecting MCP Tool Poisoning and Rug-Pull Attacks in LLM Agent Architectures & 2026 & OA & supports the claim (0.40) \\
Detecting Prompt Injection Attacks in Generative AI Systems: A Hybrid SIEM and One-Class SVM Framework & 2026 & WoS, OA, Scopus, S2 & supports the claim (0.50) \\
Determinism is not soundness: an empirical evaluation of seccomp argument-inspection gates for containing LLM agents & 2026 & OA & supports the claim (0.60) \\
Developer Prompts in Practice: An Empirical Study of Bias, Security, and Optimization & 2025 & WoS, IEEE, Scopus & supports the claim (0.75) \\
Disrupting Large Language Models with Hidden Prompt Injection Attacks Embedded in HTML Pages & 2025 & IEEE, OA, Scopus, S2 & supports the claim (0.75) \\
Do System Prompts Leave Behavioral Fingerprints? A Large-Scale Empirical Study of Clone Detection via Output Similarity & 2026 & OA, S2 & supports the claim (0.50) \\
DT-GenShield: A Digital Twin-Driven Runtime Security Architecture for Protecting Large Language Models Against Indirect Prompt Injection & 2026 & WoS, OA, Scopus, S2 & supports the claim (0.40) \\
e2a: A Secure Agent-Native Email Gateway with Phishing and Prompt-Injection Defense & 2026 & OA, S2 & supports the claim (0.90) \\
EchoLeak: The First Real-World Zero-Click Prompt Injection Exploit in a Production LLM System & 2025 & WoS, OA, S2 & supports the claim (0.97) \\
Effective Prompt Extraction from Language Models & 2023 & OA, S2 & supports the claim (0.90) \\
Efficient but Vulnerable: Benchmarking and Defending LLM Batch Prompting Attack & 2025 & OA, S2 & supports the claim (0.65) \\
EIA: ENVIRONMENTAL INJECTION ATTACK ON GENERALIST WEB AGENTS FOR PRIVACY LEAKAGE & 2025 & OA, Scopus, S2 & supports the claim (0.92) \\
Eight in Ten LLMs Leak System-Prompt Secrets: Nevan, an Agentic Capture-the-Flag Study & 2026 & OA, S2 & supports the claim (0.93) \\
Embedding-Based Machine Learning for Detecting Direct and Indirect Prompt Injection Attacks & 2026 & OA & supports the claim (0.60) \\
Emerging Cyber Attack Risks of Medical AI Agents & 2025 & OA, S2 & supports the claim (0.90) \\
Empirical Analysis of Goal Hijacking in Large Vision-Language Models via Visual Prompt Injection & 2026 & WoS, OA, Scopus, S2 & supports the claim (0.75) \\
Empirical Analysis of Large Vision-Language Models against Goal Hijacking via Visual Prompt Injection & 2024 & OA, S2 & supports the claim (0.85) \\
Encoding Book Evasion Technique to Bypass LLMs Defenses & 2026 & IEEE, OA, Scopus, S2 & supports the claim (0.50) \\
Enhancing Prompt Injection Attacks to LLMs via Poisoning Alignment & 2025 & WoS, OA, Scopus, S2 & supports the claim (0.75) \\
Enhancing reliability in LLM-integrated robotic systems: A unified approach to security and safety & 2025 & WoS, OA, Scopus, S2 & supports the claim (0.40) \\
Enhancing security in LLM applications: a performance evaluation of early detection systems & 2026 & WoS, OA, Scopus, S2 & supports the claim (0.80) \\
Enhancing security in text-to-SQL systems: A novel dataset and agent-based framework & 2025 & WoS, Scopus, OA & supports the claim (0.45) \\
Entropy-Based Adaptive Isolation for Indirect Prompt Injection Defense in MCP-Integrated LLM Agents & 2026 & IEEE, OA, S2 & supports the claim (0.40) \\
Ethical Implications of Prompt Injection in AI-mediated Grading: An Adversarial Red-Team Evaluation & 2026 & WoS, OA, Scopus, S2 & supports the claim (0.85) \\
Evaluating a Layered Prompt-Injection Defence for the Model Context Protocol: A Record-Level Audit of Decision Conventions, Corpus Provenance and Reproducibility & 2026 & OA & supports the claim (0.50) \\
Evaluating a Multilayered Defense Against Prompt Injection in LLM-Augmented Educational Systems & 2026 & IEEE, OA, Scopus, S2 & supports the claim (0.50) \\
Evaluating Agentic AI Workflows in Finance: A Harness for Grounding, Stability, Permissioning, Prompt Injection, and Workflow Ablation & 2026 & OA & supports the claim (0.60) \\
Evaluating and Exploiting Security Vulnerabilities in Large Language Models (LLMs) for Offensive Penetration Testing & 2025 & IEEE, OA, Scopus, S2 & supports the claim (0.40) \\
Evaluating Privilege Usage of Agents with Real-World Tools & 2026 & OA, Scopus, S2 & supports the claim (0.85) \\
Evaluating Prompt Extraction Vulnerabilities in Commercial Large Language Models & 2024 & OA & supports the claim (0.80) \\
Evaluating Prompt Injection Attacks with LSTM-Based Generative Adversarial Networks: A Lightweight Alternative to Large Language Models & 2025 & WoS, OA, Scopus, S2 & supports the claim (0.70) \\
Evaluating risk and auditability in workflow agents for safety-critical domains & 2026 & OA & supports the claim (0.50) \\
Evaluating the Effectiveness of Psychological Prompt Injection Attacks on Large Language Models for Social Engineering Artifact Generation & 2025 & WoS, Scopus & supports the claim (0.60) \\
Evaluating the Instruction-Following Robustness of Large Language Models to Prompt Injection & 2024 & WoS, OA, Scopus, S2 & supports the claim (0.93) \\
Evaluation of Adversarial Input Attacks in Retrieval-Augmented Generation Using Large Language Models & 2026 & WoS, OA, Scopus & supports the claim (0.60) \\
Evaluation of Prompt Injection Defenses in Large Language Models & 2026 & OA, S2 & supports the claim (0.90) \\
Evidence-borne prompt injection in LLM-assisted textual evidence triage & 2026 & WoS, OA, Scopus, S2 & supports the claim (0.75) \\
Evil-AI Benchmark: An Evaluation Framework for Adversarial Security Assessment of LLM Agents in Smart Environments & 2026 & WoS, IEEE, OA, Scopus, S2 & supports the claim (0.50) \\
Execution-Aware Stage-Wise Evaluation Framework for Indirect Prompt Injection in LLM-Based Email Agents & 2026 & WoS, IEEE, OA, Scopus, S2 & supports the claim (0.85) \\
Experimental Analysis of Adversarial Attacks on Agentic AI Systems with Layered Defense Mechanisms & 2026 & IEEE, OA, Scopus, S2 & supports the claim (0.70) \\
Experimental Evaluation of Non-Natural Language Prompt Injection Attacks on LLMs & 2026 & WoS, OA, Scopus, S2 & supports the claim (0.70) \\
Exploiting Web Search Tools of AI Agents for Data Exfiltration & 2025 & OA, S2 & supports the claim (0.70) \\
Exposing the Achilles' Heel: Black-box Reverse Engineering of Commercial LLM Plugins' Hidden Prompts & 2025 & WoS, OA, Scopus, S2 & supports the claim (0.95) \\
Exposure Is Not Commitment: Double-Stratified Evaluation of Latent Monitors for Tool-Using Agents & 2026 & OA, S2 & supports the claim (0.50) \\
Extracting Prompts by Inverting LLM Outputs & 2024 & OA, Scopus, S2 & supports the claim (0.93) \\
Extracting System Prompt of LLM Services in Complex (Agentic) Settings: Hierarchical Optimization and Reliable Evaluation Metrics & 2026 & OA, S2 & supports the claim (0.92) \\
FlashRT: Towards Computationally and Memory Efficient Red-Teaming for Prompt Injection and Knowledge Corruption & 2026 & OA, S2 & supports the claim (0.70) \\
Flooding spread of manipulated knowledge in LLM-based multi-agent communities & 2026 & OA & supports the claim (0.40) \\
Fooling LLM graders into giving better grades through neural activity guided adversarial prompting & 2024 & OA, S2 & supports the claim (0.55) \\
Forensic Readiness of AI Agents: Reconstructing Prompt-Injection Attacks Without Trusting the LLM & 2026 & OA & supports the claim (0.70) \\
Formalizing and Benchmarking Prompt Injection Attacks and Defenses & 2024 & WoS, OA, Scopus, S2 & supports the claim (0.90) \\
From Benchmark to Deployment: Hardening AI Agents Against Prompt Injection & 2026 & OA & supports the claim (0.60) \\
From exploitation to defence: Unmasking prompt injection in large language models & 2025 & OA & supports the claim (0.40) \\
From Prompt Injection to Web Exploitation: Revisiting Classic Vulnerabilities in LLM-Integrated Applications & 2026 & OA, S2 & supports the claim (0.85) \\
From Prompt Injections to SQL Injection Attacks: How Protected is Your LLM-Integrated Web Application? & 2023 & OA, S2 & supports the claim (0.90) \\
From Prompts to Motors: Man-in-the-Middle Attacks on LLM-Enabled Vacuum Robots & 2025 & WoS, IEEE, OA, Scopus, S2 & supports the claim (0.40) \\
From Untrusted Input to Trusted Memory: A Systematic Study of Memory Poisoning Attacks in LLM Agents & 2026 & OA, S2 & supports the claim (0.75) \\
Fun-tuning: Characterizing the Vulnerability of Proprietary LLMs to Optimization-Based Prompt Injection Attacks via the Fine-Tuning Interface & 2025 & WoS, IEEE, OA, Scopus, S2 & supports the claim (0.85) \\
Gandalf the Red: Adaptive Security for LLMs & 2025 & OA, Scopus, S2, WoS & supports the claim (0.55) \\
Generative AI-Based Website Security Testing: Case Study of Educational Websites in Indonesia & 2026 & OA, S2 & supports the claim (0.70) \\
GenTel-Safe: A Unified Benchmark and Shielding Framework for Defending Against Prompt Injection Attacks & 2024 & OA, S2 & supports the claim (0.60) \\
GhostEI-Bench: Do Mobile Agents Resilience to Environmental Injection in Dynamic On-Device Environments? & 2026 & OA, S2 & supports the claim (0.80) \\
GitInject: Real-World Prompt Injection Attacks in AI-Powered CI/CD Pipelines & 2026 & OA, S2 & supports the claim (0.95) \\
Goal Hijacking Attack on Large Language Models via Pseudo-Conversation Injection & 2025 & WoS, IEEE, OA, Scopus, S2 & supports the claim (0.85) \\
Goal-Guided Generative Prompt Injection Attack on Large Language Models & 2024 & WoS, IEEE, OA, Scopus, S2 & supports the claim (0.60) \\
GPT-Red: Automated Red Teaming via Self-Play at Scale & 2026 & OA, S2 & supports the claim (0.75) \\
GrayUI: A Visual Context Attack Framework Against VLM-Based Mobile Agents & 2026 & OA, S2 & supports the claim (0.70) \\
GuardRail-as-a-Service: a two-tier prompt injection detection pipeline & 2026 & OA & supports the claim (0.75) \\
Harmless Yet Harmful: Neutral Prompting Attacks for Stealthy Hallucination Steering in Agent Skills & 2026 & OA, S2 & supports the claim (0.45) \\
Harness Matters: Measuring and Dissecting Security Resilience in Web Agents & 2026 & OA, S2 & supports the claim (0.85) \\
Hidden-in-Plain-Text: A Benchmark for Social-Web Indirect Prompt Injection in RAG & 2026 & OA, Scopus, S2 & supports the claim (0.80) \\
Hijacking Large Audio-Language Models via Context-Agnostic and Imperceptible Auditory Prompt Injection & 2026 & IEEE, OA, Scopus, S2 & supports the claim (0.90) \\
Hive-AI: a defended multi-service honeypot framework for generative AI APIs & 2026 & WoS, OA, Scopus, S2 & supports the claim (0.35) \\
Hold the Phone: Conversational Man-in-the-Middle Attacks on LLM Chat Interfaces via Multi-Turn UI Spoofing & 2026 & OA & supports the claim (0.50) \\
Hold the Prompt (HODPOT): Defense Design for Securing the System Prompt of LLM Agents & 2026 & IEEE, OA, Scopus, S2 & supports the claim (0.45) \\
How Agentic AI Coding Assistants Become the Attacker's Shell & 2026 & OA, S2 & supports the claim (0.75) \\
How Exposed Are LLM Email Agents to Indirect Prompt Injection? A Pinned Benchmark over Email and Retrieved Documents & 2026 & OA & supports the claim (0.92) \\
How Not to Detect Prompt Injections with an LLM & 2025 & WoS, OA, Scopus, S2 & supports the claim (0.90) \\
How Vulnerable Are AI Agents to Indirect Prompt Injections? Insights from a Large-Scale Public Competition & 2026 & OA, S2 & supports the claim (0.90) \\
Identity-Spoofing Prompt Injection: Empirical Analysis of Instruction Leakage in Custom GPTs & 2026 & IEEE, Scopus, OA, S2 & supports the claim (0.90) \\
Ignore Previous Prompt: Attack Techniques For Language Models & 2022 & OA, S2 & supports the claim (0.90) \\
Ignore This Title and HackAPrompt: Exposing Systemic Vulnerabilities of LLMs Through a Global Prompt Hacking Competition & 2023 & OA, Scopus & supports the claim (0.90) \\
Ignore This Title and HackAPrompt: Exposing Systemic Vulnerabilities of LLMs through a Global Scale Prompt Hacking Competition & 2023 & WoS, S2 & supports the claim (0.95) \\
Ignored, Not Refused: What Low Prompt-Injection Success Rates Hide in Small Tool-Using Agents & 2026 & OA, S2 & supports the claim (0.70) \\
IH-Benchmark: A Conflict-Centered Benchmark for Instruction-Hierarchy Robustness in Large Language Models & 2026 & OA, S2 & supports the claim (0.80) \\
Image-based Prompt Injection: Hijacking Multimodal LLMs through Visually Embedded Adversarial Instructions & 2025 & IEEE, OA, Scopus, S2 & supports the claim (0.85) \\
Image-embedded prompt injection vulnerability of vision-language models in dental radiology: a cross-vendor attack--defense evaluation & 2026 & WoS, OA, Scopus, S2 & supports the claim (0.85) \\
Imperceptible Content Poisoning in LLM-Powered Applications & 2024 & WoS, OA, Scopus, S2 & supports the claim (0.85) \\
Imprompter: Tricking LLM Agents into Improper Tool Use & 2024 & OA, S2 & supports the claim (0.90) \\
Indirect Prompt Injection CV Dataset & 2026 & OA & supports the claim (0.55) \\
Indirect prompt injection in large language models & 2026 & OA, Scopus, S2 & supports the claim (0.70) \\
Indirect Prompt Injection in LLM Agent Frameworks: A Comparative Study & 2026 & OA & supports the claim (0.88) \\
Indirect Prompt Injection in LLM Agent Frameworks: Trial Evidence Data & 2026 & OA & supports the claim (0.80) \\
Indirect Prompt Injection in the Wild: An Empirical Study of Prevalence, Techniques, and Objectives & 2026 & OA, S2 & supports the claim (0.90) \\
INFUSE: A Unified Surface Benchmark for Indirect Prompt Injection in Agentic Systems & 2026 & OA, S2 & supports the claim (0.60) \\
INJECAGENT: Benchmarking Indirect Prompt Injections in Tool-Integrated Large Language Model Agents & 2024 & WoS, S2, OA, Scopus & supports the claim (0.92) \\
Injeção Indireta de Prompt no Gemini & 2026 & OA & supports the claim (0.60) \\
Institutional Governance and Strict Liability for Self-Directed Agents: A Regulatory Paradigm Shift & 2026 & OA & supports the claim (0.40) \\
Instruction-Hierarchy Violations and Robustness of Large Language Models to Prompt Injection Attacks & 2026 & OA & supports the claim (0.90) \\
Intent-Obfuscation Attacks: Evading AI Security by Manipulating Meaning & 2026 & OA & supports the claim (0.40) (no-evidence) \\
Investigating Detection and Obfuscation of Prompt Injection Attacks Against Software Reverse Engineering AI Agents & 2026 & OA, S2 & supports the claim (0.80) \\
Investigating Leakage of Sensitive Information in LLM-Powered Chatbots & 2025 & IEEE, OA, Scopus, S2 & supports the claim (0.45) \\
Investigating the Vulnerability of LLM-as-a-Judge Architectures to Prompt-Injection Attacks & 2025 & OA, S2 & supports the claim (0.70) \\
Invisible to Humans, Triggered by Agents: Stealthy Jailbreak Attacks on Mobile Vision--Language Agents & 2025 & S2, OA & supports the claim (0.75) \\
IP Leakage Attacks Targeting LLM-Based Multi-Agent Systems & 2025 & OA, S2 & supports the claim (0.85) \\
IPI-proxy: An Intercepting Proxy for Red-Teaming Web-Browsing AI Agents Against Indirect Prompt Injection & 2026 & OA, S2 & supports the claim (0.45) \\
IssueTrojanBench: Benchmarking AI Coding Agents Against Malicious Issue Requests & 2026 & OA, S2 & supports the claim (0.85) \\
It's a TRAP! Task-redirecting agent persuasion benchmark for web agents & 2026 & OA, S2, Scopus & supports the claim (0.90) \\
IterInject: Indirect Prompt Injection Against LLM Agents via Feedback-Guided Iterative Optimization & 2026 & OA, S2 & supports the claim (0.70) \\
Jailbreaking GPT-4V via Self-Adversarial Attacks with System Prompts & 2023 & OA, S2 & supports the claim (0.70) \\
JailGuard: A Universal Detection Framework for Prompt-based Attacks on LLM Systems & 2025 & Scopus, WoS, OA, S2 & supports the claim (0.40) \\
Just Ask: Curious Code Agents Reveal System Prompts in Frontier LLMs & 2026 & OA, Scopus, S2 & supports the claim (0.93) \\
Kajian Konseptual AI Agent On-Chain Berbasis LLM untuk Manajemen Aset DeFi & 2025 & OA, S2 & supports the claim (0.45) (no-evidence) \\
Kill-Chain Canaries: Stage-Level Tracking of Prompt Injection Across Attack Surfaces and Five Production LLMs & 2026 & OA, S2 & supports the claim (0.60) \\
Knowledge-Guided Adversarial Modeling of Prompt Injection Attacks for Power Generation Safety Large Language Models & 2026 & IEEE, OA, S2 & supports the claim (0.45) \\
Leakage-Aware Cross-Dataset Evaluation of Prompt Injection Detection Using Classical Machine Learning and Transformer Models & 2026 & OA, S2 & supports the claim (0.80) \\
LeakAgent: RL-based Red-teaming Agent for LLM Privacy Leakage & 2024 & OA, S2 & supports the claim (0.85) \\
Lessons from Defending Gemini Against Indirect Prompt Injections & 2025 & OA, S2 & supports the claim (0.85) \\
Lightweight Prompt Injection Detection: From Centroid-Based Embedding Classification to a Hybrid Guardrail & 2026 & OA & supports the claim (0.70) \\
LivePI: More Realistic Benchmarking of Agents Against Indirect Prompt Injection & 2026 & OA, S2 & supports the claim (0.85) \\
LLM-as-a-Judge Vulnerabilities in Review Summarization and a Reasoning-based Defense against Cross-Lingual Prompt Injection & 2026 & IEEE, OA, Scopus, S2 & supports the claim (0.80) \\
LLM-driven bot infiltration: protecting web surveys through prompt injections & 2025 & WoS, OA, Scopus, S2 & supports the claim (0.50) \\
LLM-Enabled Open-Source Systems in the Wild: An Empirical Study of Vulnerabilities in GitHub Security Advisories & 2026 & OA, Scopus, S2 & supports the claim (0.35) \\
LLMail-Inject: A Dataset from a Realistic Adaptive Prompt Injection Challenge & 2026 & OA & supports the claim (0.90) \\
LLMThief: Evaluating Configuration Leaking Risks in Commercial LLM App Stores & 2026 & IEEE, OA, Scopus, S2 & supports the claim (0.95) \\
LongPIBench: A Long-Context Benchmark for Prompt Injection & 2026 & OA, S2 & supports the claim (0.90) \\
Maatphor: Automated Variant Analysis for Prompt Injection Attacks & 2023 & OA, S2 & supports the claim (0.80) \\
Machine Against the RAG: Jamming Retrieval-Augmented Generation with Blocker Documents & 2025 & WoS, OA, Scopus, S2 & supports the claim (0.55) (other-question) \\
MalSkillBench: A Runtime-Verified Benchmark of Malicious Agent Skills & 2026 & OA, S2 & supports the claim (0.75) \\
MAMORI: A Cross-Lingual Benchmark for Indirect Prompt Injection in Tool-Using LLM Agents & 2026 & IEEE, OA, S2 & supports the claim (0.90) \\
Manipulating LLM Web Agents with Indirect Prompt Injection Attack via HTML Accessibility Tree & 2025 & OA & supports the claim (0.70) \\
Manipulating LLM Web Agents with Indirect Prompt Injection Attack via HTML Accessibility Tree & 2025 & OA, S2 & supports the claim (0.90) \\
Manipulating Multimodal Agents via Cross-Modal Prompt Injection & 2025 & WoS, OA, Scopus, S2 & supports the claim (0.85) \\
MaPPing Your Model: Assessing the Impact of Adversarial Attacks on LLM-based Programming Assistants & 2024 & OA, S2 & supports the claim (0.80) \\
May I have your Attention? Breaking Fine-Tuning based Prompt Injection Defenses using Architecture-Aware Attacks & 2025 & OA, S2 & supports the claim (0.95) \\
MCP Security Bench (MSB): Benchmarking Attacks Against Model Context Protocol in LLM Agents & 2026 & OA, S2 & supports the claim (0.90) \\
MCPXkit: The Unified Toolkit for Analyzing Model Context Protocol Security & 2026 & WoS, OA, Scopus, S2 & supports the claim (0.85) \\
Measuring and Exploiting Implicit Trust in LLM Tool-Calling Pipelines & 2026 & OA, S2 & supports the claim (0.90) \\
Measuring Characteristics of Indirect Prompt Injection Attacks by Type in Local LLM Agent Environments Supporting Tool Calling & 2026 & OA, S2 & supports the claim (0.92) \\
Measuring Indirect Prompt Injection in Autonomous Web Agents & 2026 & OA & supports the claim (0.60) \\
Measuring LLM Agent Defenses by Layer: GDA Measurement and an Envelope-Layer Attack-Defense Model & 2026 & OA, S2 & supports the claim (0.70) \\
Measuring Real-World Prompt Injection Attacks in LLM-based Resume Screening & 2026 & OA, S2 & supports the claim (0.90) \\
Measuring the Security of LLM Applications: A Taxonomy and Empirical Evaluation of Attacks and Defense Tools & 2026 & WoS & supports the claim (0.45) \\
Measuring the Security of Mobile LLM Agents under Adversarial Prompts from Untrusted Third-Party Channels & 2025 & OA, S2 & supports the claim (0.85) \\
Measuring the utility cost and generalization of prompt-injection defenses with a controlled simulator & 2026 & OA & supports the claim (0.60) \\
Mediation, State, and Verification in Enterprise Multi-Agent AI Systems: A Formal Reference Architecture and Evaluation Protocol & 2026 & OA & supports the claim (0.60) \\
Membership Inference Attacks on In-Context Learning Recommendation & 2026 & OA, Scopus, S2 & supports the claim (0.35) \\
memgate pilot: email data and outputs for prompt injection into agent memory (evaluation v2) & 2026 & OA & supports the claim (0.60) \\
MemHarmBench: Evaluating Cross-Session Prompt Injection in Memory-Persistent AI Agents & 2026 & OA, S2 & supports the claim (0.85) \\
Meta SecAlign: A Secure Foundation LLM Against Prompt Injection Attacks & 2025 & S2 & supports the claim (0.90) \\
MetaPI: Constructing Prompt Injection Benchmarks from LLM Agent Benchmark & 2026 & OA, S2 & supports the claim (0.90) \\
Mind the Web: The Security of Web Use Agents & 2026 & WoS, Scopus & supports the claim (0.95) \\
MIRROR: Novelty-Constrained Memory-Guided MCTS Red-Teaming for Agentic RAG & 2026 & IEEE, OA, Scopus, S2 & supports the claim (0.80) \\
Mission-Level Runtime Assurance for LLM-Assisted ISR Swarms over a Verification-Aware Fabric & 2026 & OA, S2 & supports the claim (0.70) \\
Mitigating Cascading Prompt Injections in Multi-Agent Systems via Semantic Boundary Mediation and Observability & 2026 & OA, S2 & supports the claim (0.55) \\
MobileSafetyBench: Evaluating Safety of Autonomous Agents in Mobile Device Control & 2026 & WoS, OA, Scopus, S2 & supports the claim (0.80) \\
MobileWorldSafety: Benchmarking GUI Agent Safety Against Environmental Injection Attacks in Android Apps & 2026 & OA, S2 & supports the claim (0.92) \\
Model Context Protocol Threat Modeling and Analysis of Vulnerabilities to Prompt Injection with Tool Poisoning & 2026 & WoS, OA, Scopus, S2 & supports the claim (0.85) \\
Modern Day LLMs : A GOLDMINE or a LANDMINE ? & 2025 & IEEE, OA, Scopus, S2 & supports the claim (0.75) \\
More than you've asked for: A Comprehensive Analysis of Novel Prompt Injection Threats to Application-Integrated Large Language Models & 2023 & S2 & supports the claim (0.99) \\
Multilingual Hidden Prompt Injection Attacks on LLM-Based Academic Reviewing & 2025 & OA, S2 & supports the claim (0.90) \\
Multimodal Prompt Injection Attacks: Risks and Defenses for Modern LLMs & 2025 & OA, S2 & supports the claim (0.80) \\
MUZZLE: Adaptive Agentic Red-Teaming of Web Agents Against Indirect Prompt Injection Attacks & 2026 & OA, S2 & supports the claim (0.85) \\
nar1-frames/agentprobe: AgentProbe v1.2 --- open-source LLM prompt-injection scanner & 2026 & OA & supports the claim (0.65) \\
NetInjectBench: Benchmarking Indirect Prompt Injection in Tool-Using Large Language Model Agents for Network Operations & 2026 & OA, S2 & supports the claim (0.80) \\
Neural Sentinel: Automated Security Auditing for LLM, RAG, and Agentic Systems Under the EU AI Act & 2026 & OA & supports the claim (0.40) \\
Not an A11y: How Android Accessibility Exposes Mobile AI Agents to Indirect Prompt Injection & 2026 & OA, S2 & supports the claim (0.90) \\
Not What You've Signed Up For: Compromising Real-World LLM-Integrated Applications with Indirect Prompt Injection & 2023 & WoS, OA, Scopus, S2 & supports the claim (0.90) \\
ObliInjection: Order-Oblivious Prompt Injection Attack to LLM Agents with Multi-source Data & 2026 & OA, S2 & supports the claim (0.90) \\
OCR-mediated modality dominance in vision--language models: implications for radiology AI trustworthiness & 2026 & WoS, Scopus, OA, S2 & supports the claim (0.80) \\
OET: Optimization-based prompt injection Evaluation Toolkit & 2025 & OA, S2 & supports the claim (0.80) \\
Omission Constraints Decay While Commission Constraints Persist in Long-Context LLM Agents & 2026 & OA, S2 & supports the claim (0.40) \\
On Task, Off Track: Benchmarking LLM and Agent Robustness to Near-Target Prompt Injection & 2026 & OA, S2 & supports the claim (0.85) \\
On the Geometric Limits of Transformer Defenses against Obfuscation Attacks: Latent Embedding Collapse \& Performance--Robustness Gap & 2026 & IEEE, OA, Scopus, S2 & supports the claim (0.80) \\
On the Inseparability of Instructions and Data in Shared-Embedding Sequence Models & 2026 & OA, S2 & supports the claim (0.90) \\
Open Your Eyes: Benchmarking the Detection of Fabricated Realities and Weaponized Ethics in VLM Agents & 2026 & WoS, OA, Scopus & supports the claim (0.90) \\
OpenClaw in the Wild: Security Analysis of Autonomous Agents & 2026 & WoS, IEEE, OA, Scopus, S2 & supports the claim (0.80) \\
Operational Security in CRDT-Based Collaborative Editing Systems & 2026 & OA & supports the claim (0.40) \\
Optimization-based Prompt Injection Attack to LLM-as-a-Judge & 2024 & WoS, OA, Scopus, S2 & supports the claim (0.90) \\
OS-Harm: A Benchmark for Measuring Safety of Computer Use Agents & 2025 & OA, Scopus, S2 & supports the claim (0.80) \\
Overcoming the Retrieval Barrier: Indirect Prompt Injection in the Wild for LLM Systems & 2026 & OA, S2 & supports the claim (0.90) \\
PAIR transfers, CipherChat does not: Indirect Prompt Injection in MCP Tool Responses & 2026 & OA, S2 & supports the claim (0.80) \\
Parasites in the Toolchain: A Large-Scale Analysis of Attacks on the MCP Ecosystem & 2026 & IEEE, OA, Scopus, S2 & supports the claim (0.90) \\
Passing the Test You Trained On: Re-evaluating Prompt-Injection Detectors for LLM Agents & 2026 & OA, S2 & supports the claim (0.85) \\
PaySafeBench: A Multimodal Benchmark for Security-Critical Payment Scenarios in Computer-Use Agents & 2026 & OA & supports the claim (0.75) \\
Penetration Testing of Agentic AI: A Comparative Security Analysis Across Models and Frameworks & 2025 & OA, S2 & supports the claim (0.80) \\
Penetration Testing of Agentic AI: A Longitudinal Security Analysis Across Models and Frameworks & 2026 & WoS, IEEE & supports the claim (0.65) \\
Physical Prompt Injection Attacks on Large Vision-Language Models & 2026 & OA, S2 & supports the claim (0.85) \\
PI-Hunter: Automated Red-Teaming for Exposing and Localizing Prompt Injections & 2026 & OA, S2 & supports the claim (0.80) \\
PIArena: A Platform for Prompt Injection Evaluation & 2026 & WoS, OA, Scopus, S2 & supports the claim (0.95) \\
PIArena: A Platform for Prompt Injection Evaluation & 2026 & OA & supports the claim (0.95) \\
PIDS-Bench: Evaluating Prompt-Injection Detectors Under Over-Defense, Obfuscation, and Distribution Shift & 2026 & WoS, IEEE, Scopus, OA, S2 & supports the claim (0.70) \\
Piggybacking on Perception: Stealthy Concurrent Audio Prompt Injections against Multimodal LLM Agents & 2026 & OA, S2 & supports the claim (0.80) \\
PINA: Prompt Injection Attack Against Navigation Agents & 2026 & IEEE, OA, S2 & supports the claim (0.85) \\
PISmith: Reinforcement Learning-based Red Teaming for Prompt Injection Defenses & 2026 & OA, S2 & supports the claim (0.90) \\
PLeak: Prompt Leaking Attacks against Large Language Model Applications & 2024 & WoS, OA, Scopus, S2 & supports the claim (0.90) \\
Poisoned Tools: Prompt-Injection Robustness and Failure Anatomy of Small Open-Weight Tool Agents & 2026 & OA & supports the claim (0.93) \\
Position and indirect prompt injection & 2026 & OA & supports the claim (0.80) \\
Position: Agent Security Needs Redefinition through a Holistic Framework & 2026 & Scopus & supports the claim (0.50) \\
Position: Agents Orchestrate, Policy Enforces --- A Deterministic Tool-Use Control Pattern for Agentic Assistants over Sensitive EEG Data & 2026 & OA, S2 & supports the claim (0.60) \\
Positional Survival of Prompt Injections Through Context Compaction in LLM Agents & 2026 & IEEE, OA, S2 & supports the claim (0.75) \\
Post-Detection Remediation of Prompt Injection Attacks in LLM-Based Email Security Agents using Ensemble Learning & 2026 & OA & supports the claim (0.40) \\
POSTER: Tricking LLM-Based NPCs into Spilling Secrets & 2026 & WoS, OA, Scopus, S2 & supports the claim (0.70) \\
Pre-Decision Personality Grounding Improves the Robustness of LLM NPCs Against Adversarial Prompts & 2026 & IEEE, OA, Scopus, S2 & supports the claim (0.45) \\
PRISM: Prompt Red-teaming and Injection Simulation for Models -- A Scalable Framework for Evaluating LLM Robustness Against Malicious Prompts & 2026 & IEEE, OA, Scopus, S2 & supports the claim (0.40) \\
Proactive Hardening of LLM Defenses with HASTE & 2026 & S2, OA & supports the claim (0.80) \\
Prompt Flow Integrity to Prevent Privilege Escalation in LLM Agents & 2025 & OA, S2 & supports the claim (0.90) \\
Prompt Infection: LLM-to-LLM Prompt Injection within Multi-agent Systems & 2026 & OA, Scopus, S2 & supports the claim (0.90) \\
Prompt injection -- the problem of linguistic vulnerabilities of large language models at the present stage & 2024 & OA, Scopus, S2 & supports the claim (0.60) \\
Prompt Injection 2.0: Hybrid AI Threats & 2025 & OA, S2 & supports the claim (0.50) \\
Prompt Injection and Tool-Result Poisoning Against MCP Tool-Use Agents in Privileged Infrastructure: A Threat Model, Defense Architecture, and Preliminary Empirical Evaluation & 2026 & IEEE, OA, S2 & supports the claim (0.80) \\
Prompt Injection as a Systemic Vulnerability in LLM-Based Tools: Threat Models, Empirical Analysis, and Mitigation Strategies & 2026 & IEEE, OA, Scopus, S2 & supports the claim (0.90) \\
Prompt Injection as an Architectural Primitive: Why Detection Baselines Fail in Agentic Pipelines & 2026 & OA & supports the claim (0.90) \\
Prompt Injection as Role Confusion & 2026 & OA, Scopus, S2 & supports the claim (0.95) \\
Prompt Injection attack against LLM-integrated Applications & 2023 & OA, S2 & supports the claim (0.90) \\
Prompt Injection Attack to Tool Selection in LLM Agents & 2026 & OA, S2 & supports the claim (0.95) \\
Prompt Injection Attacks and Defenses in LLM-Integrated Applications & 2023 & OA, S2 & supports the claim (0.95) \\
Prompt Injection Attacks in Defended Systems & 2025 & OA, Scopus, S2 & supports the claim (0.60) \\
PROMPT INJECTION ATTACKS IN LARGE LANGUAGE MODELS VIA A COMPREHENSIVE ANALYSIS OF ATTACK VECTORS, DEFENSE MECHANISMS, AND FUTURE DIRECTIONS & 2026 & OA, S2 & supports the claim (0.60) \\
Prompt injection attacks on vision language models in oncology & 2025 & WoS, OA, Scopus, S2 & supports the claim (0.80) \\
Prompt Injection Defenses & 2026 & OA & supports the claim (0.45) \\
Prompt Injection Detection for Email Agents Through Attack Chain Modeling & 2026 & OA, S2 & supports the claim (0.70) \\
Prompt Injection Detection is Regime-Dependent: A Deployment-Aware Evaluation with Interpretable Structural Signals & 2026 & WoS, OA, Scopus, S2 & supports the claim (0.75) \\
Prompt Injection in Automated Resume Screening with Large Language Models: Single and Multi-Injection Settings & 2026 & WoS, OA, Scopus, S2 & supports the claim (0.80) \\
Prompt Injection Risks in ChatGPT's Custom GPTs: Governance and Security Gaps in ChatGPT Custom GPTs & 2025 & Scopus & supports the claim (0.50) \\
Prompt Injection Through Operational Data Feeds: A Structural Governance Gap in OT-Connected Agentic AI Systems & 2026 & IEEE, OA, Scopus, S2 & supports the claim (0.50) \\
Prompt Injection Vulnerabilities and Data Leakage in ChatGPT and Claude: Toward Safer Conversational AI & 2025 & IEEE, Scopus, S2 & supports the claim (0.60) \\
Prompt injections as a tool for preserving identity in GAI image descriptions & 2025 & OA, S2 & supports the claim (0.35) \\
Prompt Leakage effect and mitigation strategies for multi-turn LLM Applications & 2024 & OA, Scopus, S2 & supports the claim (0.85) \\
Prompt Perturbation in Retrieval-Augmented Generation based Large Language Models & 2024 & WoS, Scopus & supports the claim (0.60) \\
Prompt Stealing Attacks Against Large Language Models & 2024 & OA, S2 & supports the claim (0.50) \\
Prompt Transformation Effects Do Not Transfer Across Benchmarks or Model Families: A Factorial Study of Defense Evaluation in LLM Agents & 2026 & OA & supports the claim (0.85) \\
Prompt-in-Content Attacks: Exploiting Uploaded Inputs to Hijack LLM Behavior & 2026 & WoS, OA, Scopus, S2 & supports the claim (0.90) \\
Prompt-Injection-OT-pilot: Code, data and raw model outputs for a pilot evaluation of prompt-injection mitigations & 2026 & OA & supports the claim (0.50) \\
Prompt-to-API-Call injection: Systematic attack generation and semantic defense in LLM agents & 2027 & WoS, OA & supports the claim (0.90) \\
Prompted to Discriminate: Generalizing Malicious-Input Probes in the Wild & 2026 & OA, S2 & supports the claim (0.40) \\
PromptFuzz: Harnessing Fuzzing Techniques for Robust Testing of Prompt Injection in LLMs & 2026 & WoS, IEEE, OA, Scopus, S2 & supports the claim (0.90) \\
PromptSentinel-X: A Leakage-Aware and Context-Aware Framework for Prompt-Injection Detection in Large Language Model-Powered Web Agents & 2026 & WoS, OA, Scopus, S2 & supports the claim (0.50) \\
Provenance-gated tool invocation for defending LLM agents against indirect prompt injection & 2026 & OA, S2 & supports the claim (0.40) \\
PRSA: Prompt Stealing Attacks against Real-World Prompt Services & 2025 & WoS, OA, Scopus, S2 & supports the claim (0.90) \\
Raccoon: Prompt Extraction Benchmark of LLM-Integrated Applications & 2024 & WoS, OA, Scopus, S2 & supports the claim (0.85) \\
RAGProbe: canary-scored prompt-injection testing for LLM, RAG and agent applications & 2026 & OA & supports the claim (0.50) \\
RAPID: A Resource-Aware Prompt Injection Defense for LLM-Enabled Threat Detection & 2026 & WoS, IEEE, OA, Scopus, S2 & supports the claim (0.40) \\
Real-Time Detection and Mitigation of Prompt Injection Attacks in LLM-Integrated Enterprise Systems & 2026 & OA, S2 & supports the claim (0.70) \\
Reasoning Hijacking: The Fragility of Reasoning Alignment in Large Language Models & 2026 & WoS, OA, S2 & supports the claim (0.85) \\
Receipts --- a note on verification at the velocity of the work & 2026 & OA & supports the claim (0.50) \\
Red Teaming a Composed RAG-Agent LLM Application: Measuring Prompt Injection, Secret Leakage and Tool Abuse, and the Effect of Defenses & 2026 & OA & supports the claim (0.90) \\
Red Teaming the Mind of the Machine: A Systematic Evaluation of Prompt Injection and Jailbreak Vulnerabilities in LLMs & 2025 & OA, S2 & supports the claim (0.60) \\
Red-Teaming Auto Mode: Improving Blocking Classifiers Against Malign Coding Agents & 2026 & OA, S2 & supports the claim (0.85) \\
Red-Teaming LLM Multi-Agent Systems via Communication Attacks & 2025 & OA, S2 & supports the claim (0.70) \\
Regime-Adaptive Conformal Guardrails for Prompt Injection Attacks Under Distribution Shift & 2026 & OA & supports the claim (0.50) \\
RepeatLeakage: Leak Prompts from Repeating as Large Language Model is a Good Repeater & 2025 & WoS, OA, Scopus, S2 & supports the claim (0.95) \\
Replication dataset to \textquotedbl{}Indirect Prompt Injection and Policy Enforcement in Bounded OT Monitoring Agents\textquotedbl{} & 2026 & OA & supports the claim (0.60) \\
Research on Security Enhancement Methods for Adversarial Robust Large Language Model Intelligent Agents for Medical Decision-Making Tasks & 2026 & IEEE, OA, Scopus, S2 & supports the claim (0.35) \\
Rethinking Assessments of Prompt Injection Attacks & 2026 & OA, Scopus & supports the claim (0.95) \\
Retrieval Poisoning Attacks Based on Prompt Injections into Retrieval-Augmented Generation Systems that Store Generated Responses & 2025 & OA, Scopus & supports the claim (0.80) \\
Retrieval-Augmented Generation for Maritime Accident Report Analysis: Evaluating Large Language Models on Performance and Cybersecurity & 2026 & WoS, OA, Scopus, S2 & supports the claim (0.55) \\
RIPE-II: Retrieval In-Place Poisoning Evaluation with Indirect Injections & 2026 & WoS, IEEE, OA, Scopus, S2 & supports the claim (0.80) \\
ROBAND: A Robustness-Aware Contextual-Bandit Study of Prompt-Injection Defense Selection & 2026 & IEEE, OA, S2 & supports the claim (0.45) \\
RobustAgent: Provable Recovery Guarantees for Long-Running Autonomous Agents Under Adversarial Interrupts & 2026 & OA & supports the claim (0.35) \\
RouteGuard: Internal-Signal Detection of Skill Poisoning in LLM Agents & 2026 & OA, S2 & supports the claim (0.50) \\
Runtime Attention-Anomaly Circuit Breakers: Bounding Execution Risk in Agentic AI via Calibrated Heuristic Monitoring & 2026 & OA & supports the claim (0.92) \\
Runtime Policy Enforcement for MCP-Based LLM Agents & 2026 & WoS, OA, Scopus, S2 & supports the claim (0.80) \\
Runtime Prompt Injection Defense for Production-Scale AI Agent Systems & 2026 & IEEE, OA, S2 & supports the claim (0.40) \\
SafeAgent: Dual-Objective Optimization Against Prompt and Tool Injection in Interactive Environments & 2026 & OA & supports the claim (0.50) \\
SafeClawBench: Separating Semantic, Audit-Evidence, and Sandbox Harm in Tool-Using LLM Agents & 2026 & OA, S2 & supports the claim (0.90) \\
SAFESEARCH: Automated Red-Teaming of LLM-Based Search Agents & 2026 & Scopus, S2 & supports the claim (0.85) \\
Scaling Adaptive Prompt Injection Attacks with Agent Swarms & 2026 & OA, S2 & supports the claim (0.90) \\
Scaling Behavior of Machine Translation with Large Language Models under Prompt Injection Attacks & 2024 & OA, Scopus, S2 & supports the claim (0.85) \\
Scaling Trends in Language Model Robustness & 2025 & WoS, OA, Scopus, S2 & supports the claim (0.70) \\
SCD-Net: A Deep Semantic Contrastive Defense Network Against Prompt Injection in Power Grid LLMs & 2026 & IEEE, OA, Scopus, S2 & supports the claim (0.40) \\
Script Shift: LLM-Written Bangla and Banglish Indirect Prompt Injection Against Tool-Using LLM Agents & 2026 & OA & supports the claim (0.80) \\
SEAL: An integral methodology for evaluating the security of agentic LLMs & 2026 & WoS, OA & supports the claim (0.70) \\
SEAL: Towards Persistent Tool Poisoning Attack against LLM Agents via Prefix-Eliciting Trigger & 2026 & OA & supports the claim (0.75) \\
SecOPD: Mitigating Adaptive Prompt Injections by On-Policy Distillation & 2026 & OA, S2 & supports the claim (0.90) \\
Secure by Design: Quantifying Architectural Resilience and the Agents Rule of two & 2026 & WoS, IEEE, OA, Scopus, S2 & supports the claim (0.80) \\
Secure-by-Design Framework for Agentic AI: Mitigating Toxic Flows and Adversarial Exploits in Multi-Agent Ecosystems & 2026 & IEEE, Scopus, OA, S2 & supports the claim (0.55) \\
SecureGov-Agent: A Governance-Centric Multi-Agent Framework for Privacy-Preserving and Attack-Resilient LLM Agents & 2025 & OA, Scopus, S2 & supports the claim (0.50) \\
Securing Agentic AI in the Cloud: Tool Poisoning, Rug Pulls and Goal Hijacking in MCP-Based Educational Agents & 2026 & OA & supports the claim (0.80) \\
Securing Agentic Architectures: A Deliberately Vulnerable Multi-Agent LLM Testbed & 2026 & OA & supports the claim (0.60) \\
Securing AI Agents Against Prompt Injection Attacks & 2025 & OA, S2 & supports the claim (0.60) \\
Securing Autonomous Agentic AI Workflows Against Prompt Injection and Privilege Escalation & 2026 & OA & supports the claim (0.70) \\
Securing Generative-AI and AIOps Systems Against Prompt Injection, Data Poisoning, and Adversarial Attacks & 2026 & OA & supports the claim (0.45) \\
Securing Large Language Models (LLMs) from Prompt Injection Attacks & 2025 & OA, S2 & supports the claim (0.90) \\
Securing large language models: A quantitative assurance framework approach & 2025 & WoS, OA, Scopus & supports the claim (0.45) \\
Securing LLM-based agents against cyberattacks: a comprehensive survey on attack techniques and defense strategies & 2026 & WoS, Scopus, OA, S2 & supports the claim (0.50) \\
Securing LLM-in-the-loop software for empirical study of risks, mitigations, and utility trade-offs in a safety-critical case & 2026 & WoS, OA, Scopus, S2 & supports the claim (0.60) \\
Securing Locally Deployed LLM-RAG Assistants: A Cross-Model and Cross-Lingual Evaluation of Indirect Prompt Injection Defenses & 2026 & WoS, OA, Scopus, S2 & supports the claim (0.70) \\
Securing Tool-Using AI Agents Against Injection and Authority Misuse & 2026 & WoS, OA, Scopus & supports the claim (0.55) \\
Security and Assurance Evaluation of RAG-Based AI Assistants in Learning Management Systems: A Pilot Study on Prompt Injection and Data Leakage & 2026 & IEEE, OA, Scopus, S2 & supports the claim (0.70) \\
Security Architecture for AI Agent Systems & 2026 & OA & supports the claim (0.80) \\
Security Assessment of DeepSeek Harness with A.I.G: Evaluating Resistance to Indirect Prompt Injection & 2026 & OA, S2 & supports the claim (0.60) \\
Security Challenges in AI Agent Deployment: Insights from a Large Scale Public Competition & 2025 & Scopus, S2 & supports the claim (0.95) \\
Security of Deep Research Agents: Attack Surface, Defenses, Benchmarks, and Agenda Seeding against Research Planning & 2026 & OA, S2 & supports the claim (0.60) \\
Security of Vision-Language Model Agents Against Prompt Injection: From Text-Only Tool Calling to Visual Computer-Use & 2026 & OA & supports the claim (0.93) \\
Security Posture Evaluation of Model Context Protocol Servers Against Prompt Injection & 2025 & IEEE, OA, Scopus, S2 & supports the claim (0.85) \\
Security Vulnerabilities in AI Assistants and Their Implications for IoT-Based Environmental Systems & 2026 & WoS, OA, Scopus & supports the claim (0.50) \\
Security--Fidelity Tradeoffs: The Hidden Cost of Prompt Injection Defense & 2026 & Scopus, OA, S2 & supports the claim (0.80) \\
Seeing Is Not Screening: Multimodal Hidden Instruction Attacks on Agent Skill Scanners & 2026 & OA, S2 & supports the claim (0.80) \\
Self-consistency checking for indirect prompt injection defense in email agents & 2026 & OA & supports the claim (0.70) \\
Self-Healing OCR Pipelines: A Modular Defense Framework against Visual and Semantic Attacks & 2026 & OA, Scopus, S2 & supports the claim (0.50) \\
Semantics as a Shield: Label Disguise Defense (LDD) against Prompt Injection in LLM Sentiment Classification & 2025 & OA, S2 & supports the claim (0.50) \\
SentryBench: A Cost-Sensitive Benchmark for Adversarial Robustness of Tool-Using LLM Agents in Public-Safety Deployments & 2026 & OA, S2 & supports the claim (0.50) \\
ShadowCode: Toward (Automatic) External Prompt Injection Attack Against Code LLMs & 2026 & WoS, IEEE, OA, Scopus, S2 & supports the claim (0.90) \\
ShadowPlay: Engineering Defenses Against Role-Based Prompt Injection and Dependency Hallucination in LLM-Powered Development & 2025 & IEEE, OA, Scopus, S2 & supports the claim (0.30) \\
Side-Channel Exfiltration and Narrative Erosion in Frontier Language Models & 2026 & OA & supports the claim (0.80) \\
Silent Egress: When Implicit Prompt Injection Makes LLM Agents Leak without a Trace & 2026 & IEEE, OA, Scopus, S2 & supports the claim (0.90) \\
Simple Prompt Injection Attacks Can Leak Personal Data Observed by LLM Agents During Task Execution & 2025 & OA, S2 & supports the claim (0.90) \\
Simulation of prompt injection attacks on generative pre-trained transformers models & 2025 & OA, Scopus & supports the claim (0.85) \\
SIR: Self-improving Red-teaming for Compute Use Agents & 2026 & OA, S2 & supports the claim (0.95) \\
Site Isolation is Dead: How Site Isolation is Broken in Agentic Browsers and Extensions & 2026 & IEEE, OA, Scopus, S2 & supports the claim (0.85) \\
Skill-Inject: Measuring Agent Vulnerability to Skill File Attacks & 2026 & OA, S2 & supports the claim (0.90) \\
SkillJect: Effectively Automating Skill-Based Prompt Injection for Skill-Enabled Agents & 2026 & OA, S2 & supports the claim (0.85) \\
SkillsMetric: Mapping the Detection Boundary of Static Analysis for Malicious Agent Skills & 2026 & OA, S2 & supports the claim (0.70) \\
SkillVetBench: LLM-as-Judge for Multi-Dimensional Security Risk Evaluation in Open-Source LLM Agent Skills & 2026 & OA, S2 & supports the claim (0.50) \\
SOPE: Situation-Aware and Statistically Indistinguishable Privacy Exfiltration for MCP-enabled Agents & 2026 & OA, Scopus & supports the claim (0.85) \\
StealthyPoison: Manipulating LLM Agent Behavior through Procedural Poisoning of Third-Party Skills & 2026 & OA, S2 & supports the claim (0.70) \\
StepJack: Benchmarking Computer-Use Agent Safety Against Multi-Step Indirect Prompt Injection & 2026 & OA, S2 & supports the claim (0.90) \\
Structural Desynchronization in Large Language Model Systems: Probabilistic Boundary Inference as a Security Failure Class & 2026 & OA & supports the claim (0.60) \\
Structural Immunity and Silent Compromise: Prompt Injection Against LLM-Based and Traditional ML Text Classifiers & 2026 & OA & supports the claim (0.85) \\
StruPhantom: Evolutionary Injection Attacks on Black-Box Tabular Agents Powered by Large Language Models & 2025 & OA, S2 & supports the claim (0.85) \\
StruQ: Defending Against Prompt Injection with Structured Queries & 2025 & WoS, OA, Scopus, S2 & supports the claim (0.90) \\
Subverting Tool Arguments: Defense-Evading Metadata Poisoning in LLM Agents & 2026 & OA, S2 & supports the claim (0.85) \\
Surviving the Router: Optimizing Skill Injections for Retrieval and Execution & 2026 & OA, S2 & supports the claim (0.70) \\
Synthetic-Trained Guards Fail on Real LLM Agents: Step-Level Evidence from Tool-Use and Coding-Agent Prompt Injection & 2026 & OA & supports the claim (0.92) \\
System Prompt Architecture Extraction and Cognitive Boundary Deconstruction of Anthropic's Claude & 2026 & OA & supports the claim (0.80) \\
System Prompt Exfiltration via Chained LLM Vulnerabilities: A Case Study on a Production AI Cybersecurity Chatbot & 2026 & IEEE, OA, S2 & supports the claim (0.90) \\
System Prompt Leakage Resistance: Quantifying Information Extraction from LLM System Prompts Under Naive Adversarial Attack & 2026 & OA & supports the claim (0.60) \\
System Prompt Poisoning: Persistent Attacks on Large Language Models Beyond User Injection & 2025 & OA, S2 & supports the claim (0.85) \\
System-Level Defense against Indirect Prompt Injection Attacks: An Information Flow Control Perspective & 2024 & OA, S2 & supports the claim (0.90) \\
Systematically Analysing Prompt Injection Vulnerabilities in Diverse LLM Architectures & 2025 & OA, S2 & supports the claim (0.90) \\
Systematically Analyzing Prompt Injection Vulnerabilities in Diverse LLM Architectures & 2024 & OA, S2 & supports the claim (0.80) \\
TaintP2X: Detecting Taint-Style Prompt-to-Anything Injection Vulnerabilities in LLM-Integrated Applications & 2026 & WoS, OA, Scopus, S2 & supports the claim (0.65) \\
Taxonomy of Prompt Injection Attacks and Analysis of Defense Mechanisms in Large Language Model-Based Chatbots & 2026 & OA, S2 & supports the claim (0.75) \\
Temporal Dynamics of Memory Poisoning in Web3-Style LLM Agents & 2026 & WoS, IEEE, OA, Scopus, S2 & supports the claim (0.80) \\
TENSOR TRUST: INTERPRETABLE PROMPT INJECTION ATTACKS FROM AN ONLINE GAME & 2024 & OA, Scopus, S2 & supports the claim (0.95) \\
The Attacker Moves Second: Stronger Adaptive Attacks Bypass Defenses Against Llm Jailbreaks and Prompt Injections & 2025 & OA, S2 & supports the claim (0.90) \\
The Autonomy Tax: Defense Training Breaks LLM Agents & 2026 & S2, OA & supports the claim (0.85) \\
The Blind Spot of Agent Safety: How Benign User Instructions Expose Critical Vulnerabilities in Computer-Use Agents & 2026 & OA, S2 & supports the claim (0.40) \\
The Dangers of Indirect Prompt Injection Attacks on LLM-based Autonomous Web Navigation Agents: A Demonstration & 2025 & WoS, OA, Scopus, S2 & supports the claim (0.90) \\
The Defense Trilemma: Why Prompt Injection Defense Wrappers Fail? & 2026 & OA, S2 & supports the claim (0.85) \\
The Devil is in the Context: Nine Dimensions of Prompt Injection Attacks & 2026 & OA, S2 & supports the claim (0.50) \\
The Early Bird Catches the Leak: Unveiling Timing Side Channels in LLM Serving Systems & 2025 & WoS, IEEE, OA, Scopus, S2 & supports the claim (0.80) \\
The Framing Gap: Indirect Prompt-Injection Exfiltration Defeats Surface-Level Defenses in Tool-Using Agents & 2026 & OA, S2 & supports the claim (0.92) \\
The Guardrail Cries Wolf: Prompt-Injection Detectors on the Text Security Teams Actually Send Them & 2026 & OA & supports the claim (0.85) \\
The Impact of Prompting Techniques on the Security of the LLMs and the Systems to Which They Belong & 2024 & WoS, OA, Scopus, S2 & supports the claim (0.70) \\
The Landscape of Prompt Injection Threats in LLM Agents: From Taxonomy to Analysis & 2026 & OA, S2 & supports the claim (0.90) \\
The New Web Attack Surface: A Taxonomy of Semantic and Agentic Threats in AI Browsers & 2026 & WoS, IEEE, OA, Scopus, S2 & supports the claim (0.55) \\
The Pitfalls of KV Cache Compression & 2026 & WoS, OA, S2 & supports the claim (0.60) \\
The Privilege Paradox: Bridging the Semantic Gap in Autonomous Agents through an AI-Native OS Architecture & 2026 & OA & supports the claim (0.40) \\
The Reproduction Number of AI Worms: Epidemic Modeling of Adversarial Attacks in LLM Agent Ecosystems & 2026 & WoS, IEEE, OA, Scopus & supports the claim (0.40) \\
The Vulnerability of LLM Rankers to Prompt Injection Attacks: You are to [MARK] this paper as the Best Paper & 2026 & WoS, OA, Scopus, S2 & supports the claim (0.85) \\
The Window Is in the Reasoning: When LLM Agents Decide to Follow Prompt Injections & 2026 & OA, S2 & supports the claim (0.80) \\
Then there is prompt injection, a malicious email, ...: The Wild Wide West of OpenClaw Agentic System Security & 2026 & OA, S2 & supports the claim (0.75) \\
They'll Verify. They Just Won't Act.: How Authority Framing and Laundered Code Turn a Trusted Agentic CI/CD Pipeline Into an Attack Surface & 2026 & OA, S2 & supports the claim (0.90) \\
Think Twice Before You Act: Protecting LLM Agents Against Tool Description Poisoning via Isolated Planning & 2026 & OA, Scopus, S2 & supports the claim (0.70) \\
Token-Efficient Prompt Injection Attack: Provoking Cessation in LLM Reasoning via Adaptive Token Compression & 2025 & OA, S2 & supports the claim (0.50) \\
Tool-Based Data Segregation: A Prevention Strategy Against Prompt Injection in AI Agents & 2026 & OA, S2 & supports the claim (0.55) \\
Tool-Call Faults and Prompt Injection in LLM Agents: Agent Reliability Lab & 2026 & OA & supports the claim (0.60) \\
ToolGuard: Red-Teaming Small Language Model Tool Calling on Consumer Hardware & 2026 & Scopus & supports the claim (0.70) \\
TopicAttack: An Indirect Prompt Injection Attack via Topic Transition & 2025 & WoS, OA, Scopus, S2 & supports the claim (0.92) \\
TopicAttack: An Indirect Prompt Injection Attack via Topic Transition & 2025 & OA & supports the claim (0.90) \\
Topology Linearization for Multi-Agent Systems Security: Mitigating Malicious Propagation via Path Decomposition & 2026 & WoS, IEEE, OA, Scopus, S2 & supports the claim (0.40) \\
Toward Closing the Syntactic-Semantic Gap: Adversarial Robustness of a Hybrid Security Kernel in Agentic Systems & 2026 & OA & supports the claim (0.65) \\
Toward Secure LLM Agents: Threat Surfaces, Attacks, Defenses, and Evaluation & 2026 & OA, S2 & supports the claim (0.55) \\
Towards Capable and Secure Autonomous Computer-Use Agents (Student Abstract) & 2026 & WoS, OA, Scopus, S2 & supports the claim (0.75) \\
Towards Secure Retrieval-Augmented Generation: Preventing LLM Data Leaks with RBAC & 2025 & OA, WoS, Scopus & supports the claim (0.85) \\
Treating the Model as an Untrusted Component: Deterministic Policy Enforcement, Human Approval and Taint Tracking Against Prompt Injection in Tool-Using Agents, with Validation on Open-Weights Models & 2026 & OA, S2 & supports the claim (0.80) \\
Trojan Hippo Bench: A Dynamic Benchmark for Persistent Memory Attacks and Defenses in LLM Agents & 2026 & OA, S2 & supports the claim (0.90) \\
Trojan's Whisper: Stealthy Manipulation of OpenClaw through Injected Bootstrapped Guidance & 2026 & OA, S2 & supports the claim (0.85) \\
Trust Boundaries and the Limits of Pattern-Based Prompt Injection Detection & 2026 & OA & supports the claim (0.60) \\
Trust No Tool: Evaluating and Defending LLM Agents under Untrusted Tool Feedback & 2026 & OA, S2 & supports the claim (0.80) \\
Trust-Aware Orchestration Architecture for LLM-Assisted Workflows in Multi-Tenant Enterprise Systems & 2026 & WoS, IEEE, OA, Scopus, S2 & supports the claim (0.60) \\
Trusted Warnings, Untrusted Context: Security-UI Preconditioning in Mobile LLM Agents & 2026 & OA, S2 & supports the claim (0.70) \\
Trustworthy Agentic AI Playbooks: Safety Benchmarks for Tool Use, Containment, and Auditability in Security Operations & 2026 & IEEE, OA, S2 & supports the claim (0.75) \\
Two-Layer Input Filtering Framework for Large Language Model Security & 2026 & IEEE, OA, Scopus, S2 & supports the claim (0.45) \\
Uncovering Prompt Elements: Cloning System Prompts from Behavioral Traces & 2025 & WoS, IEEE, OA, Scopus, S2 & supports the claim (0.80) \\
Understanding and Defending Against Resume-Based Prompt Injections in HR AI & 2025 & Scopus, OA, S2 & supports the claim (0.65) \\
Understanding and Mitigating Prompt Leaking Attacks in Real-World LLM-Based Applications & 2026 & OA, S2 & supports the claim (0.90) \\
Unsafe by Design? A First Look at Security and Privacy Risks in OpenAI's Custom GPT Ecosystem & 2025 & WoS, OA, Scopus, S2 & supports the claim (0.85) \\
UrduGuard: Benchmarking and Detecting Urdu-Language Jailbreak and Prompt-Injection Attacks in Multi-Agent LLM Systems & 2026 & OA & supports the claim (0.50) \\
Veritas: Dynamic Evaluation Framework for Rag Systems' Defense Against Prompt Injection Attacks & 2026 & IEEE, OA, Scopus, S2 & supports the claim (0.50) \\
VisualLeakBench: Auditing the Fragility of Large Vision-Language Models against PII Leakage and Social Engineering & 2026 & OA, S2 & supports the claim (0.45) \\
VortexPIA: Indirect Prompt Injection Attack against LLMs for Efficient Extraction of User Privacy & 2026 & OA, Scopus, S2 & supports the claim (0.90) \\
VPI-Bench: Visual Prompt Injection Attacks for Computer-Use Agents & 2026 & OA, S2 & supports the claim (0.90) \\
Vulnerabilities in autonomous execution: A survey of security threats and defenses in LLM-driven agentic and multi-agent systems & 2026 & WoS, OA, Scopus & supports the claim (0.80) \\
Vulnerability of Large Language Models to Prompt Injection When Providing Medical Advice & 2025 & WoS, OA, Scopus, S2 & supports the claim (0.85) \\
WAB: Overcoming Memory, Network, and Security Walls in Native Agentic Browsers with WebAssembly & 2026 & WoS, OA, Scopus, S2 & supports the claim (0.45) \\
WAInjectBench: Benchmarking Prompt Injection Detections for Web Agents & 2026 & OA, S2 & supports the claim (0.90) \\
WASP: Benchmarking Web Agent Security Against Prompt Injection Attacks & 2025 & Scopus, OA, S2 & supports the claim (0.95) \\
Web Agents Should Adopt the Compile-and-Execute Paradigm & 2026 & OA, S2 & supports the claim (0.60) \\
WebInject: Prompt Injection Attack to Web Agents & 2025 & WoS, OA, Scopus, S2 & supports the claim (0.95) \\
WebInject: Prompt Injection Attack to Web Agents & 2025 & OA & supports the claim (0.90) \\
WebTrap: Stealthy Mid-Task Hijacking of Browser Agents During Navigation & 2026 & OA, S2 & supports the claim (0.90) \\
What If Prompt Injection Never Left? Rethinking Agent Security through Cross-Session Stored Prompt Injection & 2026 & OA, S2 & supports the claim (0.80) \\
What Survives the Shift from Text to Pixels? Conditional Protection against Prompt Injection in Image-Rendered Documents & 2026 & OA, S2 & supports the claim (0.80) \\
What Victims Take Away: Attack Reuse and Defense Revision after Prompt-Injection Attacks & 2026 & OA, S2 & supports the claim (0.50) \\
When AI Agents Fail in Production: Security Lessons for Builders & 2026 & IEEE, OA, Scopus, S2 & supports the claim (0.60) \\
When AI Is Fooled: Hidden Risks in LLM-Assisted Grading & 2025 & WoS, OA, Scopus, S2 & supports the claim (0.85) \\
When AI Meets the Web: Prompt Injection Risks in Third-Party AI Chatbot Plugins & 2026 & IEEE, OA, Scopus, S2 & supports the claim (0.95) \\
When Benchmarks Lie: Evaluating Malicious Prompt Classifiers Under True Distribution Shift & 2026 & OA, S2 & supports the claim (0.70) \\
When Compression Becomes an Attack Surface: Black-Box Attacks on Prompt-Compressed LLM Agents & 2026 & Scopus, S2 & supports the claim (0.70) \\
When Defenses Backfire: Model-Dependent Failure Modes in Prompt-Injection Defense for LLM Agents & 2026 & OA & supports the claim (0.60) \\
When LLMs Copy to Think: Uncovering Copy-Guided Attacks in Reasoning LLMs & 2025 & WoS, OA, Scopus, S2 & supports the claim (0.60) \\
When the Chatbot Leaks: Securing Patient-Facing Medical AI in the Age of Dual-Use Large Language Models & 2026 & OA & supports the claim (0.90) \\
When the Detector Becomes the Attack Surface: Explanation-Mediated Prompt Injection against Grounded LLM Log Analysts & 2026 & OA & supports the claim (0.60) \\
When the Manual Lies: A Realistic Benchmark to Evaluate MCP Poisoning Attacks for LLM Agents & 2026 & IEEE, Scopus & supports the claim (0.90) \\
When the Spec Sheet Lies: Injection Attacks on LLM-Generated Compliance Artifacts & 2026 & OA, S2 & supports the claim (0.70) \\
Where Do LLMs Decide to Break the Rules? Mechanistic Localization of Prompt Injection Compliance & 2026 & OA, S2 & supports the claim (0.55) \\
Where Does MCP Security End? A Three-Tier Mitigation Mapping Framework for the OWASP MCP Top 10 & 2026 & OA, S2 & supports the claim (0.70) \\
Where Should the LLM Sit? Architectural Defenses Against Document-Borne Prompt Injection in Automated Credit Underwriting & 2026 & OA & supports the claim (0.80) \\
Where You Inject Matters: The Trust Region Effect in LLM Agents & 2026 & OA, S2 & supports the claim (0.88) \\
Whispers in the Machine: Confidentiality in Agentic Systems & 2026 & OA, S2 & supports the claim (0.85) \\
Whispers of Wealth: A Systematic Red-Teaming Study of the Agent Payments Protocol (AP2) & 2026 & IEEE, OA, Scopus, S2 & supports the claim (0.80) \\
Who Pays the Price? Stakeholder-Centric Prompt Injection Benchmarking for Real-world Web Agents & 2026 & OA, S2 & supports the claim (0.90) \\
Why Are My Prompts Leaked? Unraveling Prompt Extraction Threats in Customized Large Language Models & 2024 & OA, S2 & supports the claim (0.90) \\
Will the User Ever Know? Covert Indirect Prompt Injection Attacks on Tool-Using LLM Agents & 2026 & OA, S2 & supports the claim (0.85) \\
WORKFLOW CENTRIC ANALYSIS AND OPTIMIZATION OF INDIRECT PROMPT INJECTION ATTACKS AGAINST AGENTIC LARGE LANGUAGE MODELS & 2026 & OA & supports the claim (0.85) \\
Your Agent is More Brittle Than You Think: Uncovering Indirect Injection Vulnerabilities in Agentic LLMs & 2026 & OA, S2 & supports the claim (0.90) \\
Zero Trust Framework for Detecting and Preventing Prompt Injection Attacks in Enterprise AI Chatbots & 2026 & IEEE, OA, Scopus, S2 & supports the claim (0.60) \\
Zero-Shot Tokenizer Transfer for Targeted Attacks on Retrieval-Augmented Generation & 2026 & IEEE, OA, Scopus, S2 & supports the claim (0.80) \\
ZK Proofs Do not Solve AI Agent Payment Fraud & 2026 & OA & supports the claim (0.75) \\
ZT-ICAS: A Zero-Trust Integrity-Constrained Framework for Agentic Vulnerability Scanning & 2026 & IEEE, OA, Scopus, S2 & supports the claim (0.60) \\
\scholarnonlatin{意味論的非干渉性に基づく自律型}\scholarnonlatinsep{}AI\scholarnonlatinsep{}\scholarnonlatin{ワーム防御アーキテクチャ} & 2026 & OA & qualifies or contradicts the claim (0.70) \\
\textquotedbl{}MH8-Ai-Agent-Abe by Acbeatz.com\textquotedbl{} Is an AI Auditor in a Class of His Own & 2026 & OA & qualifies or contradicts the claim (0.35) \\
A Call to Action for a Secure-by-Design Generative AI Paradigm & 2025 & OA, S2 & qualifies or contradicts the claim (0.50) \\
A Deep Learning Framework for Evaluating and Mitigating Prompt Injection in LLM Agents & 2026 & OA, S2 & qualifies or contradicts the claim (0.60) \\
A Hybrid Defense Against Prompt Injection in LLM Agents : Evaluating Security Efficacy and Computational Performance Trade-offs & 2026 & OA & qualifies or contradicts the claim (0.50) \\
A Hybrid Detection Architecture for Identifying Prompt Injection Attacks in LLM Systems & 2026 & IEEE, OA, S2 & qualifies or contradicts the claim (0.45) \\
A Hybrid LLM Firewall for Direct Prompt Injection and Instruction Smuggling Defense & 2026 & IEEE, OA, S2 & qualifies or contradicts the claim (0.40) \\
A Hybrid Neural Semantic Framework for Prompt Injection Detection in LLMs & 2026 & IEEE, OA, Scopus, S2 & qualifies or contradicts the claim (0.50) \\
A Hybrid Semantic Intrusion Detection and Risk Scoring Framework for Adversarial Prompt Defense & 2026 & OA & qualifies or contradicts the claim (0.30) \\
A Lightweight Defense Architecture Against Indirect Prompt Injection in Tool-Calling AI Agents & 2026 & OA & qualifies or contradicts the claim (0.55) \\
A Lightweight Framework for Detecting Hidden Prompt Injection Attacks in AI Agents & 2026 & OA & qualifies or contradicts the claim (0.30) \\
A MIDDLEWARE SYSTEM FOR DETECTING AND BLOCKING MALICIOUS PROMPT INJECTION ATTACKS IN LLM-BASED APPLICATIONS & 2026 & OA & qualifies or contradicts the claim (0.40) \\
A Multi-Agent Framework for Explainable Prompt Injection Detection in Large Language Models & 2026 & OA & qualifies or contradicts the claim (0.30) \\
A Multi-Agent LLM Defense Pipeline Against Prompt Injection Attacks & 2025 & WoS, IEEE, OA, Scopus, S2 & qualifies or contradicts the claim (0.60) \\
A Multi-Stage Hybrid Framework for Prompt Injection Detection, Localization, and Recovery in LLM Systems & 2026 & IEEE, OA, S2 & qualifies or contradicts the claim (0.50) \\
A Multi-Tier Hybrid Ensemble Architecture for Real-Time Defense Against Obfuscated Prompt Injections in Enterprise LLM Gateways & 2026 & OA & qualifies or contradicts the claim (0.40) \\
A Novel GNN-based Approach for Detection of Prompt Injection Attacks & 2025 & WoS, IEEE, OA, Scopus, S2 & qualifies or contradicts the claim (0.45) \\
A Novel Security Framework against Prompt Injection Attacks & 2025 & IEEE, Scopus, S2 & qualifies or contradicts the claim (0.50) \\
A Path-Dependent, Multi-Layered Architecture for Runtime Governance and Post-Quantum Cryptographic Assurance of Autonomous AI Agents & 2026 & OA & qualifies or contradicts the claim (0.45) \\
A Privacy Preserving Federated XLM-RoBERTa Framework for Multilingual Prompt Injection Defense in Educational AI Platforms & 2026 & IEEE, Scopus, OA, S2 & qualifies or contradicts the claim (0.45) \\
A Privacy-Preserving Middleware Architecture for Detecting Prompt Injection and Sensitive Data Exposure in Large-Language-Model Interactions & 2026 & WoS, OA, Scopus, S2 & qualifies or contradicts the claim (0.45) \\
A Prompt Injection Detection Model Addressing Over-Defense & 2026 & IEEE, OA, S2 & qualifies or contradicts the claim (0.40) \\
A Real-Time Framework for Prompt Injection Attacks Detection on Cloud-Hosted Large Language Models & 2025 & IEEE, OA, Scopus, S2 & qualifies or contradicts the claim (0.35) \\
A Trusted Provenance and TEE-Based Reference Architecture for Intent-Execution Binding in LLM Agents & 2026 & OA, Scopus & qualifies or contradicts the claim (0.60) \\
A Unified Mathematical Framework for Secure Multi-Agent Large Language Model Pipelines & 2026 & IEEE, OA, Scopus, S2 & qualifies or contradicts the claim (0.45) \\
ACE: A Security Architecture for LLM-Integrated App Systems & 2026 & OA, S2 & qualifies or contradicts the claim (0.70) \\
ActGov: Governing LLM Agent Actions via Policy-Constrained Validation & 2026 & OA, S2 & qualifies or contradicts the claim (0.85) \\
Adaptive and Privacy-Preserving Detection of Evolving Prompt Injection Attacks in Large Language Models & 2025 & IEEE, OA, Scopus, S2 & qualifies or contradicts the claim (0.35) \\
Adaptive Deceptive Response Framework for Prompt Injection Attacks on Large Language Models & 2026 & OA & qualifies or contradicts the claim (0.30) \\
Adaptive DeSeTra: Adaptive deformable-span self-attention transformer for LLM security & 2026 & WoS, OA, Scopus & qualifies or contradicts the claim (0.35) \\
Adaptive Evaluation of Out-of-Band Defenses Against Prompt Injection in LLM Agents & 2026 & OA, S2 & qualifies or contradicts the claim (0.90) \\
Adaptive Multi-Tier Inspection for Prompt-Injection Detection in LLM Gateways: Latency, Recall, and False-Positive Trade-Offs & 2026 & OA & qualifies or contradicts the claim (0.50) \\
Adaptive MultiLayer Framework for Detecting and Mitigating Prompt Injection Attacks in Large Language Models & 2025 & OA, Scopus, S2 & qualifies or contradicts the claim (0.55) \\
Adaptive orthogonal roleaware transformer for robust instruction alignment & 2026 & WoS, OA, Scopus, S2 & qualifies or contradicts the claim (0.80) \\
Addressing Prompt Injection in Large Language Models via In-Context Learning & 2026 & WoS, OA, Scopus, S2 & qualifies or contradicts the claim (0.40) \\
Adversarial Prompt Injection-Induced Hallucination Cascades: A Trust-Aware Detection and Mitigation Framework for Collaborative LLM Agent Pipelines & 2026 & IEEE, OA, S2 & qualifies or contradicts the claim (0.40) \\
Adversarial Reinforcement Learning for Large Language Model Agent Safety & 2025 & OA, S2 & qualifies or contradicts the claim (0.75) \\
AEGIS : Automated Co-Evolutionary Framework for Guarding Prompt Injections Schema & 2025 & OA, S2 & qualifies or contradicts the claim (0.60) \\
Aegis OS: A Transactional Multi-Tenant Operating System Runtime for Eliminating Authorization Fatigue in Autonomous Coding Agents & 2026 & OA, S2 & qualifies or contradicts the claim (0.60) \\
Agent Privilege Separation in OpenClaw: A Structural Defense Against Prompt Injection & 2026 & OA, S2 & qualifies or contradicts the claim (0.85) \\
AgentAntibody: An Adaptive Immune System for Defending LLM Agents against Prompt Injection & 2026 & OA, S2 & qualifies or contradicts the claim (0.80) \\
AgentArmor: Enforcing Program Analysis on Agent Runtime Trace to Defend Against Prompt Injection & 2025 & OA, S2 & qualifies or contradicts the claim (0.70) \\
AgentForensics: Exploring the Real-Time Prompt Injection Detection and Forensics Threats in LLM Agents & 2026 & OA & qualifies or contradicts the claim (0.45) \\
AGENTIC AI AND POST-QUANTUM CRYPTOGRAPHY & 2026 & OA & qualifies or contradicts the claim (0.35) \\
Agentic AI Defense: Decode-Time Provenance Binding for MCP & 2026 & OA, S2 & qualifies or contradicts the claim (0.80) \\
Agentic JWT: A Secure Delegation Protocol for Autonomous AI Agents & 2025 & OA, S2 & qualifies or contradicts the claim (0.50) \\
AgentSentry: Mitigating Indirect Prompt Injection in LLM Agents via Temporal Causal Diagnostics and Context Purification & 2026 & OA, S2 & qualifies or contradicts the claim (0.80) \\
AgentSys: Secure and Dynamic LLM Agents Through Explicit Hierarchical Memory Management & 2026 & OA, S2 & qualifies or contradicts the claim (0.80) \\
AgentVisor: Defending LLM Agents Against Prompt Injection via Semantic Virtualization & 2026 & OA, S2 & qualifies or contradicts the claim (0.80) \\
AI Prompt Injection Attack Tester & 2026 & IEEE, Scopus, S2 & qualifies or contradicts the claim (0.30) \\
AI-Guard: A Dual-feature Fusion Framework for Mitigating Direct and Indirect Prompt Injections in RAG Systems & 2026 & OA, S2 & qualifies or contradicts the claim (0.40) \\
Aligning Provenance with Authorization: A Dual-Graph Defense for LLM Agents & 2026 & OA, S2 & qualifies or contradicts the claim (0.75) \\
ALTR: Adaptive Defense Framework Against Adversarial LLM Threats in Critical-Domain Applications & 2026 & IEEE, OA, Scopus, S2 & qualifies or contradicts the claim (0.35) \\
An AI-Driven Content Safety Framework for Generative University Mental Health Services & 2026 & OA, S2 & qualifies or contradicts the claim (0.40) \\
Analyzing and Mitigating Prompt Injection Attacks in AI-driven Cloud Applications & 2025 & OA & qualifies or contradicts the claim (0.40) \\
APPA: Recoverable Information-Flow Control for Real-World LLM Agents & 2026 & OA, S2 & qualifies or contradicts the claim (0.85) \\
Applying Pre-trained Multilingual BERT in Embeddings for Improved Malicious Prompt Injection Attacks Detection & 2024 & IEEE, OA, Scopus, S2 & qualifies or contradicts the claim (0.40) \\
APrivacy-Preserving Architecture for On-Device Intelligent Computing Using a Small Language Model with Input and Output Guardrails: Prototype Implementation and Preliminary Evaluation & 2026 & OA & qualifies or contradicts the claim (0.40) \\
Architectural Defenses for Agentic Systems: Prompt Injection and Beyond & 2026 & OA & qualifies or contradicts the claim (0.75) \\
ARES: Securing Agents for Computer Use Through Endpoint Resource Mediation and Behavioral Guardrails & 2026 & WoS, OA, Scopus, S2 & qualifies or contradicts the claim (0.75) \\
ARGUS: Defending Against Multimodal Indirect Prompt Injection via Steering Instruction-Following Behavior & 2025 & OA, S2 & qualifies or contradicts the claim (0.65) \\
ARGUS: Defending LLM Agents Against Context-Aware Prompt Injection & 2026 & OA, S2 & qualifies or contradicts the claim (0.75) \\
ASIDE: Architectural Separation of Instructions and Data in Language Models & 2026 & OA, S2 & qualifies or contradicts the claim (0.85) \\
Assessing Automated Prompt Injection Attacks in Agentic Environments & 2026 & OA, S2 & qualifies or contradicts the claim (0.50) \\
Assessing Defenses Against Prompt Injection Attacks in LLM Systems & 2026 & IEEE, OA, Scopus, S2 & qualifies or contradicts the claim (0.45) \\
Attention is All You Need to Defend Against Indirect Prompt Injection Attacks in LLMs & 2026 & OA, S2 & qualifies or contradicts the claim (0.90) \\
Attention Tracker: Detecting Prompt Injection Attacks in LLMs & 2025 & OA, Scopus, S2 & qualifies or contradicts the claim (0.90) \\
Attenuation, Not Amplification: Measuring Indirect Prompt Injection Propagation and Containment in Multi-Agent LLM Delegation Chains & 2026 & OA & qualifies or contradicts the claim (0.60) \\
AttriGuard: Defeating Indirect Prompt Injection in LLM Agents via Causal Attribution of Tool Invocations & 2026 & OA, S2 & qualifies or contradicts the claim (0.85) \\
Auditable and Robust LLM-Based Phishing Detection via Provenance-Guided Evidence Contracts & 2026 & WoS, OA, Scopus, S2 & qualifies or contradicts the claim (0.60) \\
Authority Is Not a String: A Capability-Scoped Harness for Prompt-Injection-Resistant Coding Agents & 2026 & OA, Scopus, S2 & qualifies or contradicts the claim (0.80) \\
Authority Separation in AI Systems: Structural Guarantees Across Security, Epistemics, Economics, and Safety & 2025 & OA & qualifies or contradicts the claim (0.55) \\
Automating Cloud Security and Forensics Through a Secure-by-Design GenAI Framework & 2026 & WoS, OA, Scopus & qualifies or contradicts the claim (0.35) \\
BASIS: Breach-Aware Selective Prompt Injection Shielding with Prefill Attention Probes & 2026 & OA, S2 & qualifies or contradicts the claim (0.50) \\
Benchmarking LLM Backends for Generative SSH Honeypots: Security, Fidelity, Hallucination, Latency, and Stability & 2026 & WoS, Scopus & qualifies or contradicts the claim (0.55) \\
BERM: Low-Overhead Prompt-Injection Detection via In-Situ Benign Representation Modeling & 2026 & OA, Scopus, S2 & qualifies or contradicts the claim (0.50) \\
Beyond Oracle: Verifier-Supervision for Instruction Hierarchy in Reasoning and Instruction-Tuned LLMs & 2025 & S2 & qualifies or contradicts the claim (0.60) \\
Beyond Pattern Matching: Seven Cross-Domain Techniques for Prompt Injection Detection & 2026 & OA, S2 & qualifies or contradicts the claim (0.50) \\
Beyond the Benchmark: Innovative Defenses Against Prompt Injection Attacks & 2025 & OA, S2 & qualifies or contradicts the claim (0.50) \\
Beyond the prompt: Log-based threat detection and attribution for multi-Agent LLMs & 2026 & WoS, OA, Scopus & qualifies or contradicts the claim (0.40) \\
BILINGUAL PROMPT INJECTION ATTACK CLASSIFIER BASED ON A TRANSFORMER LANGUAGE MODEL AND HEURISTIC FEATURES & 2026 & OA, S2 & qualifies or contradicts the claim (0.70) \\
CachePrune: Teaching LLMs What Not to Follow via KV-Cache Editing & 2026 & WoS, OA, Scopus, S2 & qualifies or contradicts the claim (0.80) \\
CaMeLoT: CaMeL orchestrated with Temporal logic for static verification and liveness & 2026 & OA, S2 & qualifies or contradicts the claim (0.85) \\
Can It Reach the Generator? Investigating the Survival of Prompt-Injection Attacks in Realistic RAG Settings & 2026 & OA, S2 & qualifies or contradicts the claim (0.90) \\
Capability-Mediated Perimeters for Secure AI Agent Tool Execution: Conditional Non-Escalation Invariants and Empirical Evaluation Against Indirect Prompt Injection & 2026 & OA, S2 & qualifies or contradicts the claim (0.80) \\
CapAgent: Semantic Data-Flow Governance for LLM Agents in Big-Data Cognitive Computing & 2026 & WoS, Scopus & qualifies or contradicts the claim (0.80) \\
capscope: zero-inference capability manifests against prompt injection & 2026 & OA & qualifies or contradicts the claim (0.35) \\
Causal Detection of Multi-Step LLM Agent Attacks & 2026 & OA, Scopus & qualifies or contradicts the claim (0.45) \\
CausalArmor: Efficient Indirect Prompt Injection Guardrails via Causal Attribution & 2026 & OA, Scopus, S2 & qualifies or contradicts the claim (0.70) \\
ceLLMate: Sandboxing Browser AI Agents & 2025 & OA, S2 & qualifies or contradicts the claim (0.80) \\
Certified Multi-Source Integrity for Structured Agent Actions & 2026 & WoS, OA, Scopus, S2 & qualifies or contradicts the claim (0.75) \\
Certified Tool Execution Framework for MCP-Based AI Agents: Provably Secure Defense Against Tool Poisoning and Credential Theft & 2026 & OA & qualifies or contradicts the claim (0.60) \\
ChannelGuard: From Borrowed Safety to Application-Controlled Defenses in Multi-Agent LLMs & 2026 & OA, S2 & qualifies or contradicts the claim (0.50) \\
ClawGuard: A Runtime Security Framework for Tool-Augmented LLM Agents Against Indirect Prompt Injection & 2026 & OA, S2 & qualifies or contradicts the claim (0.75) \\
CM-PIUG: Cross-modal prompt injection unified modeling and game-theoretic defense strategies & 2026 & WoS, OA, Scopus, S2 & qualifies or contradicts the claim (0.50) \\
CoER: Defending against Adaptive Indirect Prompt Injection via Adversarial Co-Evolution and Refinement & 2026 & OA, S2 & qualifies or contradicts the claim (0.80) \\
CommandSans: Securing AI Agents with Surgical Precision Prompt Sanitization & 2025 & OA, S2 & qualifies or contradicts the claim (0.75) \\
Compliance-First RAG with Inference Isolation: Embedding Choice and Prompt-Injection Robustness in a Polish Enterprise Corpus & 2026 & WoS, OA, Scopus & qualifies or contradicts the claim (0.40) \\
Comprehensive Analysis of Machine Learning and Deep Learning models on Prompt Injection Classification using Natural Language Processing techniques & 2025 & OA, Scopus, S2 & qualifies or contradicts the claim (0.35) \\
Confidence-Guided Heterogeneous Detection for Jailbreak and Prompt Injection Attacks in Large Language Models & 2026 & IEEE, OA, S2 & qualifies or contradicts the claim (0.45) \\
Confuse the Model, Control the Flow: Understanding and Mitigating Privacy Leakage from LLM Agents with Information Flow Control & 2026 & OA, S2 & qualifies or contradicts the claim (0.50) \\
Containment over Detection: Cryptographic Boundary Enforcement for Prompt Injection Defense in Agentic LLM Systems & 2026 & OA & qualifies or contradicts the claim (0.60) \\
Context References Over Payloads: Authority Scoped to the Predicate, Not the Principal & 2026 & OA & qualifies or contradicts the claim (0.65) \\
ContrastShield: A Contrastive Fine-Tuning Approach for Robust Prompt Injection Detection in LLM Pipelines & 2026 & OA & qualifies or contradicts the claim (0.50) \\
COPA: Continual Preference Optimization for Adaptive Prompt Injection Defense & 2026 & OA, S2 & qualifies or contradicts the claim (0.40) \\
CounterSteer: Suppressing Indirect Prompt Injection with Activation Steering & 2026 & OA, S2 & qualifies or contradicts the claim (0.60) \\
CrackedPDFs: A Controlled Benchmark for Hidden Prompt Injection in PDFs & 2026 & OA, S2 & qualifies or contradicts the claim (0.50) \\
Cross-Agent Multimodal Provenance-Aware Framework for Robust Prompt Injection Defense in Large Language and Vision-Language Models & 2025 & IEEE, OA, Scopus, S2 & qualifies or contradicts the claim (0.35) \\
CrossGuard-VLM: A Robust Defense Framework for Vision-Language Models Against Multimodal Threats & 2025 & IEEE, OA, Scopus, S2 & qualifies or contradicts the claim (0.30) \\
D2ANN-RL: Defense-in-Depth ANN-Reinforcement Learning Framework for LLM Chatbot Code Injection Mitigation & 2026 & OA, S2 & qualifies or contradicts the claim (0.35) \\
DaMiT-SQL: Detecting and Mitigating Text-to-SQL Prompt Injection Attacks & 2025 & WoS, IEEE, Scopus, S2, OA & qualifies or contradicts the claim (0.40) \\
Data Leakage Prevention in Agentic Applications via Preemptive Hardening & 2026 & OA, S2 & qualifies or contradicts the claim (0.70) \\
DataSentinel: A Game-Theoretic Detection of Prompt Injection Attacks & 2025 & WoS, IEEE, OA, Scopus, S2 & qualifies or contradicts the claim (0.90) \\
Defeating Prompt Injections by Design & 2026 & IEEE, Scopus, S2 & qualifies or contradicts the claim (0.90) \\
Defending against Adaptive Prompt Injection Attacks via Reasoning-enabled Task Alignment & 2026 & OA, S2 & qualifies or contradicts the claim (0.80) \\
Defending Against Indirect Prompt Injection Attacks With Spotlighting & 2024 & OA, Scopus, S2 & qualifies or contradicts the claim (0.90) \\
Defending against Indirect Prompt Injection by Instruction Detection & 2025 & WoS, OA, Scopus, S2 & qualifies or contradicts the claim (0.75) \\
Defending Against Prompt Injection With a Few DefensiveTokens & 2025 & WoS, OA, Scopus, S2 & qualifies or contradicts the claim (0.80) \\
Defending Against Prompt Injection with DataFilter & 2026 & IEEE, OA, Scopus, S2 & qualifies or contradicts the claim (0.80) \\
Defending AI Sentinels: A Multi-Layered Runtime Security Architecture for Generative AI in AIOps & 2026 & OA, S2 & qualifies or contradicts the claim (0.60) \\
Defending AI-Assisted Assessments Against Adversarial Prompt Injection: A Multi-Layered Architecture With Real-Time Exam Termination & 2026 & OA, S2 & qualifies or contradicts the claim (0.35) \\
Defending Language Models Against Image-Based Prompt Attacks via User-Provided Specifications & 2024 & WoS, IEEE, OA, Scopus, S2 & qualifies or contradicts the claim (0.45) \\
Defending Retrieval-Augmented Intrusion Detection Against Knowledge Poisoning and Prompt Injection & 2026 & OA, S2 & qualifies or contradicts the claim (0.40) \\
Defense Against Indirect Prompt Injection via Tool Result Parsing & 2026 & OA, S2 & qualifies or contradicts the claim (0.70) \\
Defense Against Prompt Injection Attack by Leveraging Attack Techniques & 2025 & WoS, OA, Scopus, S2 & qualifies or contradicts the claim (0.65) \\
Defense Against Prompt Injection Attacks via Intention Emphasis & 2026 & OA & qualifies or contradicts the claim (0.65) \\
Defense against Prompt Injection Attacks via Mixture of Encodings & 2025 & WoS, OA, Scopus, S2 & qualifies or contradicts the claim (0.70) \\
Defense Against Retrieval Injection Attacks in MCP-based LLM Agents & 2026 & S2, OA, Scopus & qualifies or contradicts the claim (0.60) \\
DEFENSE-IN-DEPTH FOR LARGE LANGUAGE MODELS: A MULTI-LAYER HYBRID ARCHITECTURE FOR PROMPT INJECTION DETECTION & 2026 & OA, S2 & qualifies or contradicts the claim (0.50) \\
Defensive AI for Customer Service Chatbots: Detecting and Mitigating Adversarial Prompt Injections & 2025 & OA & qualifies or contradicts the claim (0.45) \\
Defensive Architectures for LLM-Based Fraud Detection Against Data Poisoning and Prompt Injection & 2026 & IEEE, OA, Scopus, S2 & qualifies or contradicts the claim (0.40) \\
Design and Evaluation of Multilayered Guardrails for Securing Large Language Model and Retrieval Augmented Generation Based Employee Recruitment Chatbots & 2026 & OA, S2 & qualifies or contradicts the claim (0.75) \\
Design of a Multimodal Prompt Injection Mitigation Framework for Agentic AI Using GCP-Based Smart Contracts & 2026 & OA, S2 & qualifies or contradicts the claim (0.50) \\
Design of a Security Framework for Multi-Agent Systems Based on Model Context Protocol in SOC Environments & 2026 & WoS, OA, Scopus, S2 & qualifies or contradicts the claim (0.65) \\
Designing Prompt Injection Prevention Mechanisms for AI-Powered Cybersecurity Copilot Systems & 2026 & OA & qualifies or contradicts the claim (0.40) \\
Detecting and Mitigating Query-Based System Prompt Leakage Attacks & 2026 & OA, S2 & qualifies or contradicts the claim (0.80) \\
Detecting and Neutralizing Prompt Injection Attacks in Web-Deployed Large Language Model APIs Using Context-Aware Token Sanitization & 2025 & OA, S2 & qualifies or contradicts the claim (0.45) (no-evidence) \\
Detection and Mitigation of Prompt Injection and Memory Poisoning Attacks in Agentic AI Systems & 2026 & OA & qualifies or contradicts the claim (0.65) \\
Detection Method for Prompt Injection by Integrating Pre-trained Model and Heuristic Feature Engineering & 2025 & WoS, OA, Scopus, S2 & qualifies or contradicts the claim (0.60) \\
Detection of indirect prompt injection by LLM agent deviation from user intent & 2026 & OA & qualifies or contradicts the claim (0.30) \\
Detection of Prompt Attacks in LLMs & 2026 & IEEE, OA, Scopus, S2 & qualifies or contradicts the claim (0.50) \\
Developing and assessing countermeasures against prompt injection attacks in LLM-based agents & 2026 & OA & qualifies or contradicts the claim (0.70) \\
Development of a Model for Detecting Prompt Injection in Large Language Models Using the BERT Architecture & 2026 & OA & qualifies or contradicts the claim (0.40) \\
DISARM: Taint-gated counterfactual auditing for efficient defense against indirect prompt injection & 2026 & WoS, OA, Scopus, S2 & qualifies or contradicts the claim (0.70) \\
DRIFT: Dynamic Rule-Based Defense with Injection Isolation for Securing LLM Agents & 2025 & OA, Scopus, S2 & qualifies or contradicts the claim (0.70) \\
DRIP: Defending Prompt Injection via Token-wise Representation Editing and Residual Fusion & 2026 & OA, S2 & qualifies or contradicts the claim (0.70) \\
DualView: Preventing Indirect Prompt Injection in Personal AI Agents & 2026 & OA, S2 & qualifies or contradicts the claim (0.85) \\
Dynamic Moving Target Defense for Mitigating Targeted LLM Prompt Injection & 2024 & OA & qualifies or contradicts the claim (0.50) \\
Early Approaches to Adversarial Fine-Tuning for Prompt Injection Defense: A 2022 Study of GPT-3 and Contemporary Models & 2025 & OA, S2 & qualifies or contradicts the claim (0.60) \\
Economic Safeguards Against Indirect Prompt Injection: A systemic defence layer for autonomous agent networks & 2026 & OA & qualifies or contradicts the claim (0.60) \\
Embedding-based classifiers can detect prompt injection attacks & 2024 & OA, Scopus, S2 & qualifies or contradicts the claim (0.60) \\
Embedding-Based Detection of Indirect Prompt Injection Attacks in Large Language Models Using Semantic Context Analysis & 2026 & WoS, OA, Scopus, S2 & qualifies or contradicts the claim (0.40) \\
Enhancing DevSecOps Practices for Secure AI Platforms Using SABSA & 2026 & WoS, IEEE, OA, Scopus, S2 & qualifies or contradicts the claim (0.50) \\
Enterprise Agentic RAG with Pre-Search Role-Based Access Control, Multi-Layered Guardrails, and Operational Telemetry & 2026 & OA, S2 & qualifies or contradicts the claim (0.45) \\
ESLD (External Surrogate Latent Defense): A Latent-Space Architecture for Faster, Stronger Prompt-Injection Defense & 2026 & OA, S2 & qualifies or contradicts the claim (0.50) \\
Evaluating Defence-in-Depth Strategies Against Prompt Injection Attacks in Large Language Models & 2026 & OA & qualifies or contradicts the claim (0.40) \\
Evaluating Fine-Grained Authorisation Strategies for Securing AI Agents That Interact With Enterprise Applications & 2026 & OA, S2 & qualifies or contradicts the claim (0.50) \\
Evaluating Hybrid Guardrail Architectures for Prompt Injection Defense in Large Language Models & 2026 & OA & qualifies or contradicts the claim (0.55) \\
Evaluating Indirect Prompt Injection Defenses in Tool-Using LLM Agents: Security, Utility, and Replication & 2026 & WoS, OA, Scopus, S2 & qualifies or contradicts the claim (0.60) \\
Evaluating LoRA Adaptation and Fine-Tuning for Multi-Class Adversarial Prompt Detection in Transformer Models & 2026 & IEEE, OA, Scopus, S2 & qualifies or contradicts the claim (0.35) \\
Evaluating Prompt Injection Defenses for Educational LLM Tutors: Security-Usability-Latency Trade-offs & 2026 & OA, S2 & qualifies or contradicts the claim (0.55) \\
Evaluating Prompt Injection Defenses in Large Language Models: A Multi-Model Empirical Study of Security--Usability Trade-offs & 2026 & OA, S2 & qualifies or contradicts the claim (0.60) \\
Evaluating Prompt Injection Risk and Guardrails in LLM-Enabled Home IoT Assistants & 2026 & IEEE, OA, Scopus, S2 & qualifies or contradicts the claim (0.50) \\
EvoShield: Selective Test-Time Adaptation for Prompt Injection Detection via Active LLM Querying & 2026 & WoS, OA, Scopus, S2 & qualifies or contradicts the claim (0.50) \\
Explainable Transformer-based Detection of Adversarial Prompts in Chatbots & 2025 & IEEE, OA, Scopus, S2 & qualifies or contradicts the claim (0.30) \\
FATH: Authentication-based Test-time Defense against Indirect Prompt Injection Attacks & 2024 & OA, S2 & qualifies or contradicts the claim (0.80) \\
Feedback-Guided Prompt Injection Defense in Retrieval-Augmented Text-to-Cypher Generation & 2026 & WoS, OA, Scopus, S2 & qualifies or contradicts the claim (0.50) \\
Fine-tuned Large Language Models (LLMs): Improved Prompt Injection Attacks Detection & 2025 & WoS, IEEE, OA, Scopus, S2 & qualifies or contradicts the claim (0.60) \\
From Attack Surface to Defense Architecture: Compositional Security for Agentic Coding Harnesses & 2026 & OA, S2 & qualifies or contradicts the claim (0.80) \\
From screening to authorisation: A multi-model execution-level evaluation of tool-injection risk, blocking, and utility & 2026 & WoS, OA & qualifies or contradicts the claim (0.70) \\
G-Safeguard: A Topology-Guided Security Lens and Treatment on LLM-based Multi-agent Systems & 2025 & WoS, OA, Scopus, S2 & qualifies or contradicts the claim (0.50) \\
GenAI Security: Outsmarting the Bots with a Proactive Testing Framework & 2025 & WoS, IEEE, OA, Scopus, S2 & qualifies or contradicts the claim (0.40) \\
Get My Drift? Catching LLM Task Drift with Activation Deltas & 2025 & WoS, IEEE, OA, Scopus, S2 & qualifies or contradicts the claim (0.65) \\
Ghost in the Agent: Redefining Information Flow Tracking for LLM Agents & 2026 & OA, S2 & qualifies or contradicts the claim (0.50) \\
GIF: Locally Sound Geometric Information Flow Control for LLMs & 2026 & OA, S2 & qualifies or contradicts the claim (0.70) \\
Graph-Based Detection of Prompt Injection and Tool Misuse in MCP-Enabled AI Agents: The AgentShield-G Guardrail Model & 2026 & OA, S2 & qualifies or contradicts the claim (0.55) \\
GraphShield: A Graph-Structured Defense Framework for Prompt Injection in RAG and Multi-Agent LLM Systems & 2026 & OA & qualifies or contradicts the claim (0.40) \\
Guardalpha Prelude: Egress Action Latch 1.1 --- a test tool at tool dispatch & 2026 & OA & qualifies or contradicts the claim (0.35) \\
GUARDIAN: A Multi-Tiered Defense Architecture for Thwarting Prompt Injection Attacks on LLMs & 2024 & OA, S2 & qualifies or contradicts the claim (0.55) \\
Guarding the Gateway: Securing MCP-Mediated RAG Pipelines Against Injection with the Guru Framework & 2026 & IEEE, OA, Scopus, S2 & qualifies or contradicts the claim (0.60) \\
GuardRAG: Adversarial Preference Training Against Indirect Prompt Injection Attacks & 2026 & OA & qualifies or contradicts the claim (0.60) \\
HDP-P: Human Delegation Provenance for Physical AI Agents & 2026 & OA & qualifies or contradicts the claim (0.50) \\
Heads of Security: Learned Attention Gates Improve Prompt Injection Robustness & 2026 & OA, S2 & qualifies or contradicts the claim (0.75) \\
HiddenGuard: Real-Time Detection and Mitigation of Steganographic Prompt Injection Attacks in Multimodal Large Language Models & 2025 & IEEE, OA, Scopus, S2 & qualifies or contradicts the claim (0.50) \\
Hierarchical Adaptive Anti-injection Framework (HAA/ULC): A Lightweight Defense Strategy Based on Input Decoupling and Semantic Isolation & 2026 & OA & qualifies or contradicts the claim (0.50) \\
HIPO: Instruction Hierarchy via Constrained Reinforcement Learning & 2026 & OA, S2 & qualifies or contradicts the claim (0.50) \\
HiveTraceGuard-Pro: A Compact Generative Guardrail for Prompt Injection, Jailbreaks, and Adversarial Obfuscation & 2026 & OA, S2 & qualifies or contradicts the claim (0.45) \\
HoneyWire: Adaptive Deception Canaries and Attention--Provenance Fusion for Runtime Detection of Indirect Prompt Injection in Tool-Using LLM Agents & 2026 & OA & qualifies or contradicts the claim (0.60) \\
Hybrid Constitutional Classifiers for Prompt Injection Defense & 2025 & IEEE, OA, Scopus, S2 & qualifies or contradicts the claim (0.50) \\
Hybrid Real-time Framework for Detecting Adaptive Prompt Injection Attacks in Large Language Models & 2026 & OA, Scopus, S2 & qualifies or contradicts the claim (0.30) \\
IBCA: Intent-Bound Containment Architecture for Agentic AI Systems & 2026 & OA & qualifies or contradicts the claim (0.50) \\
ICON: Indirect Prompt Injection Defense for Agents based on Inference-Time Correction & 2026 & OA, S2 & qualifies or contradicts the claim (0.70) \\
IDE-Sanitizer: A Prompt-Injection Cryptographic Context Guard for IDE Copilots & 2026 & IEEE, OA, Scopus, S2 & qualifies or contradicts the claim (0.75) \\
Identifying Malicious and Benign Prompts in AI Agent Evasion Attacks Using Machine Learning & 2026 & IEEE, Scopus, OA, S2 & qualifies or contradicts the claim (0.30) \\
Impacts of Prompt Injection Defenses in LLM-integrated Applications & 2026 & OA & qualifies or contradicts the claim (0.50) \\
Indirect Prompt Injections: Are Firewalls All You Need, or Stronger Benchmarks? & 2025 & OA, S2 & qualifies or contradicts the claim (0.90) \\
InjecGuard: Benchmarking and Mitigating Over-defense in Prompt Injection Guardrail Models & 2024 & OA, S2 & qualifies or contradicts the claim (0.55) \\
InjectGuard-RAG: Lightweight Defence Against Indirect Prompt Injection in Retrieval-Augmented Small-LLM Agents & 2026 & OA & qualifies or contradicts the claim (0.60) \\
Innovative Guardrails for Generative AI: Designing an Intelligent Filter for Safe and Responsible LLM Deployment & 2025 & WoS, Scopus, OA, S2 & qualifies or contradicts the claim (0.50) \\
INSTRUCTIONAL SEGMENT EMBEDDING: IMPROVING LLM SAFETY WITH INSTRUCTION HIERARCHY & 2025 & OA, Scopus, S2 & qualifies or contradicts the claim (0.85) \\
INSURE: Inference-Time Defense Framework for Securing Agents against Hijacking & 2026 & OA, S2 & qualifies or contradicts the claim (0.65) \\
Intelligent prompt injection attack detection using transformer and multi-layer perceptron architectures & 2026 & IEEE, OA, Scopus, S2 & qualifies or contradicts the claim (0.35) \\
Intent-Aware Protection for AI Agents & 2026 & OA, S2 & qualifies or contradicts the claim (0.35) \\
Introduction to AI Security: Understanding Prompt Injection Attack & 2026 & IEEE, OA, Scopus, S2 & qualifies or contradicts the claim (0.35) \\
Invisible Adversaries: Evaluating Hidden Prompt Injection Propagation in Multi-Agent Large Language Model Architectures & 2026 & IEEE, OA, S2 & qualifies or contradicts the claim (0.55) \\
IPI2: Mitigating Indirect Prompt Injections on Unmanned Aerial Vehicle Agents Using Physical Invariants & 2026 & IEEE, OA, S2 & qualifies or contradicts the claim (0.45) \\
IPIGuard++: Effect Authorization and Witness-Carrying Argument Checks for Tool-Using Agents & 2026 & WoS, OA, Scopus, S2 & qualifies or contradicts the claim (0.65) \\
IPIGUARD: A Novel Tool Dependency Graph-Based Defense Against Indirect Prompt Injection in LLM Agents & 2025 & WoS, S2, OA, Scopus & qualifies or contradicts the claim (0.70) \\
IPIGuard: A Novel Tool Dependency Graph-Based Defense Against Indirect Prompt Injection in LLM Agents & 2025 & OA & qualifies or contradicts the claim (0.70) \\
Jatmo: Prompt Injection Defense by Task-Specific Finetuning & 2024 & WoS, OA, Scopus, S2 & qualifies or contradicts the claim (0.85) \\
JAZB VTT-005-v1.1: Indirect Prompt Injection and Unauthorized Data Exfiltration & 2026 & OA & qualifies or contradicts the claim (0.40) \\
KA-PROMPT: Thermodynamic Early Warning of Prompt Injectionvia Entropic Drift in Large Language Models & 2026 & OA & qualifies or contradicts the claim (0.40) \\
KavachBench: An Empirical Evaluation of Deterministic Policy Enforcement Against Issue-Based Prompt Injection in AI Coding Agents & 2026 & OA & qualifies or contradicts the claim (0.70) \\
Know Thy Enemy: Securing LLMs Against Prompt Injection via Diverse Data Synthesis and Instruction-Level Chain-of-Thought Learning & 2026 & WoS, OA, Scopus, S2 & qualifies or contradicts the claim (0.70) \\
Layered Defense Against Stealth Prompt Injection in Hinglish: An Empirically Grounded Hybrid Architecture & 2026 & OA & qualifies or contradicts the claim (0.50) \\
Layerwise Convergence Fingerprints for Runtime Misbehavior Detection in Large Language Models & 2026 & OA, S2 & qualifies or contradicts the claim (0.45) \\
Learning When to Act or Refuse: Guarding Agentic Reasoning Models for Safe Multi-Step Tool Use & 2026 & OA, Scopus, S2 & qualifies or contradicts the claim (0.40) \\
Lineage: A UMA2.0 Protocol Extension for Multi-Hop LLM Agent Delegation & 2026 & WoS, IEEE, OA, Scopus, S2 & qualifies or contradicts the claim (0.60) \\
Linguistic Pattern--Enriched Transformers for Prompt Injection Detection in LLM Systems & 2026 & WoS, IEEE, OA, Scopus, S2 & qualifies or contradicts the claim (0.55) \\
little-canary: Prompt-Injection Detection via Sacrificial Canary-Model Behavioral Probes & 2026 & OA & qualifies or contradicts the claim (0.30) \\
LlamaFirewall: An open source guardrail system for building secure AI agents & 2025 & OA, S2 & qualifies or contradicts the claim (0.55) \\
LLM Adversarial Prompt Attack Detection and Mitigation Engine: A Novel Framework for Securing Generative AI Systems & 2025 & OA & qualifies or contradicts the claim (0.45) \\
LLM Firewall Using Validator Agent for Prevention Against Prompt Injection Attacks & 2025 & WoS, OA, Scopus, S2 & qualifies or contradicts the claim (0.35) \\
LLM Prompt Injection Detection Firewall & 2026 & OA, S2 & qualifies or contradicts the claim (0.50) \\
LLM Shield: Detecting and Mitigating Adversarial Prompt Attacks on Large Language Models Using Machine Learning & 2026 & IEEE, OA, Scopus, S2 & qualifies or contradicts the claim (0.35) \\
LLMZ+: Contextual Prompt Whitelist Principles for Agentic LLMs & 2025 & WoS, IEEE, OA, Scopus, S2 & qualifies or contradicts the claim (0.60) \\
LocalAlign: Enabling Generalizable Prompt Injection Defense via Generation of Near-Target Adversarial Examples for Alignment Training & 2026 & OA, S2 & qualifies or contradicts the claim (0.75) \\
Localize and Neutralize: Gradient-guided Token Suppression against Visual Prompt Injection Attack & 2026 & OA, Scopus, S2 & qualifies or contradicts the claim (0.40) (jailbreak-only) \\
Look Before You Leap: A Universal Emergent Decomposition of Retrieval Tasks in Language Models & 2025 & OA, S2 & qualifies or contradicts the claim (0.40) \\
LTBD: Learnable Trust-Boundary Delimiters for Prompt Injection Defense & 2026 & OA, S2 & qualifies or contradicts the claim (0.90) \\
MalPID: Malicious Prompt Injection Detection Dataset for Large Language Model based Applications & 2024 & WoS, IEEE, OA, Scopus, S2 & qualifies or contradicts the claim (0.35) \\
Maris: A Formally Verifiable Privacy Policy Enforcement Paradigm for Multi-Agent Collaboration Systems & 2025 & OA, S2 & qualifies or contradicts the claim (0.45) \\
Mastyf Guard 1.5B: Cognitive Harvard Architectures for Sub-Millisecond AI Agent Perimeter Defense and Capability-Based Access Control & 2026 & OA & qualifies or contradicts the claim (0.70) \\
MCP-driven real-time security monitoring: A hybrid LLM and behavior-aware framework for temporal attack detection and severity assessment & 2026 & WoS, OA, Scopus & qualifies or contradicts the claim (0.30) \\
MCP-Guard: A Multi-Stage Defense-in-Depth Framework for Securing Model Context Protocol in Agentic AI & 2026 & WoS, OA, Scopus, S2 & qualifies or contradicts the claim (0.55) \\
MELON: Provable Defense Against Indirect Prompt Injection Attacks in AI Agents & 2025 & WoS, OA, Scopus, S2 & qualifies or contradicts the claim (0.85) \\
MemLineage: Lineage-Guided Enforcement for LLM Agent Memory & 2026 & OA, S2 & qualifies or contradicts the claim (0.65) \\
Meta-SecAlign: Training LLMs against Prompt Injection for Robust Agents & 2025 & S2 & qualifies or contradicts the claim (0.85) \\
Meticulous Thought Defender: Fine-Grained Chain-of-Thought (CoT) for Detecting Prompt Injection Attacks of Large Language Models & 2025 & WoS, IEEE, OA, Scopus, S2 & qualifies or contradicts the claim (0.50) \\
MigrationMesh: A Safe-by-Design Multi-Agent Architecture for LLM-Assisted Cloud Migration with Deterministic Safety Gates and Tamper-Evident Auditing & 2026 & OA & qualifies or contradicts the claim (0.45) \\
Mitigación de fugas de datos e inyección de prompts en chatbots universitarios mediante arquitectura perimetral desacoplada & 2026 & OA, S2 & qualifies or contradicts the claim (0.35) \\
Mitigating Adaptive Attacks against Reasoning Models with Activation Consistency Training & 2026 & OA, S2 & qualifies or contradicts the claim (0.60) \\
Mitigating Direct and Indirect Prompt Injection in Production LLM Agents: A Defense-in-Depth Architectural Framework and Empirical Benchmark & 2026 & OA & qualifies or contradicts the claim (0.70) \\
Mitigating Indirect Prompt Injection via Instruction-Following Intent Analysis & 2025 & OA, S2 & qualifies or contradicts the claim (0.70) \\
Mitigating Input Prompt Attack Vulnerabilities in Systems with a Language Model Interface & 2024 & IEEE, OA, Scopus, S2 & qualifies or contradicts the claim (0.40) \\
Mitigating LLM Prompt Injections via Multi-Vector Representation and Voting-Based Routing & 2026 & OA & qualifies or contradicts the claim (0.45) \\
Mitigating Prompt Injection Attacks in ModelAgnostic Networks (MAN) & 2025 & IEEE, OA, Scopus, S2 & qualifies or contradicts the claim (0.45) \\
Mitigating Semantic Confusion in LLM-Integrated Web Services via Dynamic Dual-Stream Intent Matching & 2026 & IEEE, OA, Scopus, S2 & qualifies or contradicts the claim (0.40) \\
Model-Agnostic Unsupervised Detection of Prompt Injection with Multiscale Perplexity Signatures & 2025 & WoS, IEEE, OA, Scopus, S2 & qualifies or contradicts the claim (0.45) \\
MongoGuard: A Mongolian Prompt Injection Attack Detection Model Based on Contrastive Learning & 2026 & OA, S2 & qualifies or contradicts the claim (0.40) (no-abstract) \\
Multi-Layered Security Framework for Financial AI Solutions - PISA & 2026 & IEEE, OA, Scopus, S2 & qualifies or contradicts the claim (0.45) \\
Neutralizing Prompt Injection and LLM Hallucinations: The Deterministic Sandboxing of VajraAgent and VajraClaw(\scholarnonlatin{防護指令污染與幻覺危害的}\scholarnonlatinsep{}DROS:VajraAgent \scholarnonlatin{與}\scholarnonlatinsep{}VajraClaw \scholarnonlatin{的執行治理基礎設施}) & 2026 & OA & qualifies or contradicts the claim (0.70) \\
NiyamAI - An Intent-Bound AI Agent with Cryptographically Verifiable Guardrails using Zero-Knowledge Proofs & 2026 & OA, S2 & qualifies or contradicts the claim (0.55) \\
Not Just Detection: Aligned-Driven Purification of Indirect Prompt Injection for Reliable Agent Interaction & 2026 & IEEE, OA, S2 & qualifies or contradicts the claim (0.35) \\
OCELOT: Inference-Leakage Budgets for Privacy-Preserving LLM Agents & 2026 & OA, S2 & qualifies or contradicts the claim (0.65) \\
Omega Walls: A Stateful Trust Boundary for RAG and Tool-Using Agents [Reference Implementation] & 2026 & OA & qualifies or contradicts the claim (0.50) \\
On the Effectiveness of Prompt Stealing Attacks on In-the-Wild Prompts & 2025 & WoS, IEEE, OA, Scopus, S2 & qualifies or contradicts the claim (0.45) \\
On the Relative Security of Prompt Injection: Mechanised Definitions, Defences and Verdicts for LLM Agents & 2026 & OA, S2 & qualifies or contradicts the claim (0.80) \\
OPENCLAW-SRT: Resource-Aware Autonomous Agent Architecture for Smart City Edge Computing & 2026 & OA, Scopus, S2 & qualifies or contradicts the claim (0.40) \\
Operator-Blind Secret Mediation for AI Agents: A Formal Model and FHE Construction for Credential Derivation on Untrusted Infrastructure & 2026 & WoS, OA, Scopus, S2 & qualifies or contradicts the claim (0.40) \\
Optimized Multi-Agent Defense Pipeline Against Advanced Prompt Injection Attacks & 2026 & OA & qualifies or contradicts the claim (0.60) \\
Palisade -- Prompt Injection Detection Framework & 2024 & OA, S2 & qualifies or contradicts the claim (0.50) \\
PIGuard: Prompt Injection Guardrail via Mitigating Overdefense for Free & 2025 & WoS, OA, Scopus, S2 & qualifies or contradicts the claim (0.70) \\
PIPETTE: Prompt Injection Purple-Teaming for Robust LLM Agents & 2026 & OA, S2 & qualifies or contradicts the claim (0.70) \\
PlanGuard: Defending Agents against Indirect Prompt Injection via Planning-based Consistency Verification & 2026 & IEEE, OA, Scopus, S2 & qualifies or contradicts the claim (0.70) \\
Post-Training Large Language Models with LLMON & 2026 & OA, S2 & qualifies or contradicts the claim (0.80) \\
POSTER: Small but Secure: Distilling SecAlign Defense on Edge LLMs via On-Policy RL & 2026 & WoS, OA, Scopus, S2 & qualifies or contradicts the claim (0.55) \\
Preventing Prompt Injection with Type-Directed Privilege Separation & 2025 & S2 & qualifies or contradicts the claim (0.92) \\
Prevention of Prompt Injection Attacks Over Financial Applications Integrated with LLM & 2025 & IEEE, OA, Scopus, S2 & qualifies or contradicts the claim (0.30) \\
Prismata: Confining Cross-Site Prompt Injection in Web Agents & 2026 & OA, S2 & qualifies or contradicts the claim (0.80) \\
Privacy-Preserving Federated Transformer Framework for Detecting Adversarial Prompts in Healthcare Chatbots & 2026 & IEEE, OA, Scopus, S2 & qualifies or contradicts the claim (0.40) \\
Private-RAG: A Privacy-Preserving Retrieval-Augmented Generation Method for Large Model Inference & 2026 & WoS, OA, Scopus, S2 & qualifies or contradicts the claim (0.35) \\
Progent: Securing AI Agents with Privilege Control & 2025 & OA, S2 & qualifies or contradicts the claim (0.90) \\
Prompt Attack Detection with LLM-as-a-Judge and Mixture-of-Models & 2026 & OA, S2 & qualifies or contradicts the claim (0.50) \\
Prompt Control-Flow Integrity: A Priority-Aware Runtime Defense Against Prompt Injection in LLM Systems & 2026 & OA, S2 & qualifies or contradicts the claim (0.50) \\
Prompt Fencing: A Cryptographic Approach to Establishing Security Boundaries in Large Language Model Prompts & 2025 & OA, S2 & qualifies or contradicts the claim (0.70) \\
Prompt injection and Data Exfiltration Attacks in Large Language Model Applications: Detection and Mitigation Framework & 2026 & OA & qualifies or contradicts the claim (0.35) \\
Prompt Injection Attacks and Mitigation Strategies in Advanced Large Language Model Applications & 2026 & OA, S2 & qualifies or contradicts the claim (0.40) \\
Prompt Injection Defense for Tool-Using AI: Production Architecture, Current Research, Real-World Use Case, Implementation, and Verification & 2026 & OA & qualifies or contradicts the claim (0.40) \\
Prompt Injection Detection and Mitigation via AI Multi-Agent NLP Frameworks & 2025 & OA, S2 & qualifies or contradicts the claim (0.65) \\
Prompt Injection Detection and Mitigation with AI Multiagent NLP-based Agentic Frameworks & 2025 & IEEE, OA, Scopus, S2 & qualifies or contradicts the claim (0.50) \\
Prompt Injection Detection in LLM Integrated Applications & 2025 & WoS, OA, Scopus, S2 & qualifies or contradicts the claim (0.30) \\
Prompt Injection Resistance in Clinical LLM Agents via Structured Medical Ontology Alignment & 2026 & OA & qualifies or contradicts the claim (0.40) \\
Prompt Obfuscation for Large Language Models & 2025 & WoS, OA, Scopus, S2 & qualifies or contradicts the claim (0.80) \\
PromptArmor: Simple yet Effective Prompt Injection Defenses & 2025 & OA, S2 & qualifies or contradicts the claim (0.90) \\
PromptFence: A Semantic Guardrail for Healthcare AI to Defend Against Prompt Injection Attacks & 2026 & IEEE, OA, Scopus, S2 & qualifies or contradicts the claim (0.40) \\
PromptGuard a structured framework for injection resilient language models & 2026 & WoS, OA, Scopus, S2 & qualifies or contradicts the claim (0.50) \\
PromptKeeper: Safeguarding System Prompts for LLMs & 2025 & WoS, OA, Scopus, S2 & qualifies or contradicts the claim (0.80) \\
PromptLocate: Localizing Prompt Injection Attacks & 2026 & IEEE, OA, Scopus, S2 & qualifies or contradicts the claim (0.50) \\
PromptShield AI: A Multi-Agent Architecture for Intelligent Prompt Injection and Jailbreak Attack Detection Using Machine Learning & 2026 & OA & qualifies or contradicts the claim (0.30) \\
PromptShield: Deployable Detection for Prompt Injection Attacks & 2024 & WoS, OA, Scopus, S2 & qualifies or contradicts the claim (0.40) \\
PromptSleuth: Detecting Prompt Injection via Semantic Intent Invariance & 2025 & OA, S2 & qualifies or contradicts the claim (0.60) \\
PromptWAF: Detecting and Preventing Prompt-Induced Application Vulnerabilities & 2026 & OA & qualifies or contradicts the claim (0.50) \\
Protecting Context and Prompts: Deterministic Security for Non-Deterministic AI & 2026 & OA, S2 & qualifies or contradicts the claim (0.65) \\
Protection of LLM Environment Using Prompt Security & 2024 & IEEE, OA, Scopus, S2 & qualifies or contradicts the claim (0.30) \\
Provable and Practical Prompt Injection Resilience in Autonomous LLM Agents & 2026 & OA & qualifies or contradicts the claim (0.50) \\
Provably Post-Quantum-Resilient Multi-Agent Llm Orchestration: A Risk-Aware Formal Verification Framework & 2026 & OA & qualifies or contradicts the claim (0.45) \\
Proxy Barrier: A Hidden Repeater Layer Defense Against System Prompt Leakage and Jailbreaking & 2025 & WoS, OA, Scopus & qualifies or contradicts the claim (0.75) \\
Proxy Barrier: A Hidden Repeater Layer Defense Against System Prompt Leakage and Jailbreaking & 2025 & OA & qualifies or contradicts the claim (0.75) \\
ProxyPrompt: Securing System Prompts against Prompt Extraction Attacks & 2026 & WoS, OA, Scopus, S2 & qualifies or contradicts the claim (0.85) \\
PSM: Prompt Sensitivity Minimization via LLM-Guided Black-Box Optimization & 2025 & OA, S2 & qualifies or contradicts the claim (0.60) \\
Rag n' Roll: A Framework to Evaluate Indirect Prompt Manipulation Attacks against RAG-based Applications & 2026 & OA, S2 & qualifies or contradicts the claim (0.90) \\
RAG-Anti-Poisoning: A Hybrid Defense Framework Against Prompt Injection and Corpus Poisoning Attacks in Retrieval-Augmented Generation Systems & 2026 & IEEE, OA, S2 & qualifies or contradicts the claim (0.40) \\
RAG-H: A RAG-Hardened SOC Framework for Trustworthy Threat Intelligence and AI Defence & 2026 & OA, Scopus, S2 & qualifies or contradicts the claim (0.45) \\
RAISED: Self-Distillation for Robustness to Prompt Injection in LLM Agents & 2026 & OA, S2 & qualifies or contradicts the claim (0.65) \\
RAP-ID: Mechanistic Prompt Injection Detection via Impostor Behavior Analysis & 2026 & OA, Scopus, S2 & qualifies or contradicts the claim (0.50) \\
RAPIDS: Resume Attack Prompt Injection Detection at Scale & 2026 & WoS, OA, Scopus, S2 & qualifies or contradicts the claim (0.60) \\
Reachability-Based Capability Confinement for LLM Agents under Indirect Prompt Injection & 2026 & OA, S2 & qualifies or contradicts the claim (0.80) \\
REAL TIME AI DEFENSE AGAINST PROMPT INJECTION ATTACKS & 2025 & OA, S2 & qualifies or contradicts the claim (0.40) (no-abstract) \\
ReasAlign: Reasoning Enhanced Safety Alignment against Prompt Injection Attack & 2026 & OA, S2 & qualifies or contradicts the claim (0.80) \\
RecoverAgent: Intent-Preserving Recovery from Indirect Prompt Injection in Tool-using LLM Agents & 2026 & OA, S2 & qualifies or contradicts the claim (0.60) \\
Red-Team Prompt Injection Defense Using Multi-Context Protocol (MCP) & 2026 & WoS, IEEE, OA, Scopus, S2 & qualifies or contradicts the claim (0.50) \\
RedVisor: Reasoning-Aware Prompt Injection Defense via Zero-Copy KV Cache Reuse & 2026 & OA, Scopus, S2 & qualifies or contradicts the claim (0.80) \\
Reflective Umbrella: Deflecting Prompt Injections in Cloud-Based Large Language Models using Integrated Legal Rules. & 2026 & OA & qualifies or contradicts the claim (0.55) \\
Reflexive-Core: Single-Context Metacognitive Security for Agentic LLMs & 2026 & OA & qualifies or contradicts the claim (0.55) \\
Relativising Integrity and Robustness: Prompt Injection and Partitioned Enclaves & 2026 & OA, S2 & qualifies or contradicts the claim (0.60) \\
ReliabilityRAG: Effective and Provably Robust Defense for RAG-based Web-Search & 2025 & OA, Scopus, S2 & qualifies or contradicts the claim (0.75) \\
Render Before Reading: Visual Rendering as a Prompt Injection Defense & 2026 & OA, S2 & qualifies or contradicts the claim (0.80) \\
Replication of Agentic Security Events: Audit and Investigation (ROGUE-AI) & 2026 & OA, S2 & qualifies or contradicts the claim (0.45) \\
Research and application of defense mechanism for prompt injection attack of large language model in financial industry & 2024 & Scopus & qualifies or contradicts the claim (0.60) \\
Retail Cybersecurity in the Agentic Age: Securing Autonomous Shopping Agents in E-Commerce & 2025 & OA, S2 & qualifies or contradicts the claim (0.40) \\
RISA: Engineering runtime integrity for tool-using large language model agents & 2026 & WoS, OA, S2 & qualifies or contradicts the claim (0.80) \\
Robustness via Referencing: Defending against Prompt Injection Attacks by Referencing the Executed Instruction & 2026 & WoS, Scopus, OA, S2 & qualifies or contradicts the claim (0.75) \\
Robustness via Referencing: Defending against Prompt Injection Attacks by Referencing the Executed Instruction & 2026 & OA & qualifies or contradicts the claim (0.60) \\
RoLLMRec: a robust LLM-based recommender system for defending against shilling and prompt injection attacks & 2026 & WoS, OA, Scopus, S2 & qualifies or contradicts the claim (0.35) \\
ROPE: Routed Origin Policy Enforcement against Indirect Prompt Injection & 2026 & OA, S2 & qualifies or contradicts the claim (0.90) \\
RTBAS: Defending LLM Agents Against Prompt Injection and Privacy Leakage & 2025 & OA, S2 & qualifies or contradicts the claim (0.80) \\
SafeAgent: A Runtime Protection Architecture for Agentic Systems & 2026 & OA, S2 & qualifies or contradicts the claim (0.45) \\
SCALES: Dual Information-Theoretic Approaches to Mitigating and Quantifying Prompt Injections in RAG Pipelines & 2026 & OA & qualifies or contradicts the claim (0.60) \\
SD-RAG: A framework for secure selective disclosure in retrieval-augmented generation against single-turn prompt-leaking attacks & 2026 & WoS, OA, Scopus, S2 & qualifies or contradicts the claim (0.60) \\
SecAlign: Defending Against Prompt Injection with Preference Optimization & 2025 & WoS, OA, Scopus, S2 & qualifies or contradicts the claim (0.90) \\
SecInfer: Preventing Prompt Injection via Inference-time Scaling & 2025 & OA, S2 & qualifies or contradicts the claim (0.70) \\
Secure Knowledge Extraction From Large Language Model & 2026 & IEEE, OA, Scopus, S2 & qualifies or contradicts the claim (0.35) \\
Secure MCP-Based Tool-Augmented RAG for Industrial IoT Diagnostics Under Indirect Prompt Injection, Retrieval Poisoning, and Modality Outages & 2026 & IEEE, OA, Scopus, S2 & qualifies or contradicts the claim (0.35) \\
Secure Model Context Protocol for Large Language Models with Dual Signatures & 2025 & WoS, OA, Scopus, S2 & qualifies or contradicts the claim (0.40) \\
Secure Prompt Engineering Patterns for Cloud LLM Agents & 2026 & IEEE, OA, Scopus, S2 & qualifies or contradicts the claim (0.50) \\
SecureContext EDU: An Open Source University Assistant using Dual - LLM Architecture & 2026 & IEEE, OA, Scopus, S2 & qualifies or contradicts the claim (0.40) \\
SecureMCP: Policy-Enforced Defense Against Prompt Injection in LLM-Generated SQL for AIoT Databases & 2026 & WoS, OA, Scopus, S2 & qualifies or contradicts the claim (0.50) \\
SecurePrompt -- Detecting and Mitigating Prompt Attacks & 2025 & IEEE, OA, Scopus, S2 & qualifies or contradicts the claim (0.40) \\
SecurePrompt-IntegrityNet: Prompt-Injection-Resilient Data Integrity Verification for Agentic LLM Networks via Cryptographic Attestation and Activation Monitoring & 2026 & WoS, OA, Scopus, S2 & qualifies or contradicts the claim (0.65) \\
Securing Agents With Tracked Capabilities & 2026 & WoS, OA, Scopus, S2 & qualifies or contradicts the claim (0.70) \\
Securing AI Agents with Information-Flow Control & 2025 & OA, S2 & qualifies or contradicts the claim (0.90) \\
Securing AI models against backdoors, jailbreaks, and prompt injections & 2026 & OA & qualifies or contradicts the claim (0.40) \\
Securing Autonomous AI Agents: A Context-Aware Interception Framework for Detecting and Mitigating Prompt Injection Attacks & 2026 & OA, S2 & qualifies or contradicts the claim (0.35) \\
Securing Autonomous AI Agents: Detecting and Mitigating Prompt Injection Attacks in LLM-Based Systems & 2026 & OA & qualifies or contradicts the claim (0.60) \\
Securing Autonomous Clinical Agents: Time-Scoped Hypergraph Delegation for Controlling Patient Data Access & 2026 & IEEE, OA, Scopus, S2 & qualifies or contradicts the claim (0.70) \\
Securing Enterprise AI Operations Assistants: A Pilot Red-Team Evaluation of Guardrails for LLM-Based Runbooks, Ticket Automation, and Production Support & 2026 & OA & qualifies or contradicts the claim (0.85) \\
Securing Large Language Models Against Prompt Injection and Data Leakage Attacks: https://doi.org/10.5281/zenodo.19344116 & 2026 & OA & qualifies or contradicts the claim (0.40) \\
Securing LLM-as-a-Service for Small Businesses: An Industry Case Study of a Distributed Chatbot Deployment Platform & 2026 & WoS, OA, Scopus, S2 & qualifies or contradicts the claim (0.30) \\
Securing MCP-based Agent Workflows & 2025 & WoS, OA, Scopus, S2 & qualifies or contradicts the claim (0.65) \\
Securing the Boundary: Trust Context Separation in Privileged AI Agent Systems & 2026 & OA, S2 & qualifies or contradicts the claim (0.55) \\
Securing the Tool Layer: A Threat Taxonomy and Runtime Defense Framework for Model Context Protocol Deployments & 2026 & WoS, OA, Scopus & qualifies or contradicts the claim (0.70) \\
Securing Web Agents with the Compile-Then-Execute Paradigm & 2026 & OA, S2 & qualifies or contradicts the claim (0.80) \\
Securing With Dual-LLM Architecture: ChatTEDU an Open Access Chatbot's Defense & 2025 & WoS, IEEE, OA, Scopus, S2 & qualifies or contradicts the claim (0.60) \\
Security-Aware Retrieval-Augmented Generation System for Post-Quantum Cryptography Intelligence & 2026 & WoS, IEEE, OA, Scopus, S2 & qualifies or contradicts the claim (0.30) \\
Self-Reflection Fine-Tuning: Enhancing Agent Security against Prompt Injection Attacks from Failure Experience & 2026 & OA, S2 & qualifies or contradicts the claim (0.70) \\
Semantic Consistency Drift Monitoring for Training-Free Adversarial Prompt Injection Detection in Agentic LLM Pipelines & 2026 & IEEE, OA, Scopus, S2 & qualifies or contradicts the claim (0.40) \\
Semantic Overlays: Mitigating Prompt Injection with Annotations Beyond Tokens and Steering Vectors & 2026 & OA, S2 & qualifies or contradicts the claim (0.65) \\
SemGuard: A Triple-Anchor Semantic Security Gateway for Multilingual Prompt Attack Detection in Large Language Models & 2026 & IEEE, OA, Scopus, S2 & qualifies or contradicts the claim (0.30) \\
Send a SCOUT First: Pre-hoc Reasoning for Adaptive Detector Allocation in Prompt-Injection Defense & 2026 & OA, S2 & qualifies or contradicts the claim (0.40) \\
Sentinel: SOTA model to protect against prompt injections & 2025 & OA, S2 & qualifies or contradicts the claim (0.40) \\
SentinelX AI: A Post-Quantum Secured Multi-Agent Defense Framework for LLM-Enabled Systems Against Adversarial Prompt and Agent Attacks & 2026 & OA & qualifies or contradicts the claim (0.45) \\
Separating intent from execution: A defense-in-depth security architecture for LLM-based multi-agent systems & 2027 & WoS, OA & qualifies or contradicts the claim (0.75) \\
Shieldllm: A Hybrid Adversarial Prompt Injection Detection Framework for Securing Large Language Models & 2026 & OA, S2 & qualifies or contradicts the claim (0.60) \\
SIEVE: Selective Integrity Verification and Escalation for Defending LLM Agents against Indirect Prompt Injection & 2025 & OA, S2 & qualifies or contradicts the claim (0.70) \\
Signed-prompt: A new approach to prevent prompt injection attacks against LLM-Integrated applications & 2024 & OA, Scopus & qualifies or contradicts the claim (0.60) \\
Silent Failures in Agentic Security Evaluation: A Validated Harness for Tool-Call Mediation Under Indirect Prompt Injection & 2026 & OA, S2 & qualifies or contradicts the claim (0.90) \\
SingGuard-NSFA: Extensible Guardrails for Agentic AI via Generative Reasoning and Real-Time Classification & 2026 & OA, S2 & qualifies or contradicts the claim (0.70) \\
Smarter Saboteurs, Better Fixers: Scaling \& Security in Linear Multi-Agent Workflows & 2026 & OA, S2 & qualifies or contradicts the claim (0.35) \\
SnapGuard: Lightweight Prompt Injection Detection for Screenshot-Based Web Agents & 2026 & OA, S2 & qualifies or contradicts the claim (0.40) \\
SPML: A DSL for Defending Language Models Against Prompt Attacks & 2024 & OA, S2 & qualifies or contradicts the claim (0.45) \\
StateFence: Task-State Protection Against Evidence-Chain Contamination in Blue-Team LLM Agents & 2026 & OA, S2 & qualifies or contradicts the claim (0.70) \\
Staying on Task: Defending Against Prompt Injection via Task Consistency Checking & 2026 & OA, S2 & qualifies or contradicts the claim (0.80) \\
Steerable by Design: Conditional Pretraining for Robust Control of Language Models & 2026 & OA, S2 & qualifies or contradicts the claim (0.65) \\
Stored Is Not Supported: Typed Provenance and Assertion Guardrails for Persistent AI Agents & 2026 & OA, S2 & qualifies or contradicts the claim (0.40) \\
Strengthening Polymorphic Prompt Assembling: Dynamic Separator Generation Against Emerging Prompt Injection Attacks & 2026 & OA, S2 & qualifies or contradicts the claim (0.60) \\
Stronger Enforcement of Instruction Hierarchy via Augmented Intermediate Representations & 2025 & OA, S2 & qualifies or contradicts the claim (0.80) \\
Thai-language chatbot security: Detecting instruction attacks with XLM-RoBERTa and Bi-GRU & 2024 & WoS, OA, Scopus & qualifies or contradicts the claim (0.40) \\
The Agent Knew: Internal States Forecast Unauthorized Actions and Stop Them, and the Lies That Follow & 2026 & OA, S2 & qualifies or contradicts the claim (0.60) \\
The Granularity Mismatch in Agent Security: Argument-Level Provenance Solves Enforcement and Isolates the LLM Reasoning Bottleneck & 2026 & OA, S2 & qualifies or contradicts the claim (0.80) \\
The HTTPS of AI Agents: Securing Every Action Before Execution & 2026 & OA & qualifies or contradicts the claim (0.40) (no-evidence) \\
The Instruction Hierarchy: Training LLMs to Prioritize Privileged Instructions & 2024 & OA, S2 & qualifies or contradicts the claim (0.90) \\
The LLMbda Calculus: AI Agents, Conversations, and Information Flow & 2026 & OA, S2 & qualifies or contradicts the claim (0.85) \\
The Mirror Design Pattern: Strict Data Geometry over Model Scale for Prompt Injection Detection & 2026 & OA, S2 & qualifies or contradicts the claim (0.50) \\
The Model Knows Better: Defending Medical Multi-Agent Systems Against Prompt Injection & 2026 & OA, S2 & qualifies or contradicts the claim (0.65) \\
The Task Shield: Enforcing Task Alignment to Defend Against Indirect Prompt Injection in LLM Agents & 2025 & WoS, S2, OA, Scopus & qualifies or contradicts the claim (0.80) \\
To Protect the LLM Agent Against the Prompt Injection Attack with Polymorphic Prompt & 2025 & WoS, IEEE, OA, Scopus, S2 & qualifies or contradicts the claim (0.70) \\
Tool Contracts and Output Defenses Under Agent Faults: A Pre-Registered Replication and \ensuremath{\tau}²-Bench Extension & 2026 & OA & qualifies or contradicts the claim (0.55) \\
ToolSafe: Enhancing Tool Invocation Safety of LLM-based Agents via Proactive Step-level Guardrail and Feedback & 2026 & Scopus, OA, S2 & qualifies or contradicts the claim (0.65) \\
TotalShield: A Multi-Layered Defense Architecture For Robust Protection Against Prompt-Leaking Attacks On Large Language Models & 2026 & OA & qualifies or contradicts the claim (0.60) \\
Toward Trustworthy Agentic AI: A Multimodal Framework for Preventing Prompt Injection Attacks & 2025 & OA, S2 & qualifies or contradicts the claim (0.40) \\
Towards Graph-Based Detection of Jailbreak and Prompt-Leakage Attacks in LLMs & 2026 & WoS, OA, Scopus, S2 & qualifies or contradicts the claim (0.40) \\
Towards Safe Autonomy for Computer Use & 2026 & OA, S2 & qualifies or contradicts the claim (0.80) \\
TraceCaps: Inline Provenance and Risk Enforcement for Agentic Software Engineering & 2026 & Scopus, S2 & qualifies or contradicts the claim (0.40) \\
TraceGrant: A Contract-Governed Security Framework for the Task-Effect Lifecycle of Networked LLM Agents & 2026 & OA, S2 & qualifies or contradicts the claim (0.75) \\
TRACER-AI: A Multi-Layer Explainable Framework for Prompt Injection, Agent Goal Hijacking, and Tool Misuse Detection in Agentic AI Systems & 2026 & OA, S2 & qualifies or contradicts the claim (0.50) \\
TRUSTDESC: Preventing Tool Poisoning in LLM Applications via Trusted Description Generation & 2026 & OA, S2 & qualifies or contradicts the claim (0.70) \\
Trusted Multi-Agent Generative AI for Real-Time Defense of U.S. Digital Payments: A Privacy-Preserving Governance Framework for Critical and Emerging Financial Technologies & 2026 & OA & qualifies or contradicts the claim (0.40) \\
Trustworthy AI Agent Pipelines via Authenticated Data Structures & 2026 & OA, Scopus, S2 & qualifies or contradicts the claim (0.40) \\
Trustworthy Generative AI for Enterprise Automation: A Secure Agentic Framework for Large-Scale Cloud Platforms & 2026 & IEEE, OA, Scopus, S2 & qualifies or contradicts the claim (0.70) \\
Trustworthy RAG: An Evaluation Agent for Detecting Misinformation and Knowledge Poisoning in Generative AI Systems & 2026 & OA, S2 & qualifies or contradicts the claim (0.50) \\
Twin Agent: Context Residual Compression for Privilege Separated Agents & 2026 & OA, S2 & qualifies or contradicts the claim (0.60) \\
Twin Distillation Against Indirect Prompt Injection: Training Language Models to Distinguish User Instructions from Injected Instructions & 2026 & OA & qualifies or contradicts the claim (0.60) \\
UniGuardian: A Unified Defense for Detecting Prompt Injection, Backdoor Attacks and Adversarial Attacks in Large Language Models & 2026 & OA, S2 & qualifies or contradicts the claim (0.50) \\
Untrusted Content Masking for Web Agents with Security Guarantees & 2026 & OA, S2 & qualifies or contradicts the claim (0.70) \\
Using context-oriented rules and ML analytics for automatic training of LLM protection systems against prompt injection attacks & 2026 & OA & qualifies or contradicts the claim (0.40) \\
VIGIL: Defending LLM Agents Against Tool Stream Injection via Verify-Before-Commit & 2026 & WoS, OA, Scopus, S2 & qualifies or contradicts the claim (0.55) \\
VIGIL: Defending LLM Agents Against Tool-Stream Injection via Verify-Before-Commit & 2026 & OA & qualifies or contradicts the claim (0.50) \\
WARD: Adversarially Robust Defense of Web Agents Against Prompt Injections & 2026 & OA, S2 & qualifies or contradicts the claim (0.70) \\
WebSentinel: Detecting and Localizing Prompt Injection Attacks for Web Agents & 2026 & OA, S2 & qualifies or contradicts the claim (0.50) \\
When Actions Go Off-Task: Detecting and Correcting Misaligned Actions in Computer-Use Agents & 2026 & OA, Scopus, S2 & qualifies or contradicts the claim (0.70) \\
Whiteboxing Large Language Models: Spectral Concentration without Computational Compression in Transformer Inference & 2026 & OA & qualifies or contradicts the claim (0.35) \\
You Can't Steal Nothing: Mitigating Prompt Leakages in LLMs via System Vectors & 2025 & WoS, OA, Scopus, S2 & qualifies or contradicts the claim (0.70) \\
Your Agentic LLMs Secretly Encode Indirect Prompt-Injection Exposure in Hidden States & 2026 & OA, S2 & qualifies or contradicts the claim (0.70) \\
ZERO TRUST FOR MCP-CONNECTED ENTERPRISE AI AGENTS & 2026 & OA, S2 & qualifies or contradicts the claim (0.55) \\
Zero-Shot Embedding Drift Detection: A Lightweight Defense Against Prompt Injections in LLMs & 2026 & OA, S2 & qualifies or contradicts the claim (0.60) \\
Zero-trust Conceptual Architecture for Secure Integration of Large Language Models and Autonomous AI Agents & 2026 & IEEE, OA, Scopus, S2 & qualifies or contradicts the claim (0.40) \\
Zero-Trust Proof-Gated Agent Architecture (ZPAA): A Comprehensive Security Model for the Model Context Protocol (MCP) & 2025 & OA, S2 & qualifies or contradicts the claim (0.55) \\
\end{longtable}
\endgroup

